\documentclass[10pt]{report}
\usepackage[slovak]{babel}
\usepackage{lmodern}
\usepackage{graphicx}
\usepackage{varwidth}
\usepackage{enumitem}
\usepackage{amsmath}
\usepackage{amssymb}
\usepackage{mathtools}
\usepackage{pifont}
\usepackage{arydshln}
\usepackage{tikz}
\usepackage{lscape}
\usepackage{bm}
\usepackage{nicefrac}
\usepackage{multicol}
\usepackage{vwcol}
\usepackage{physics}
\usepackage{fontsize}
\usepackage[landscape]{geometry}
\geometry{a4paper, total={296mm,210mm}, left=-5mm, top=0mm}
\renewcommand{\labelenumi}{\bfseries(\alph{enumi})\phantom{x}}
\newcommand\omicron{o}
\hfuzz=50pt
\setlist[enumerate]{leftmargin=0pt,itemindent=34pt}
\pagenumbering{gobble}
\setlength{\tabcolsep}{0pt}
\begin{document}
\thispagestyle{empty}
\begin{tabular}{c:c}
\begin{minipage}[c][104.5mm][t]{0.5\linewidth}
\begin{center}
\vspace{7mm}
{\huge Monotónnosť a extrémy, skupina \textit{Alpha $\alpha$} -\romannumeral1}\\[5mm]
\begin{minipage}{0.95\linewidth}
\begin{center}
{\small \begin{multicols}{2} \raggedright
\textbf{(a)} Určte monotónnosť funkcie \(f(x)\) v okolí bodu \(x_0\). \\[1mm]
\textbf{(b)} Určte intervaly, na ktorých je funkcia \(g(x)\) klesajúca. \\[1mm]
\textbf{(c)} Určte počet stacionárnych bodov funkcie \(h(t)\). \par \columnbreak
\textbf{(d)} Nadobúda funkcia \(v(x)\) lokálny extrém? \\[1mm]
\textbf{(e)} Určte rozdiel hodnôt lokálneho maxima a lokálneho minima funkcie \(\varphi(x)\). \\[1mm]
\textbf{(f)} Určte súčin \(x\)-ových súradníc bodov globálneho maxima a globálneho minima funkcie \(\omega(x)\) na zadanom intervale. \\[1mm]
\end{multicols}
}
\end{center}
\end{minipage}
\\[4mm]
\begin{minipage}{0.79\linewidth}
\begin{center}
\begin{varwidth}{\linewidth}
\begin{enumerate}
\normalsize
\item $f(x)=\frac{x}{8}+\frac{8}{x} \enspace , \enspace x_0=7$\quad \dotfill\; ???\;\dotfill \quad \text{klesajúca}
\item $g(x)=x^2\, e^{-x+3}$\quad \dotfill\; ???\;\dotfill \quad $\left\langle2 , 0\right\rangle$
\item $h(t)=\cfrac{\ln t}{t^3}$\quad \dotfill\; ???\;\dotfill \quad 1
\item $v(x)=2-\sqrt[3]{(2x+2)^2}$\quad \dotfill\; ???\;\dotfill \quad $\text{nie}$
\item $\varphi(x)=\cfrac{x^2+4x+28}{x+1}$\quad \dotfill\; ???\;\dotfill \quad $20$
\item $\omega(x)=-\frac{2}{3}x^3+\frac{13}{2}x^2-15x+5 \enspace , \enspace x\in\left\langle-1 , 6\right\rangle$\quad \dotfill\; ???\;\dotfill \quad $-\frac{3}{2}$
\end{enumerate}
\end{varwidth}
\end{center}
\end{minipage}
\begin{minipage}{0.20\linewidth}
\begin{center}
{\Huge\bfseries 1.} \\[2mm]
\includegraphics[height=37mm]{images/braille.png}
{\small Braillovo písmeno}
\end{center}
\end{minipage}
\end{center}
\end{minipage}
&
\begin{minipage}[c][104.5mm][t]{0.5\linewidth}
\begin{center}
\vspace{7mm}
{\huge Monotónnosť a extrémy, skupina \textit{Alpha $\alpha$} -\romannumeral2}\\[5mm]
\begin{minipage}{0.95\linewidth}
\begin{center}
{\small \begin{multicols}{2} \raggedright
\textbf{(a)} Určte monotónnosť funkcie \(f(x)\) v okolí bodu \(x_0\). \\[1mm]
\textbf{(b)} Určte intervaly, na ktorých je funkcia \(g(x)\) klesajúca. \\[1mm]
\textbf{(c)} Určte počet stacionárnych bodov funkcie \(h(t)\). \par \columnbreak
\textbf{(d)} Nadobúda funkcia \(v(x)\) lokálny extrém? \\[1mm]
\textbf{(e)} Určte rozdiel hodnôt lokálneho maxima a lokálneho minima funkcie \(\varphi(x)\). \\[1mm]
\textbf{(f)} Určte súčin \(x\)-ových súradníc bodov globálneho maxima a globálneho minima funkcie \(\omega(x)\) na zadanom intervale. \\[1mm]
\end{multicols}
}
\end{center}
\end{minipage}
\\[4mm]
\begin{minipage}{0.79\linewidth}
\begin{center}
\begin{varwidth}{\linewidth}
\begin{enumerate}
\normalsize
\item $f(x)=\frac{x}{3}+\frac{3}{x} \enspace , \enspace x_0=1$\quad \dotfill\; ???\;\dotfill \quad \text{klesajúca}
\item $g(x)=x^2\, e^{7x-6}$\quad \dotfill\; ???\;\dotfill \quad $\left\langle-\frac{2}{7} , 0\right\rangle$
\item $h(t)=\cfrac{\ln t}{7t^4}$\quad \dotfill\; ???\;\dotfill \quad 1
\item $v(x)=2-\sqrt[3]{(x+6)^2}$\quad \dotfill\; ???\;\dotfill \quad $\text{nie}$
\item $\varphi(x)=\cfrac{x^2-2x-2}{x-3}$\quad \dotfill\; ???\;\dotfill \quad $-2$
\item $\omega(x)=\frac{5}{3}x^3-\frac{1}{2}x^2-4x-2 \enspace , \enspace x\in\left\langle-1 , 3\right\rangle$\quad \dotfill\; ???\;\dotfill \quad $2$
\end{enumerate}
\end{varwidth}
\end{center}
\end{minipage}
\begin{minipage}{0.20\linewidth}
\begin{center}
{\Huge\bfseries 2.} \\[2mm]
\includegraphics[height=37mm]{images/braille.png}
{\small Braillovo písmeno}
\end{center}
\end{minipage}
\end{center}
\end{minipage}
\\ \hdashline
\begin{minipage}[c][104.5mm][t]{0.5\linewidth}
\begin{center}
\vspace{7mm}
{\huge Monotónnosť a extrémy, skupina \textit{Alpha $\alpha$} -\romannumeral3}\\[5mm]
\begin{minipage}{0.95\linewidth}
\begin{center}
{\small \begin{multicols}{2} \raggedright
\textbf{(a)} Určte monotónnosť funkcie \(f(x)\) v okolí bodu \(x_0\). \\[1mm]
\textbf{(b)} Určte intervaly, na ktorých je funkcia \(g(x)\) klesajúca. \\[1mm]
\textbf{(c)} Určte počet stacionárnych bodov funkcie \(h(t)\). \par \columnbreak
\textbf{(d)} Nadobúda funkcia \(v(x)\) lokálny extrém? \\[1mm]
\textbf{(e)} Určte rozdiel hodnôt lokálneho maxima a lokálneho minima funkcie \(\varphi(x)\). \\[1mm]
\textbf{(f)} Určte súčin \(x\)-ových súradníc bodov globálneho maxima a globálneho minima funkcie \(\omega(x)\) na zadanom intervale. \\[1mm]
\end{multicols}
}
\end{center}
\end{minipage}
\\[4mm]
\begin{minipage}{0.79\linewidth}
\begin{center}
\begin{varwidth}{\linewidth}
\begin{enumerate}
\normalsize
\item $f(x)=\frac{x}{7}+\frac{7}{x} \enspace , \enspace x_0=2$\quad \dotfill\; ???\;\dotfill \quad \text{klesajúca}
\item $g(x)=x^2\, e^{x-1}$\quad \dotfill\; ???\;\dotfill \quad $\left(-\infty , 0\right\rangle\enspace , \enspace \left\langle-2 , \infty\right)$
\item $h(t)=\cfrac{\ln t}{8t^4}$\quad \dotfill\; ???\;\dotfill \quad 1
\item $v(x)=-4-\sqrt[3]{(3x+1)^2}$\quad \dotfill\; ???\;\dotfill \quad $\text{nie}$
\item $\varphi(x)=\cfrac{x^2-2x+33}{x-3}$\quad \dotfill\; ???\;\dotfill \quad $-24$
\item $\omega(x)=2x^3-\frac{7}{2}x^2+x-5 \enspace , \enspace x\in\left\langle0 , 3\right\rangle$\quad \dotfill\; ???\;\dotfill \quad $1$
\end{enumerate}
\end{varwidth}
\end{center}
\end{minipage}
\begin{minipage}{0.20\linewidth}
\begin{center}
{\Huge\bfseries 3.} \\[2mm]
\includegraphics[height=37mm]{images/braille.png}
{\small Braillovo písmeno}
\end{center}
\end{minipage}
\end{center}
\end{minipage}
&
\begin{minipage}[c][104.5mm][t]{0.5\linewidth}
\begin{center}
\vspace{7mm}
{\huge Monotónnosť a extrémy, skupina \textit{Alpha $\alpha$} -\romannumeral4}\\[5mm]
\begin{minipage}{0.95\linewidth}
\begin{center}
{\small \begin{multicols}{2} \raggedright
\textbf{(a)} Určte monotónnosť funkcie \(f(x)\) v okolí bodu \(x_0\). \\[1mm]
\textbf{(b)} Určte intervaly, na ktorých je funkcia \(g(x)\) klesajúca. \\[1mm]
\textbf{(c)} Určte počet stacionárnych bodov funkcie \(h(t)\). \par \columnbreak
\textbf{(d)} Nadobúda funkcia \(v(x)\) lokálny extrém? \\[1mm]
\textbf{(e)} Určte rozdiel hodnôt lokálneho maxima a lokálneho minima funkcie \(\varphi(x)\). \\[1mm]
\textbf{(f)} Určte súčin \(x\)-ových súradníc bodov globálneho maxima a globálneho minima funkcie \(\omega(x)\) na zadanom intervale. \\[1mm]
\end{multicols}
}
\end{center}
\end{minipage}
\\[4mm]
\begin{minipage}{0.79\linewidth}
\begin{center}
\begin{varwidth}{\linewidth}
\begin{enumerate}
\normalsize
\item $f(x)=\frac{x}{6}+\frac{6}{x} \enspace , \enspace x_0=1$\quad \dotfill\; ???\;\dotfill \quad \text{klesajúca}
\item $g(x)=x^2\, e^{6x-2}$\quad \dotfill\; ???\;\dotfill \quad $\left(-\infty , 0\right\rangle\enspace , \enspace \left\langle-\frac{1}{3} , \infty\right)$
\item $h(t)=\cfrac{\ln t}{3t^3}$\quad \dotfill\; ???\;\dotfill \quad 1
\item $v(x)=8-\sqrt[3]{(-9x+7)^2}$\quad \dotfill\; ???\;\dotfill \quad $\text{nie}$
\item $\varphi(x)=\cfrac{x^2+6x-2}{x-3}$\quad \dotfill\; ???\;\dotfill \quad $-18$
\item $\omega(x)=\frac{8}{3}x^3-\frac{9}{2}x^2+x-5 \enspace , \enspace x\in\left\langle0 , 3\right\rangle$\quad \dotfill\; ???\;\dotfill \quad $2$
\end{enumerate}
\end{varwidth}
\end{center}
\end{minipage}
\begin{minipage}{0.20\linewidth}
\begin{center}
{\Huge\bfseries 4.} \\[2mm]
\includegraphics[height=37mm]{images/braille.png}
{\small Braillovo písmeno}
\end{center}
\end{minipage}
\end{center}
\end{minipage}
%
\end{tabular}
\newpage
\thispagestyle{empty}
\begin{tabular}{c:c}
\begin{minipage}[c][104.5mm][t]{0.5\linewidth}
\begin{center}
\vspace{7mm}
{\huge Monotónnosť a extrémy, skupina \textit{Beta $\beta$} -\romannumeral1}\\[5mm]
\begin{minipage}{0.95\linewidth}
\begin{center}
{\small \begin{multicols}{2} \raggedright
\textbf{(a)} Určte monotónnosť funkcie \(f(x)\) v okolí bodu \(x_0\). \\[1mm]
\textbf{(b)} Určte intervaly, na ktorých je funkcia \(g(x)\) klesajúca. \\[1mm]
\textbf{(c)} Určte počet stacionárnych bodov funkcie \(h(t)\). \par \columnbreak
\textbf{(d)} Nadobúda funkcia \(v(x)\) lokálny extrém? \\[1mm]
\textbf{(e)} Určte rozdiel hodnôt lokálneho maxima a lokálneho minima funkcie \(\varphi(x)\). \\[1mm]
\textbf{(f)} Určte súčin \(x\)-ových súradníc bodov globálneho maxima a globálneho minima funkcie \(\omega(x)\) na zadanom intervale. \\[1mm]
\end{multicols}
}
\end{center}
\end{minipage}
\\[4mm]
\begin{minipage}{0.79\linewidth}
\begin{center}
\begin{varwidth}{\linewidth}
\begin{enumerate}
\normalsize
\item $f(x)=\frac{x}{5}+\frac{5}{x} \enspace , \enspace x_0=-6$\quad \dotfill\; ???\;\dotfill \quad \text{rastúca}
\item $g(x)=x^2\, e^{-7x-6}$\quad \dotfill\; ???\;\dotfill \quad $\left(-\infty , 0\right\rangle\enspace , \enspace \left\langle\frac{2}{7} , \infty\right)$
\item $h(t)=\cfrac{\ln t}{8t^3}$\quad \dotfill\; ???\;\dotfill \quad 1
\item $v(x)=1-\sqrt[5]{(-6x-3)^2}$\quad \dotfill\; ???\;\dotfill \quad $\text{áno}$
\item $\varphi(x)=\cfrac{x^2-3x-3}{x-4}$\quad \dotfill\; ???\;\dotfill \quad $-6$
\item $\omega(x)=\frac{4}{3}x^3-\frac{1}{2}x^2-3x+5 \enspace , \enspace x\in\left\langle-1 , 3\right\rangle$\quad \dotfill\; ???\;\dotfill \quad $2$
\end{enumerate}
\end{varwidth}
\end{center}
\end{minipage}
\begin{minipage}{0.20\linewidth}
\begin{center}
{\Huge\bfseries 1.} \\[2mm]
\includegraphics[height=37mm]{images/braille.png}
{\small Braillovo písmeno}
\end{center}
\end{minipage}
\end{center}
\end{minipage}
&
\begin{minipage}[c][104.5mm][t]{0.5\linewidth}
\begin{center}
\vspace{7mm}
{\huge Monotónnosť a extrémy, skupina \textit{Beta $\beta$} -\romannumeral2}\\[5mm]
\begin{minipage}{0.95\linewidth}
\begin{center}
{\small \begin{multicols}{2} \raggedright
\textbf{(a)} Určte monotónnosť funkcie \(f(x)\) v okolí bodu \(x_0\). \\[1mm]
\textbf{(b)} Určte intervaly, na ktorých je funkcia \(g(x)\) klesajúca. \\[1mm]
\textbf{(c)} Určte počet stacionárnych bodov funkcie \(h(t)\). \par \columnbreak
\textbf{(d)} Nadobúda funkcia \(v(x)\) lokálny extrém? \\[1mm]
\textbf{(e)} Určte rozdiel hodnôt lokálneho maxima a lokálneho minima funkcie \(\varphi(x)\). \\[1mm]
\textbf{(f)} Určte súčin \(x\)-ových súradníc bodov globálneho maxima a globálneho minima funkcie \(\omega(x)\) na zadanom intervale. \\[1mm]
\end{multicols}
}
\end{center}
\end{minipage}
\\[4mm]
\begin{minipage}{0.79\linewidth}
\begin{center}
\begin{varwidth}{\linewidth}
\begin{enumerate}
\normalsize
\item $f(x)=\frac{x}{5}+\frac{5}{x} \enspace , \enspace x_0=-7$\quad \dotfill\; ???\;\dotfill \quad \text{rastúca}
\item $g(x)=x^2\, e^{x+1}$\quad \dotfill\; ???\;\dotfill \quad $\left\langle-2 , 0\right\rangle$
\item $h(t)=\cfrac{\ln t}{2t^4}$\quad \dotfill\; ???\;\dotfill \quad 1
\item $v(x)=-2-\sqrt[7]{(-2x+2)^2}$\quad \dotfill\; ???\;\dotfill \quad $\text{nie}$
\item $\varphi(x)=\cfrac{x^2-5x+7}{x-3}$\quad \dotfill\; ???\;\dotfill \quad $4$
\item $\omega(x)=-2x^3+\frac{13}{2}x^2+15x-1 \enspace , \enspace x\in\left\langle-3 , 4\right\rangle$\quad \dotfill\; ???\;\dotfill \quad $\frac{5}{18}$
\end{enumerate}
\end{varwidth}
\end{center}
\end{minipage}
\begin{minipage}{0.20\linewidth}
\begin{center}
{\Huge\bfseries 2.} \\[2mm]
\includegraphics[height=37mm]{images/braille.png}
{\small Braillovo písmeno}
\end{center}
\end{minipage}
\end{center}
\end{minipage}
\\ \hdashline
\begin{minipage}[c][104.5mm][t]{0.5\linewidth}
\begin{center}
\vspace{7mm}
{\huge Monotónnosť a extrémy, skupina \textit{Beta $\beta$} -\romannumeral3}\\[5mm]
\begin{minipage}{0.95\linewidth}
\begin{center}
{\small \begin{multicols}{2} \raggedright
\textbf{(a)} Určte monotónnosť funkcie \(f(x)\) v okolí bodu \(x_0\). \\[1mm]
\textbf{(b)} Určte intervaly, na ktorých je funkcia \(g(x)\) klesajúca. \\[1mm]
\textbf{(c)} Určte počet stacionárnych bodov funkcie \(h(t)\). \par \columnbreak
\textbf{(d)} Nadobúda funkcia \(v(x)\) lokálny extrém? \\[1mm]
\textbf{(e)} Určte rozdiel hodnôt lokálneho maxima a lokálneho minima funkcie \(\varphi(x)\). \\[1mm]
\textbf{(f)} Určte súčin \(x\)-ových súradníc bodov globálneho maxima a globálneho minima funkcie \(\omega(x)\) na zadanom intervale. \\[1mm]
\end{multicols}
}
\end{center}
\end{minipage}
\\[4mm]
\begin{minipage}{0.79\linewidth}
\begin{center}
\begin{varwidth}{\linewidth}
\begin{enumerate}
\normalsize
\item $f(x)=\frac{x}{7}+\frac{7}{x} \enspace , \enspace x_0=5$\quad \dotfill\; ???\;\dotfill \quad \text{klesajúca}
\item $g(x)=x^2\, e^{3x+3}$\quad \dotfill\; ???\;\dotfill \quad $\left(-\infty , 0\right\rangle\enspace , \enspace \left\langle-\frac{2}{3} , \infty\right)$
\item $h(t)=\cfrac{\ln t}{8t^3}$\quad \dotfill\; ???\;\dotfill \quad 2
\item $v(x)=-3-\sqrt[7]{(-5x+6)^2}$\quad \dotfill\; ???\;\dotfill \quad $\text{nie}$
\item $\varphi(x)=\cfrac{x^2+6x+2}{x-1}$\quad \dotfill\; ???\;\dotfill \quad $-10$
\item $\omega(x)=-\frac{1}{3}x^3-\frac{3}{2}x^2+4x-7 \enspace , \enspace x\in\left\langle-7 , 2\right\rangle$\quad \dotfill\; ???\;\dotfill \quad $4$
\end{enumerate}
\end{varwidth}
\end{center}
\end{minipage}
\begin{minipage}{0.20\linewidth}
\begin{center}
{\Huge\bfseries 3.} \\[2mm]
\includegraphics[height=37mm]{images/braille.png}
{\small Braillovo písmeno}
\end{center}
\end{minipage}
\end{center}
\end{minipage}
&
\begin{minipage}[c][104.5mm][t]{0.5\linewidth}
\begin{center}
\vspace{7mm}
{\huge Monotónnosť a extrémy, skupina \textit{Beta $\beta$} -\romannumeral4}\\[5mm]
\begin{minipage}{0.95\linewidth}
\begin{center}
{\small \begin{multicols}{2} \raggedright
\textbf{(a)} Určte monotónnosť funkcie \(f(x)\) v okolí bodu \(x_0\). \\[1mm]
\textbf{(b)} Určte intervaly, na ktorých je funkcia \(g(x)\) klesajúca. \\[1mm]
\textbf{(c)} Určte počet stacionárnych bodov funkcie \(h(t)\). \par \columnbreak
\textbf{(d)} Nadobúda funkcia \(v(x)\) lokálny extrém? \\[1mm]
\textbf{(e)} Určte rozdiel hodnôt lokálneho maxima a lokálneho minima funkcie \(\varphi(x)\). \\[1mm]
\textbf{(f)} Určte súčin \(x\)-ových súradníc bodov globálneho maxima a globálneho minima funkcie \(\omega(x)\) na zadanom intervale. \\[1mm]
\end{multicols}
}
\end{center}
\end{minipage}
\\[4mm]
\begin{minipage}{0.79\linewidth}
\begin{center}
\begin{varwidth}{\linewidth}
\begin{enumerate}
\normalsize
\item $f(x)=\frac{x}{5}+\frac{5}{x} \enspace , \enspace x_0=-10$\quad \dotfill\; ???\;\dotfill \quad \text{rastúca}
\item $g(x)=x^2\, e^{6x-1}$\quad \dotfill\; ???\;\dotfill \quad $\left(-\infty , 0\right\rangle\enspace , \enspace \left\langle-\frac{1}{3} , \infty\right)$
\item $h(t)=\cfrac{\ln t}{3t^4}$\quad \dotfill\; ???\;\dotfill \quad 1
\item $v(x)=-8-\sqrt[5]{(-7x+4)^2}$\quad \dotfill\; ???\;\dotfill \quad $\text{nie}$
\item $\varphi(x)=\cfrac{x^2+3x-24}{x+7}$\quad \dotfill\; ???\;\dotfill \quad $8$
\item $\omega(x)=-x^3+\frac{13}{2}x^2-4x-5 \enspace , \enspace x\in\left\langle-2 , 5\right\rangle$\quad \dotfill\; ???\;\dotfill \quad $\frac{4}{9}$
\end{enumerate}
\end{varwidth}
\end{center}
\end{minipage}
\begin{minipage}{0.20\linewidth}
\begin{center}
{\Huge\bfseries 4.} \\[2mm]
\includegraphics[height=37mm]{images/braille.png}
{\small Braillovo písmeno}
\end{center}
\end{minipage}
\end{center}
\end{minipage}
%
\end{tabular}
\newpage
\thispagestyle{empty}
\begin{tabular}{c:c}
\begin{minipage}[c][104.5mm][t]{0.5\linewidth}
\begin{center}
\vspace{7mm}
{\huge Monotónnosť a extrémy, skupina \textit{Gamma $\gamma$} -\romannumeral1}\\[5mm]
\begin{minipage}{0.95\linewidth}
\begin{center}
{\small \begin{multicols}{2} \raggedright
\textbf{(a)} Určte monotónnosť funkcie \(f(x)\) v okolí bodu \(x_0\). \\[1mm]
\textbf{(b)} Určte intervaly, na ktorých je funkcia \(g(x)\) klesajúca. \\[1mm]
\textbf{(c)} Určte počet stacionárnych bodov funkcie \(h(t)\). \par \columnbreak
\textbf{(d)} Nadobúda funkcia \(v(x)\) lokálny extrém? \\[1mm]
\textbf{(e)} Určte rozdiel hodnôt lokálneho maxima a lokálneho minima funkcie \(\varphi(x)\). \\[1mm]
\textbf{(f)} Určte súčin \(x\)-ových súradníc bodov globálneho maxima a globálneho minima funkcie \(\omega(x)\) na zadanom intervale. \\[1mm]
\end{multicols}
}
\end{center}
\end{minipage}
\\[4mm]
\begin{minipage}{0.79\linewidth}
\begin{center}
\begin{varwidth}{\linewidth}
\begin{enumerate}
\normalsize
\item $f(x)=\frac{x}{4}+\frac{4}{x} \enspace , \enspace x_0=-10$\quad \dotfill\; ???\;\dotfill \quad \text{klesajúca}
\item $g(x)=x^2\, e^{-x-7}$\quad \dotfill\; ???\;\dotfill \quad $\left(-\infty , 0\right\rangle\enspace , \enspace \left\langle2 , \infty\right)$
\item $h(t)=\cfrac{\ln t}{2t^5}$\quad \dotfill\; ???\;\dotfill \quad 2
\item $v(x)=-5-\sqrt[7]{(-6x+6)^2}$\quad \dotfill\; ???\;\dotfill \quad $\text{nie}$
\item $\varphi(x)=\cfrac{x^2-5x+5}{x-4}$\quad \dotfill\; ???\;\dotfill \quad $-4$
\item $\omega(x)=-\frac{1}{3}x^3+\frac{1}{2}x^2+2x-1 \enspace , \enspace x\in\left\langle-4 , 3\right\rangle$\quad \dotfill\; ???\;\dotfill \quad $4$
\end{enumerate}
\end{varwidth}
\end{center}
\end{minipage}
\begin{minipage}{0.20\linewidth}
\begin{center}
{\Huge\bfseries 1.} \\[2mm]
\includegraphics[height=37mm]{images/braille.png}
{\small Braillovo písmeno}
\end{center}
\end{minipage}
\end{center}
\end{minipage}
&
\begin{minipage}[c][104.5mm][t]{0.5\linewidth}
\begin{center}
\vspace{7mm}
{\huge Monotónnosť a extrémy, skupina \textit{Gamma $\gamma$} -\romannumeral2}\\[5mm]
\begin{minipage}{0.95\linewidth}
\begin{center}
{\small \begin{multicols}{2} \raggedright
\textbf{(a)} Určte monotónnosť funkcie \(f(x)\) v okolí bodu \(x_0\). \\[1mm]
\textbf{(b)} Určte intervaly, na ktorých je funkcia \(g(x)\) klesajúca. \\[1mm]
\textbf{(c)} Určte počet stacionárnych bodov funkcie \(h(t)\). \par \columnbreak
\textbf{(d)} Nadobúda funkcia \(v(x)\) lokálny extrém? \\[1mm]
\textbf{(e)} Určte rozdiel hodnôt lokálneho maxima a lokálneho minima funkcie \(\varphi(x)\). \\[1mm]
\textbf{(f)} Určte súčin \(x\)-ových súradníc bodov globálneho maxima a globálneho minima funkcie \(\omega(x)\) na zadanom intervale. \\[1mm]
\end{multicols}
}
\end{center}
\end{minipage}
\\[4mm]
\begin{minipage}{0.79\linewidth}
\begin{center}
\begin{varwidth}{\linewidth}
\begin{enumerate}
\normalsize
\item $f(x)=\frac{x}{2}+\frac{2}{x} \enspace , \enspace x_0=1$\quad \dotfill\; ???\;\dotfill \quad \text{klesajúca}
\item $g(x)=x^2\, e^{-x-2}$\quad \dotfill\; ???\;\dotfill \quad $\left\langle2 , 0\right\rangle$
\item $h(t)=\cfrac{\ln t}{t^5}$\quad \dotfill\; ???\;\dotfill \quad 2
\item $v(x)=8-\sqrt[3]{(x-3)^2}$\quad \dotfill\; ???\;\dotfill \quad $\text{nie}$
\item $\varphi(x)=\cfrac{x^2+x-11}{x+5}$\quad \dotfill\; ???\;\dotfill \quad $12$
\item $\omega(x)=\frac{5}{3}x^3-\frac{11}{2}x^2-12x+1 \enspace , \enspace x\in\left\langle-1 , 5\right\rangle$\quad \dotfill\; ???\;\dotfill \quad $12$
\end{enumerate}
\end{varwidth}
\end{center}
\end{minipage}
\begin{minipage}{0.20\linewidth}
\begin{center}
{\Huge\bfseries 2.} \\[2mm]
\includegraphics[height=37mm]{images/braille.png}
{\small Braillovo písmeno}
\end{center}
\end{minipage}
\end{center}
\end{minipage}
\\ \hdashline
\begin{minipage}[c][104.5mm][t]{0.5\linewidth}
\begin{center}
\vspace{7mm}
{\huge Monotónnosť a extrémy, skupina \textit{Gamma $\gamma$} -\romannumeral3}\\[5mm]
\begin{minipage}{0.95\linewidth}
\begin{center}
{\small \begin{multicols}{2} \raggedright
\textbf{(a)} Určte monotónnosť funkcie \(f(x)\) v okolí bodu \(x_0\). \\[1mm]
\textbf{(b)} Určte intervaly, na ktorých je funkcia \(g(x)\) klesajúca. \\[1mm]
\textbf{(c)} Určte počet stacionárnych bodov funkcie \(h(t)\). \par \columnbreak
\textbf{(d)} Nadobúda funkcia \(v(x)\) lokálny extrém? \\[1mm]
\textbf{(e)} Určte rozdiel hodnôt lokálneho maxima a lokálneho minima funkcie \(\varphi(x)\). \\[1mm]
\textbf{(f)} Určte súčin \(x\)-ových súradníc bodov globálneho maxima a globálneho minima funkcie \(\omega(x)\) na zadanom intervale. \\[1mm]
\end{multicols}
}
\end{center}
\end{minipage}
\\[4mm]
\begin{minipage}{0.79\linewidth}
\begin{center}
\begin{varwidth}{\linewidth}
\begin{enumerate}
\normalsize
\item $f(x)=\frac{x}{7}+\frac{7}{x} \enspace , \enspace x_0=-8$\quad \dotfill\; ???\;\dotfill \quad \text{rastúca}
\item $g(x)=x^2\, e^{7x-6}$\quad \dotfill\; ???\;\dotfill \quad $\left(-\infty , 0\right\rangle\enspace , \enspace \left\langle-\frac{2}{7} , \infty\right)$
\item $h(t)=\cfrac{\ln t}{7t^3}$\quad \dotfill\; ???\;\dotfill \quad 1
\item $v(x)=1-\sqrt[7]{(-x-1)^2}$\quad \dotfill\; ???\;\dotfill \quad $\text{nie}$
\item $\varphi(x)=\cfrac{x^2-2x+16}{x-2}$\quad \dotfill\; ???\;\dotfill \quad $-18$
\item $\omega(x)=-\frac{1}{3}x^3-4x^2-12x+4 \enspace , \enspace x\in\left\langle-9 , -1\right\rangle$\quad \dotfill\; ???\;\dotfill \quad $54$
\end{enumerate}
\end{varwidth}
\end{center}
\end{minipage}
\begin{minipage}{0.20\linewidth}
\begin{center}
{\Huge\bfseries 3.} \\[2mm]
\includegraphics[height=37mm]{images/braille.png}
{\small Braillovo písmeno}
\end{center}
\end{minipage}
\end{center}
\end{minipage}
&
\begin{minipage}[c][104.5mm][t]{0.5\linewidth}
\begin{center}
\vspace{7mm}
{\huge Monotónnosť a extrémy, skupina \textit{Gamma $\gamma$} -\romannumeral4}\\[5mm]
\begin{minipage}{0.95\linewidth}
\begin{center}
{\small \begin{multicols}{2} \raggedright
\textbf{(a)} Určte monotónnosť funkcie \(f(x)\) v okolí bodu \(x_0\). \\[1mm]
\textbf{(b)} Určte intervaly, na ktorých je funkcia \(g(x)\) klesajúca. \\[1mm]
\textbf{(c)} Určte počet stacionárnych bodov funkcie \(h(t)\). \par \columnbreak
\textbf{(d)} Nadobúda funkcia \(v(x)\) lokálny extrém? \\[1mm]
\textbf{(e)} Určte rozdiel hodnôt lokálneho maxima a lokálneho minima funkcie \(\varphi(x)\). \\[1mm]
\textbf{(f)} Určte súčin \(x\)-ových súradníc bodov globálneho maxima a globálneho minima funkcie \(\omega(x)\) na zadanom intervale. \\[1mm]
\end{multicols}
}
\end{center}
\end{minipage}
\\[4mm]
\begin{minipage}{0.79\linewidth}
\begin{center}
\begin{varwidth}{\linewidth}
\begin{enumerate}
\normalsize
\item $f(x)=\frac{x}{4}+\frac{4}{x} \enspace , \enspace x_0=1$\quad \dotfill\; ???\;\dotfill \quad \text{klesajúca}
\item $g(x)=x^2\, e^{-4x+9}$\quad \dotfill\; ???\;\dotfill \quad $\left\langle\frac{1}{2} , 0\right\rangle$
\item $h(t)=\cfrac{\ln t}{7t^3}$\quad \dotfill\; ???\;\dotfill \quad 2
\item $v(x)=2-\sqrt[5]{(-4x-9)^2}$\quad \dotfill\; ???\;\dotfill \quad $\text{nie}$
\item $\varphi(x)=\cfrac{x^2-5x+7}{x-3}$\quad \dotfill\; ???\;\dotfill \quad $4$
\item $\omega(x)=\frac{1}{3}x^3-\frac{5}{2}x^2+4x-2 \enspace , \enspace x\in\left\langle0 , 6\right\rangle$\quad \dotfill\; ???\;\dotfill \quad $20$
\end{enumerate}
\end{varwidth}
\end{center}
\end{minipage}
\begin{minipage}{0.20\linewidth}
\begin{center}
{\Huge\bfseries 4.} \\[2mm]
\includegraphics[height=37mm]{images/braille.png}
{\small Braillovo písmeno}
\end{center}
\end{minipage}
\end{center}
\end{minipage}
%
\end{tabular}
\newpage
\thispagestyle{empty}
\begin{tabular}{c:c}
\begin{minipage}[c][104.5mm][t]{0.5\linewidth}
\begin{center}
\vspace{7mm}
{\huge Monotónnosť a extrémy, skupina \textit{Delta $\delta$} -\romannumeral1}\\[5mm]
\begin{minipage}{0.95\linewidth}
\begin{center}
{\small \begin{multicols}{2} \raggedright
\textbf{(a)} Určte monotónnosť funkcie \(f(x)\) v okolí bodu \(x_0\). \\[1mm]
\textbf{(b)} Určte intervaly, na ktorých je funkcia \(g(x)\) klesajúca. \\[1mm]
\textbf{(c)} Určte počet stacionárnych bodov funkcie \(h(t)\). \par \columnbreak
\textbf{(d)} Nadobúda funkcia \(v(x)\) lokálny extrém? \\[1mm]
\textbf{(e)} Určte rozdiel hodnôt lokálneho maxima a lokálneho minima funkcie \(\varphi(x)\). \\[1mm]
\textbf{(f)} Určte súčin \(x\)-ových súradníc bodov globálneho maxima a globálneho minima funkcie \(\omega(x)\) na zadanom intervale. \\[1mm]
\end{multicols}
}
\end{center}
\end{minipage}
\\[4mm]
\begin{minipage}{0.79\linewidth}
\begin{center}
\begin{varwidth}{\linewidth}
\begin{enumerate}
\normalsize
\item $f(x)=\frac{x}{5}+\frac{5}{x} \enspace , \enspace x_0=-6$\quad \dotfill\; ???\;\dotfill \quad \text{klesajúca}
\item $g(x)=x^2\, e^{6x-8}$\quad \dotfill\; ???\;\dotfill \quad $\left\langle-\frac{1}{3} , 0\right\rangle$
\item $h(t)=\cfrac{\ln t}{6t^5}$\quad \dotfill\; ???\;\dotfill \quad 2
\item $v(x)=3-\sqrt[5]{(5x+7)^2}$\quad \dotfill\; ???\;\dotfill \quad $\text{nie}$
\item $\varphi(x)=\cfrac{x^2-4x-11}{x+2}$\quad \dotfill\; ???\;\dotfill \quad $4$
\item $\omega(x)=\frac{7}{3}x^3-\frac{1}{2}x^2-6x+2 \enspace , \enspace x\in\left\langle-1 , 3\right\rangle$\quad \dotfill\; ???\;\dotfill \quad $3$
\end{enumerate}
\end{varwidth}
\end{center}
\end{minipage}
\begin{minipage}{0.20\linewidth}
\begin{center}
{\Huge\bfseries 1.} \\[2mm]
\includegraphics[height=37mm]{images/braille.png}
{\small Braillovo písmeno}
\end{center}
\end{minipage}
\end{center}
\end{minipage}
&
\begin{minipage}[c][104.5mm][t]{0.5\linewidth}
\begin{center}
\vspace{7mm}
{\huge Monotónnosť a extrémy, skupina \textit{Delta $\delta$} -\romannumeral2}\\[5mm]
\begin{minipage}{0.95\linewidth}
\begin{center}
{\small \begin{multicols}{2} \raggedright
\textbf{(a)} Určte monotónnosť funkcie \(f(x)\) v okolí bodu \(x_0\). \\[1mm]
\textbf{(b)} Určte intervaly, na ktorých je funkcia \(g(x)\) klesajúca. \\[1mm]
\textbf{(c)} Určte počet stacionárnych bodov funkcie \(h(t)\). \par \columnbreak
\textbf{(d)} Nadobúda funkcia \(v(x)\) lokálny extrém? \\[1mm]
\textbf{(e)} Určte rozdiel hodnôt lokálneho maxima a lokálneho minima funkcie \(\varphi(x)\). \\[1mm]
\textbf{(f)} Určte súčin \(x\)-ových súradníc bodov globálneho maxima a globálneho minima funkcie \(\omega(x)\) na zadanom intervale. \\[1mm]
\end{multicols}
}
\end{center}
\end{minipage}
\\[4mm]
\begin{minipage}{0.79\linewidth}
\begin{center}
\begin{varwidth}{\linewidth}
\begin{enumerate}
\normalsize
\item $f(x)=\frac{x}{2}+\frac{2}{x} \enspace , \enspace x_0=-6$\quad \dotfill\; ???\;\dotfill \quad \text{rastúca}
\item $g(x)=x^2\, e^{-x+5}$\quad \dotfill\; ???\;\dotfill \quad $\left(-\infty , 0\right\rangle\enspace , \enspace \left\langle2 , \infty\right)$
\item $h(t)=\cfrac{\ln t}{2t^3}$\quad \dotfill\; ???\;\dotfill \quad 2
\item $v(x)=-1-\sqrt[5]{(3x-1)^2}$\quad \dotfill\; ???\;\dotfill \quad $\text{nie}$
\item $\varphi(x)=\cfrac{x^2-3x-9}{x+2}$\quad \dotfill\; ???\;\dotfill \quad $-6$
\item $\omega(x)=\frac{1}{3}x^3+\frac{7}{2}x^2+10x-3 \enspace , \enspace x\in\left\langle-6 , 0\right\rangle$\quad \dotfill\; ???\;\dotfill \quad $4$
\end{enumerate}
\end{varwidth}
\end{center}
\end{minipage}
\begin{minipage}{0.20\linewidth}
\begin{center}
{\Huge\bfseries 2.} \\[2mm]
\includegraphics[height=37mm]{images/braille.png}
{\small Braillovo písmeno}
\end{center}
\end{minipage}
\end{center}
\end{minipage}
\\ \hdashline
\begin{minipage}[c][104.5mm][t]{0.5\linewidth}
\begin{center}
\vspace{7mm}
{\huge Monotónnosť a extrémy, skupina \textit{Delta $\delta$} -\romannumeral3}\\[5mm]
\begin{minipage}{0.95\linewidth}
\begin{center}
{\small \begin{multicols}{2} \raggedright
\textbf{(a)} Určte monotónnosť funkcie \(f(x)\) v okolí bodu \(x_0\). \\[1mm]
\textbf{(b)} Určte intervaly, na ktorých je funkcia \(g(x)\) klesajúca. \\[1mm]
\textbf{(c)} Určte počet stacionárnych bodov funkcie \(h(t)\). \par \columnbreak
\textbf{(d)} Nadobúda funkcia \(v(x)\) lokálny extrém? \\[1mm]
\textbf{(e)} Určte rozdiel hodnôt lokálneho maxima a lokálneho minima funkcie \(\varphi(x)\). \\[1mm]
\textbf{(f)} Určte súčin \(x\)-ových súradníc bodov globálneho maxima a globálneho minima funkcie \(\omega(x)\) na zadanom intervale. \\[1mm]
\end{multicols}
}
\end{center}
\end{minipage}
\\[4mm]
\begin{minipage}{0.79\linewidth}
\begin{center}
\begin{varwidth}{\linewidth}
\begin{enumerate}
\normalsize
\item $f(x)=\frac{x}{3}+\frac{3}{x} \enspace , \enspace x_0=-5$\quad \dotfill\; ???\;\dotfill \quad \text{klesajúca}
\item $g(x)=x^2\, e^{-6x+4}$\quad \dotfill\; ???\;\dotfill \quad $\left(-\infty , 0\right\rangle\enspace , \enspace \left\langle\frac{1}{3} , \infty\right)$
\item $h(t)=\cfrac{\ln t}{2t^5}$\quad \dotfill\; ???\;\dotfill \quad 2
\item $v(x)=-4-\sqrt[3]{(-3x-3)^2}$\quad \dotfill\; ???\;\dotfill \quad $\text{nie}$
\item $\varphi(x)=\cfrac{x^2+x+14}{x+2}$\quad \dotfill\; ???\;\dotfill \quad $-18$
\item $\omega(x)=2x^3-\frac{19}{2}x^2+3x-2 \enspace , \enspace x\in\left\langle0 , 5\right\rangle$\quad \dotfill\; ???\;\dotfill \quad $15$
\end{enumerate}
\end{varwidth}
\end{center}
\end{minipage}
\begin{minipage}{0.20\linewidth}
\begin{center}
{\Huge\bfseries 3.} \\[2mm]
\includegraphics[height=37mm]{images/braille.png}
{\small Braillovo písmeno}
\end{center}
\end{minipage}
\end{center}
\end{minipage}
&
\begin{minipage}[c][104.5mm][t]{0.5\linewidth}
\begin{center}
\vspace{7mm}
{\huge Monotónnosť a extrémy, skupina \textit{Delta $\delta$} -\romannumeral4}\\[5mm]
\begin{minipage}{0.95\linewidth}
\begin{center}
{\small \begin{multicols}{2} \raggedright
\textbf{(a)} Určte monotónnosť funkcie \(f(x)\) v okolí bodu \(x_0\). \\[1mm]
\textbf{(b)} Určte intervaly, na ktorých je funkcia \(g(x)\) klesajúca. \\[1mm]
\textbf{(c)} Určte počet stacionárnych bodov funkcie \(h(t)\). \par \columnbreak
\textbf{(d)} Nadobúda funkcia \(v(x)\) lokálny extrém? \\[1mm]
\textbf{(e)} Určte rozdiel hodnôt lokálneho maxima a lokálneho minima funkcie \(\varphi(x)\). \\[1mm]
\textbf{(f)} Určte súčin \(x\)-ových súradníc bodov globálneho maxima a globálneho minima funkcie \(\omega(x)\) na zadanom intervale. \\[1mm]
\end{multicols}
}
\end{center}
\end{minipage}
\\[4mm]
\begin{minipage}{0.79\linewidth}
\begin{center}
\begin{varwidth}{\linewidth}
\begin{enumerate}
\normalsize
\item $f(x)=\frac{x}{8}+\frac{8}{x} \enspace , \enspace x_0=6$\quad \dotfill\; ???\;\dotfill \quad \text{rastúca}
\item $g(x)=x^2\, e^{2x+5}$\quad \dotfill\; ???\;\dotfill \quad $\left\langle-1 , 0\right\rangle$
\item $h(t)=\cfrac{\ln t}{7t^3}$\quad \dotfill\; ???\;\dotfill \quad 1
\item $v(x)=3-\sqrt[5]{(-5x+1)^2}$\quad \dotfill\; ???\;\dotfill \quad $\text{nie}$
\item $\varphi(x)=\cfrac{x^2-4x+11}{x-5}$\quad \dotfill\; ???\;\dotfill \quad $-18$
\item $\omega(x)=\frac{1}{3}x^3+2x^2+3x-7 \enspace , \enspace x\in\left\langle-4 , 1\right\rangle$\quad \dotfill\; ???\;\dotfill \quad $-1$
\end{enumerate}
\end{varwidth}
\end{center}
\end{minipage}
\begin{minipage}{0.20\linewidth}
\begin{center}
{\Huge\bfseries 4.} \\[2mm]
\includegraphics[height=37mm]{images/braille.png}
{\small Braillovo písmeno}
\end{center}
\end{minipage}
\end{center}
\end{minipage}
%
\end{tabular}
\newpage
\thispagestyle{empty}
\begin{tabular}{c:c}
\begin{minipage}[c][104.5mm][t]{0.5\linewidth}
\begin{center}
\vspace{7mm}
{\huge Monotónnosť a extrémy, skupina \textit{Epsilon $\epsilon$} -\romannumeral1}\\[5mm]
\begin{minipage}{0.95\linewidth}
\begin{center}
{\small \begin{multicols}{2} \raggedright
\textbf{(a)} Určte monotónnosť funkcie \(f(x)\) v okolí bodu \(x_0\). \\[1mm]
\textbf{(b)} Určte intervaly, na ktorých je funkcia \(g(x)\) klesajúca. \\[1mm]
\textbf{(c)} Určte počet stacionárnych bodov funkcie \(h(t)\). \par \columnbreak
\textbf{(d)} Nadobúda funkcia \(v(x)\) lokálny extrém? \\[1mm]
\textbf{(e)} Určte rozdiel hodnôt lokálneho maxima a lokálneho minima funkcie \(\varphi(x)\). \\[1mm]
\textbf{(f)} Určte súčin \(x\)-ových súradníc bodov globálneho maxima a globálneho minima funkcie \(\omega(x)\) na zadanom intervale. \\[1mm]
\end{multicols}
}
\end{center}
\end{minipage}
\\[4mm]
\begin{minipage}{0.79\linewidth}
\begin{center}
\begin{varwidth}{\linewidth}
\begin{enumerate}
\normalsize
\item $f(x)=\frac{x}{4}+\frac{4}{x} \enspace , \enspace x_0=-6$\quad \dotfill\; ???\;\dotfill \quad \text{rastúca}
\item $g(x)=x^2\, e^{5x+3}$\quad \dotfill\; ???\;\dotfill \quad $\left(-\infty , 0\right\rangle\enspace , \enspace \left\langle-\frac{2}{5} , \infty\right)$
\item $h(t)=\cfrac{\ln t}{t^5}$\quad \dotfill\; ???\;\dotfill \quad 2
\item $v(x)=2-\sqrt[5]{(5x-5)^2}$\quad \dotfill\; ???\;\dotfill \quad $\text{áno}$
\item $\varphi(x)=\cfrac{x^2-2x+8}{x+2}$\quad \dotfill\; ???\;\dotfill \quad $-14$
\item $\omega(x)=3x^3-\frac{19}{2}x^2+2x+3 \enspace , \enspace x\in\left\langle0 , 4\right\rangle$\quad \dotfill\; ???\;\dotfill \quad $8$
\end{enumerate}
\end{varwidth}
\end{center}
\end{minipage}
\begin{minipage}{0.20\linewidth}
\begin{center}
{\Huge\bfseries 1.} \\[2mm]
\includegraphics[height=37mm]{images/braille.png}
{\small Braillovo písmeno}
\end{center}
\end{minipage}
\end{center}
\end{minipage}
&
\begin{minipage}[c][104.5mm][t]{0.5\linewidth}
\begin{center}
\vspace{7mm}
{\huge Monotónnosť a extrémy, skupina \textit{Epsilon $\epsilon$} -\romannumeral2}\\[5mm]
\begin{minipage}{0.95\linewidth}
\begin{center}
{\small \begin{multicols}{2} \raggedright
\textbf{(a)} Určte monotónnosť funkcie \(f(x)\) v okolí bodu \(x_0\). \\[1mm]
\textbf{(b)} Určte intervaly, na ktorých je funkcia \(g(x)\) klesajúca. \\[1mm]
\textbf{(c)} Určte počet stacionárnych bodov funkcie \(h(t)\). \par \columnbreak
\textbf{(d)} Nadobúda funkcia \(v(x)\) lokálny extrém? \\[1mm]
\textbf{(e)} Určte rozdiel hodnôt lokálneho maxima a lokálneho minima funkcie \(\varphi(x)\). \\[1mm]
\textbf{(f)} Určte súčin \(x\)-ových súradníc bodov globálneho maxima a globálneho minima funkcie \(\omega(x)\) na zadanom intervale. \\[1mm]
\end{multicols}
}
\end{center}
\end{minipage}
\\[4mm]
\begin{minipage}{0.79\linewidth}
\begin{center}
\begin{varwidth}{\linewidth}
\begin{enumerate}
\normalsize
\item $f(x)=\frac{x}{8}+\frac{8}{x} \enspace , \enspace x_0=2$\quad \dotfill\; ???\;\dotfill \quad \text{klesajúca}
\item $g(x)=x^2\, e^{4x+6}$\quad \dotfill\; ???\;\dotfill \quad $\left(-\infty , 0\right\rangle\enspace , \enspace \left\langle-\frac{1}{2} , \infty\right)$
\item $h(t)=\cfrac{\ln t}{2t^3}$\quad \dotfill\; ???\;\dotfill \quad 0
\item $v(x)=-3-\sqrt[7]{(-6x+8)^2}$\quad \dotfill\; ???\;\dotfill \quad $\text{nie}$
\item $\varphi(x)=\cfrac{x^2-x+23}{x-2}$\quad \dotfill\; ???\;\dotfill \quad $-20$
\item $\omega(x)=-\frac{5}{3}x^3+8x^2-3x+1 \enspace , \enspace x\in\left\langle-2 , 4\right\rangle$\quad \dotfill\; ???\;\dotfill \quad $\frac{4}{25}$
\end{enumerate}
\end{varwidth}
\end{center}
\end{minipage}
\begin{minipage}{0.20\linewidth}
\begin{center}
{\Huge\bfseries 2.} \\[2mm]
\includegraphics[height=37mm]{images/braille.png}
{\small Braillovo písmeno}
\end{center}
\end{minipage}
\end{center}
\end{minipage}
\\ \hdashline
\begin{minipage}[c][104.5mm][t]{0.5\linewidth}
\begin{center}
\vspace{7mm}
{\huge Monotónnosť a extrémy, skupina \textit{Epsilon $\epsilon$} -\romannumeral3}\\[5mm]
\begin{minipage}{0.95\linewidth}
\begin{center}
{\small \begin{multicols}{2} \raggedright
\textbf{(a)} Určte monotónnosť funkcie \(f(x)\) v okolí bodu \(x_0\). \\[1mm]
\textbf{(b)} Určte intervaly, na ktorých je funkcia \(g(x)\) klesajúca. \\[1mm]
\textbf{(c)} Určte počet stacionárnych bodov funkcie \(h(t)\). \par \columnbreak
\textbf{(d)} Nadobúda funkcia \(v(x)\) lokálny extrém? \\[1mm]
\textbf{(e)} Určte rozdiel hodnôt lokálneho maxima a lokálneho minima funkcie \(\varphi(x)\). \\[1mm]
\textbf{(f)} Určte súčin \(x\)-ových súradníc bodov globálneho maxima a globálneho minima funkcie \(\omega(x)\) na zadanom intervale. \\[1mm]
\end{multicols}
}
\end{center}
\end{minipage}
\\[4mm]
\begin{minipage}{0.79\linewidth}
\begin{center}
\begin{varwidth}{\linewidth}
\begin{enumerate}
\normalsize
\item $f(x)=\frac{x}{7}+\frac{7}{x} \enspace , \enspace x_0=5$\quad \dotfill\; ???\;\dotfill \quad \text{rastúca}
\item $g(x)=x^2\, e^{2x-2}$\quad \dotfill\; ???\;\dotfill \quad $\left\langle-1 , 0\right\rangle$
\item $h(t)=\cfrac{\ln t}{2t^3}$\quad \dotfill\; ???\;\dotfill \quad 1
\item $v(x)=2-\sqrt[3]{(-5x+1)^2}$\quad \dotfill\; ???\;\dotfill \quad $\text{nie}$
\item $\varphi(x)=\cfrac{x^2-4x-5}{x+3}$\quad \dotfill\; ???\;\dotfill \quad $-18$
\item $\omega(x)=\frac{2}{3}x^3-\frac{17}{2}x^2+30x+3 \enspace , \enspace x\in\left\langle2 , 8\right\rangle$\quad \dotfill\; ???\;\dotfill \quad $48$
\end{enumerate}
\end{varwidth}
\end{center}
\end{minipage}
\begin{minipage}{0.20\linewidth}
\begin{center}
{\Huge\bfseries 3.} \\[2mm]
\includegraphics[height=37mm]{images/braille.png}
{\small Braillovo písmeno}
\end{center}
\end{minipage}
\end{center}
\end{minipage}
&
\begin{minipage}[c][104.5mm][t]{0.5\linewidth}
\begin{center}
\vspace{7mm}
{\huge Monotónnosť a extrémy, skupina \textit{Epsilon $\epsilon$} -\romannumeral4}\\[5mm]
\begin{minipage}{0.95\linewidth}
\begin{center}
{\small \begin{multicols}{2} \raggedright
\textbf{(a)} Určte monotónnosť funkcie \(f(x)\) v okolí bodu \(x_0\). \\[1mm]
\textbf{(b)} Určte intervaly, na ktorých je funkcia \(g(x)\) klesajúca. \\[1mm]
\textbf{(c)} Určte počet stacionárnych bodov funkcie \(h(t)\). \par \columnbreak
\textbf{(d)} Nadobúda funkcia \(v(x)\) lokálny extrém? \\[1mm]
\textbf{(e)} Určte rozdiel hodnôt lokálneho maxima a lokálneho minima funkcie \(\varphi(x)\). \\[1mm]
\textbf{(f)} Určte súčin \(x\)-ových súradníc bodov globálneho maxima a globálneho minima funkcie \(\omega(x)\) na zadanom intervale. \\[1mm]
\end{multicols}
}
\end{center}
\end{minipage}
\\[4mm]
\begin{minipage}{0.79\linewidth}
\begin{center}
\begin{varwidth}{\linewidth}
\begin{enumerate}
\normalsize
\item $f(x)=\frac{x}{7}+\frac{7}{x} \enspace , \enspace x_0=-12$\quad \dotfill\; ???\;\dotfill \quad \text{rastúca}
\item $g(x)=x^2\, e^{-x-2}$\quad \dotfill\; ???\;\dotfill \quad $\left(-\infty , 0\right\rangle\enspace , \enspace \left\langle2 , \infty\right)$
\item $h(t)=\cfrac{\ln t}{t^3}$\quad \dotfill\; ???\;\dotfill \quad 2
\item $v(x)=-6-\sqrt[5]{(6x+4)^2}$\quad \dotfill\; ???\;\dotfill \quad $\text{áno}$
\item $\varphi(x)=\cfrac{x^2+x-21}{x-5}$\quad \dotfill\; ???\;\dotfill \quad $-12$
\item $\omega(x)=-\frac{1}{3}x^3-6x^2-35x-5 \enspace , \enspace x\in\left\langle-10 , -4\right\rangle$\quad \dotfill\; ???\;\dotfill \quad $28$
\end{enumerate}
\end{varwidth}
\end{center}
\end{minipage}
\begin{minipage}{0.20\linewidth}
\begin{center}
{\Huge\bfseries 4.} \\[2mm]
\includegraphics[height=37mm]{images/braille.png}
{\small Braillovo písmeno}
\end{center}
\end{minipage}
\end{center}
\end{minipage}
%
\end{tabular}
\newpage
\thispagestyle{empty}
\begin{tabular}{c:c}
\begin{minipage}[c][104.5mm][t]{0.5\linewidth}
\begin{center}
\vspace{7mm}
{\huge Monotónnosť a extrémy, skupina \textit{Zeta $\zeta$} -\romannumeral1}\\[5mm]
\begin{minipage}{0.95\linewidth}
\begin{center}
{\small \begin{multicols}{2} \raggedright
\textbf{(a)} Určte monotónnosť funkcie \(f(x)\) v okolí bodu \(x_0\). \\[1mm]
\textbf{(b)} Určte intervaly, na ktorých je funkcia \(g(x)\) klesajúca. \\[1mm]
\textbf{(c)} Určte počet stacionárnych bodov funkcie \(h(t)\). \par \columnbreak
\textbf{(d)} Nadobúda funkcia \(v(x)\) lokálny extrém? \\[1mm]
\textbf{(e)} Určte rozdiel hodnôt lokálneho maxima a lokálneho minima funkcie \(\varphi(x)\). \\[1mm]
\textbf{(f)} Určte súčin \(x\)-ových súradníc bodov globálneho maxima a globálneho minima funkcie \(\omega(x)\) na zadanom intervale. \\[1mm]
\end{multicols}
}
\end{center}
\end{minipage}
\\[4mm]
\begin{minipage}{0.79\linewidth}
\begin{center}
\begin{varwidth}{\linewidth}
\begin{enumerate}
\normalsize
\item $f(x)=\frac{x}{2}+\frac{2}{x} \enspace , \enspace x_0=-4$\quad \dotfill\; ???\;\dotfill \quad \text{klesajúca}
\item $g(x)=x^2\, e^{3x+3}$\quad \dotfill\; ???\;\dotfill \quad $\left\langle-\frac{2}{3} , 0\right\rangle$
\item $h(t)=\cfrac{\ln t}{3t^3}$\quad \dotfill\; ???\;\dotfill \quad 1
\item $v(x)=-8-\sqrt[3]{(x-6)^2}$\quad \dotfill\; ???\;\dotfill \quad $\text{nie}$
\item $\varphi(x)=\cfrac{x^2-6x-15}{x+2}$\quad \dotfill\; ???\;\dotfill \quad $-4$
\item $\omega(x)=3x^3-\frac{19}{2}x^2+2x+4 \enspace , \enspace x\in\left\langle0 , 4\right\rangle$\quad \dotfill\; ???\;\dotfill \quad $8$
\end{enumerate}
\end{varwidth}
\end{center}
\end{minipage}
\begin{minipage}{0.20\linewidth}
\begin{center}
{\Huge\bfseries 1.} \\[2mm]
\includegraphics[height=37mm]{images/braille.png}
{\small Braillovo písmeno}
\end{center}
\end{minipage}
\end{center}
\end{minipage}
&
\begin{minipage}[c][104.5mm][t]{0.5\linewidth}
\begin{center}
\vspace{7mm}
{\huge Monotónnosť a extrémy, skupina \textit{Zeta $\zeta$} -\romannumeral2}\\[5mm]
\begin{minipage}{0.95\linewidth}
\begin{center}
{\small \begin{multicols}{2} \raggedright
\textbf{(a)} Určte monotónnosť funkcie \(f(x)\) v okolí bodu \(x_0\). \\[1mm]
\textbf{(b)} Určte intervaly, na ktorých je funkcia \(g(x)\) klesajúca. \\[1mm]
\textbf{(c)} Určte počet stacionárnych bodov funkcie \(h(t)\). \par \columnbreak
\textbf{(d)} Nadobúda funkcia \(v(x)\) lokálny extrém? \\[1mm]
\textbf{(e)} Určte rozdiel hodnôt lokálneho maxima a lokálneho minima funkcie \(\varphi(x)\). \\[1mm]
\textbf{(f)} Určte súčin \(x\)-ových súradníc bodov globálneho maxima a globálneho minima funkcie \(\omega(x)\) na zadanom intervale. \\[1mm]
\end{multicols}
}
\end{center}
\end{minipage}
\\[4mm]
\begin{minipage}{0.79\linewidth}
\begin{center}
\begin{varwidth}{\linewidth}
\begin{enumerate}
\normalsize
\item $f(x)=\frac{x}{3}+\frac{3}{x} \enspace , \enspace x_0=-4$\quad \dotfill\; ???\;\dotfill \quad \text{rastúca}
\item $g(x)=x^2\, e^{-3x-1}$\quad \dotfill\; ???\;\dotfill \quad $\left\langle\frac{2}{3} , 0\right\rangle$
\item $h(t)=\cfrac{\ln t}{8t^5}$\quad \dotfill\; ???\;\dotfill \quad 0
\item $v(x)=-2-\sqrt[7]{(-3x+8)^2}$\quad \dotfill\; ???\;\dotfill \quad $\text{nie}$
\item $\varphi(x)=\cfrac{x^2-x+4}{x+3}$\quad \dotfill\; ???\;\dotfill \quad $-14$
\item $\omega(x)=-\frac{1}{3}x^3-3x^2-8x+4 \enspace , \enspace x\in\left\langle-7 , -1\right\rangle$\quad \dotfill\; ???\;\dotfill \quad $-12$
\end{enumerate}
\end{varwidth}
\end{center}
\end{minipage}
\begin{minipage}{0.20\linewidth}
\begin{center}
{\Huge\bfseries 2.} \\[2mm]
\includegraphics[height=37mm]{images/braille.png}
{\small Braillovo písmeno}
\end{center}
\end{minipage}
\end{center}
\end{minipage}
\\ \hdashline
\begin{minipage}[c][104.5mm][t]{0.5\linewidth}
\begin{center}
\vspace{7mm}
{\huge Monotónnosť a extrémy, skupina \textit{Zeta $\zeta$} -\romannumeral3}\\[5mm]
\begin{minipage}{0.95\linewidth}
\begin{center}
{\small \begin{multicols}{2} \raggedright
\textbf{(a)} Určte monotónnosť funkcie \(f(x)\) v okolí bodu \(x_0\). \\[1mm]
\textbf{(b)} Určte intervaly, na ktorých je funkcia \(g(x)\) klesajúca. \\[1mm]
\textbf{(c)} Určte počet stacionárnych bodov funkcie \(h(t)\). \par \columnbreak
\textbf{(d)} Nadobúda funkcia \(v(x)\) lokálny extrém? \\[1mm]
\textbf{(e)} Určte rozdiel hodnôt lokálneho maxima a lokálneho minima funkcie \(\varphi(x)\). \\[1mm]
\textbf{(f)} Určte súčin \(x\)-ových súradníc bodov globálneho maxima a globálneho minima funkcie \(\omega(x)\) na zadanom intervale. \\[1mm]
\end{multicols}
}
\end{center}
\end{minipage}
\\[4mm]
\begin{minipage}{0.79\linewidth}
\begin{center}
\begin{varwidth}{\linewidth}
\begin{enumerate}
\normalsize
\item $f(x)=\frac{x}{4}+\frac{4}{x} \enspace , \enspace x_0=-7$\quad \dotfill\; ???\;\dotfill \quad \text{rastúca}
\item $g(x)=x^2\, e^{5x-4}$\quad \dotfill\; ???\;\dotfill \quad $\left(-\infty , 0\right\rangle\enspace , \enspace \left\langle-\frac{2}{5} , \infty\right)$
\item $h(t)=\cfrac{\ln t}{4t^5}$\quad \dotfill\; ???\;\dotfill \quad 1
\item $v(x)=-4-\sqrt[7]{(-5x+2)^2}$\quad \dotfill\; ???\;\dotfill \quad $\text{áno}$
\item $\varphi(x)=\cfrac{x^2-x+7}{x-2}$\quad \dotfill\; ???\;\dotfill \quad $-14$
\item $\omega(x)=\frac{1}{3}x^3-3x^2+8x-5 \enspace , \enspace x\in\left\langle1 , 6\right\rangle$\quad \dotfill\; ???\;\dotfill \quad $24$
\end{enumerate}
\end{varwidth}
\end{center}
\end{minipage}
\begin{minipage}{0.20\linewidth}
\begin{center}
{\Huge\bfseries 3.} \\[2mm]
\includegraphics[height=37mm]{images/braille.png}
{\small Braillovo písmeno}
\end{center}
\end{minipage}
\end{center}
\end{minipage}
&
\begin{minipage}[c][104.5mm][t]{0.5\linewidth}
\begin{center}
\vspace{7mm}
{\huge Monotónnosť a extrémy, skupina \textit{Zeta $\zeta$} -\romannumeral4}\\[5mm]
\begin{minipage}{0.95\linewidth}
\begin{center}
{\small \begin{multicols}{2} \raggedright
\textbf{(a)} Určte monotónnosť funkcie \(f(x)\) v okolí bodu \(x_0\). \\[1mm]
\textbf{(b)} Určte intervaly, na ktorých je funkcia \(g(x)\) klesajúca. \\[1mm]
\textbf{(c)} Určte počet stacionárnych bodov funkcie \(h(t)\). \par \columnbreak
\textbf{(d)} Nadobúda funkcia \(v(x)\) lokálny extrém? \\[1mm]
\textbf{(e)} Určte rozdiel hodnôt lokálneho maxima a lokálneho minima funkcie \(\varphi(x)\). \\[1mm]
\textbf{(f)} Určte súčin \(x\)-ových súradníc bodov globálneho maxima a globálneho minima funkcie \(\omega(x)\) na zadanom intervale. \\[1mm]
\end{multicols}
}
\end{center}
\end{minipage}
\\[4mm]
\begin{minipage}{0.79\linewidth}
\begin{center}
\begin{varwidth}{\linewidth}
\begin{enumerate}
\normalsize
\item $f(x)=\frac{x}{3}+\frac{3}{x} \enspace , \enspace x_0=1$\quad \dotfill\; ???\;\dotfill \quad \text{rastúca}
\item $g(x)=x^2\, e^{4x-2}$\quad \dotfill\; ???\;\dotfill \quad $\left\langle-\frac{1}{2} , 0\right\rangle$
\item $h(t)=\cfrac{\ln t}{7t^4}$\quad \dotfill\; ???\;\dotfill \quad 2
\item $v(x)=-2-\sqrt[7]{(8x+1)^2}$\quad \dotfill\; ???\;\dotfill \quad $\text{nie}$
\item $\varphi(x)=\cfrac{x^2+x+25}{x+1}$\quad \dotfill\; ???\;\dotfill \quad $-18$
\item $\omega(x)=-x^3+\frac{7}{2}x^2+6x-4 \enspace , \enspace x\in\left\langle-3 , 4\right\rangle$\quad \dotfill\; ???\;\dotfill \quad $2$
\end{enumerate}
\end{varwidth}
\end{center}
\end{minipage}
\begin{minipage}{0.20\linewidth}
\begin{center}
{\Huge\bfseries 4.} \\[2mm]
\includegraphics[height=37mm]{images/braille.png}
{\small Braillovo písmeno}
\end{center}
\end{minipage}
\end{center}
\end{minipage}
%
\end{tabular}
\newpage
\thispagestyle{empty}
\begin{tabular}{c:c}
\begin{minipage}[c][104.5mm][t]{0.5\linewidth}
\begin{center}
\vspace{7mm}
{\huge Monotónnosť a extrémy, skupina \textit{Eta $\eta$} -\romannumeral1}\\[5mm]
\begin{minipage}{0.95\linewidth}
\begin{center}
{\small \begin{multicols}{2} \raggedright
\textbf{(a)} Určte monotónnosť funkcie \(f(x)\) v okolí bodu \(x_0\). \\[1mm]
\textbf{(b)} Určte intervaly, na ktorých je funkcia \(g(x)\) klesajúca. \\[1mm]
\textbf{(c)} Určte počet stacionárnych bodov funkcie \(h(t)\). \par \columnbreak
\textbf{(d)} Nadobúda funkcia \(v(x)\) lokálny extrém? \\[1mm]
\textbf{(e)} Určte rozdiel hodnôt lokálneho maxima a lokálneho minima funkcie \(\varphi(x)\). \\[1mm]
\textbf{(f)} Určte súčin \(x\)-ových súradníc bodov globálneho maxima a globálneho minima funkcie \(\omega(x)\) na zadanom intervale. \\[1mm]
\end{multicols}
}
\end{center}
\end{minipage}
\\[4mm]
\begin{minipage}{0.79\linewidth}
\begin{center}
\begin{varwidth}{\linewidth}
\begin{enumerate}
\normalsize
\item $f(x)=\frac{x}{6}+\frac{6}{x} \enspace , \enspace x_0=3$\quad \dotfill\; ???\;\dotfill \quad \text{klesajúca}
\item $g(x)=x^2\, e^{6x+3}$\quad \dotfill\; ???\;\dotfill \quad $\left(-\infty , 0\right\rangle\enspace , \enspace \left\langle-\frac{1}{3} , \infty\right)$
\item $h(t)=\cfrac{\ln t}{2t^3}$\quad \dotfill\; ???\;\dotfill \quad 1
\item $v(x)=-3-\sqrt[3]{(x+7)^2}$\quad \dotfill\; ???\;\dotfill \quad $\text{nie}$
\item $\varphi(x)=\cfrac{x^2+x+14}{x+2}$\quad \dotfill\; ???\;\dotfill \quad $-14$
\item $\omega(x)=\frac{2}{3}x^3-\frac{3}{2}x^2-2x-5 \enspace , \enspace x\in\left\langle-1 , 4\right\rangle$\quad \dotfill\; ???\;\dotfill \quad $8$
\end{enumerate}
\end{varwidth}
\end{center}
\end{minipage}
\begin{minipage}{0.20\linewidth}
\begin{center}
{\Huge\bfseries 1.} \\[2mm]
\includegraphics[height=37mm]{images/braille.png}
{\small Braillovo písmeno}
\end{center}
\end{minipage}
\end{center}
\end{minipage}
&
\begin{minipage}[c][104.5mm][t]{0.5\linewidth}
\begin{center}
\vspace{7mm}
{\huge Monotónnosť a extrémy, skupina \textit{Eta $\eta$} -\romannumeral2}\\[5mm]
\begin{minipage}{0.95\linewidth}
\begin{center}
{\small \begin{multicols}{2} \raggedright
\textbf{(a)} Určte monotónnosť funkcie \(f(x)\) v okolí bodu \(x_0\). \\[1mm]
\textbf{(b)} Určte intervaly, na ktorých je funkcia \(g(x)\) klesajúca. \\[1mm]
\textbf{(c)} Určte počet stacionárnych bodov funkcie \(h(t)\). \par \columnbreak
\textbf{(d)} Nadobúda funkcia \(v(x)\) lokálny extrém? \\[1mm]
\textbf{(e)} Určte rozdiel hodnôt lokálneho maxima a lokálneho minima funkcie \(\varphi(x)\). \\[1mm]
\textbf{(f)} Určte súčin \(x\)-ových súradníc bodov globálneho maxima a globálneho minima funkcie \(\omega(x)\) na zadanom intervale. \\[1mm]
\end{multicols}
}
\end{center}
\end{minipage}
\\[4mm]
\begin{minipage}{0.79\linewidth}
\begin{center}
\begin{varwidth}{\linewidth}
\begin{enumerate}
\normalsize
\item $f(x)=\frac{x}{5}+\frac{5}{x} \enspace , \enspace x_0=2$\quad \dotfill\; ???\;\dotfill \quad \text{klesajúca}
\item $g(x)=x^2\, e^{-2x+2}$\quad \dotfill\; ???\;\dotfill \quad $\left(-\infty , 0\right\rangle\enspace , \enspace \left\langle1 , \infty\right)$
\item $h(t)=\cfrac{\ln t}{7t^3}$\quad \dotfill\; ???\;\dotfill \quad 1
\item $v(x)=-6-\sqrt[5]{(-5x-3)^2}$\quad \dotfill\; ???\;\dotfill \quad $\text{nie}$
\item $\varphi(x)=\cfrac{x^2-4x+20}{x+1}$\quad \dotfill\; ???\;\dotfill \quad $-20$
\item $\omega(x)=-x^3+2x^2+4x-2 \enspace , \enspace x\in\left\langle-3 , 3\right\rangle$\quad \dotfill\; ???\;\dotfill \quad $-\frac{2}{9}$
\end{enumerate}
\end{varwidth}
\end{center}
\end{minipage}
\begin{minipage}{0.20\linewidth}
\begin{center}
{\Huge\bfseries 2.} \\[2mm]
\includegraphics[height=37mm]{images/braille.png}
{\small Braillovo písmeno}
\end{center}
\end{minipage}
\end{center}
\end{minipage}
\\ \hdashline
\begin{minipage}[c][104.5mm][t]{0.5\linewidth}
\begin{center}
\vspace{7mm}
{\huge Monotónnosť a extrémy, skupina \textit{Eta $\eta$} -\romannumeral3}\\[5mm]
\begin{minipage}{0.95\linewidth}
\begin{center}
{\small \begin{multicols}{2} \raggedright
\textbf{(a)} Určte monotónnosť funkcie \(f(x)\) v okolí bodu \(x_0\). \\[1mm]
\textbf{(b)} Určte intervaly, na ktorých je funkcia \(g(x)\) klesajúca. \\[1mm]
\textbf{(c)} Určte počet stacionárnych bodov funkcie \(h(t)\). \par \columnbreak
\textbf{(d)} Nadobúda funkcia \(v(x)\) lokálny extrém? \\[1mm]
\textbf{(e)} Určte rozdiel hodnôt lokálneho maxima a lokálneho minima funkcie \(\varphi(x)\). \\[1mm]
\textbf{(f)} Určte súčin \(x\)-ových súradníc bodov globálneho maxima a globálneho minima funkcie \(\omega(x)\) na zadanom intervale. \\[1mm]
\end{multicols}
}
\end{center}
\end{minipage}
\\[4mm]
\begin{minipage}{0.79\linewidth}
\begin{center}
\begin{varwidth}{\linewidth}
\begin{enumerate}
\normalsize
\item $f(x)=\frac{x}{6}+\frac{6}{x} \enspace , \enspace x_0=5$\quad \dotfill\; ???\;\dotfill \quad \text{klesajúca}
\item $g(x)=x^2\, e^{-8x-2}$\quad \dotfill\; ???\;\dotfill \quad $\left\langle\frac{1}{4} , 0\right\rangle$
\item $h(t)=\cfrac{\ln t}{t^5}$\quad \dotfill\; ???\;\dotfill \quad 2
\item $v(x)=-4-\sqrt[5]{(-8x-7)^2}$\quad \dotfill\; ???\;\dotfill \quad $\text{nie}$
\item $\varphi(x)=\cfrac{x^2-6x+9}{x+1}$\quad \dotfill\; ???\;\dotfill \quad $-14$
\item $\omega(x)=-\frac{5}{3}x^3+\frac{11}{2}x^2+12x+8 \enspace , \enspace x\in\left\langle-3 , 4\right\rangle$\quad \dotfill\; ???\;\dotfill \quad $\frac{4}{25}$
\end{enumerate}
\end{varwidth}
\end{center}
\end{minipage}
\begin{minipage}{0.20\linewidth}
\begin{center}
{\Huge\bfseries 3.} \\[2mm]
\includegraphics[height=37mm]{images/braille.png}
{\small Braillovo písmeno}
\end{center}
\end{minipage}
\end{center}
\end{minipage}
&
\begin{minipage}[c][104.5mm][t]{0.5\linewidth}
\begin{center}
\vspace{7mm}
{\huge Monotónnosť a extrémy, skupina \textit{Eta $\eta$} -\romannumeral4}\\[5mm]
\begin{minipage}{0.95\linewidth}
\begin{center}
{\small \begin{multicols}{2} \raggedright
\textbf{(a)} Určte monotónnosť funkcie \(f(x)\) v okolí bodu \(x_0\). \\[1mm]
\textbf{(b)} Určte intervaly, na ktorých je funkcia \(g(x)\) klesajúca. \\[1mm]
\textbf{(c)} Určte počet stacionárnych bodov funkcie \(h(t)\). \par \columnbreak
\textbf{(d)} Nadobúda funkcia \(v(x)\) lokálny extrém? \\[1mm]
\textbf{(e)} Určte rozdiel hodnôt lokálneho maxima a lokálneho minima funkcie \(\varphi(x)\). \\[1mm]
\textbf{(f)} Určte súčin \(x\)-ových súradníc bodov globálneho maxima a globálneho minima funkcie \(\omega(x)\) na zadanom intervale. \\[1mm]
\end{multicols}
}
\end{center}
\end{minipage}
\\[4mm]
\begin{minipage}{0.79\linewidth}
\begin{center}
\begin{varwidth}{\linewidth}
\begin{enumerate}
\normalsize
\item $f(x)=\frac{x}{5}+\frac{5}{x} \enspace , \enspace x_0=2$\quad \dotfill\; ???\;\dotfill \quad \text{klesajúca}
\item $g(x)=x^2\, e^{-4x+2}$\quad \dotfill\; ???\;\dotfill \quad $\left\langle\frac{1}{2} , 0\right\rangle$
\item $h(t)=\cfrac{\ln t}{5t^3}$\quad \dotfill\; ???\;\dotfill \quad 1
\item $v(x)=2-\sqrt[7]{(-2x+6)^2}$\quad \dotfill\; ???\;\dotfill \quad $\text{nie}$
\item $\varphi(x)=\cfrac{x^2-8x+37}{x-2}$\quad \dotfill\; ???\;\dotfill \quad $-22$
\item $\omega(x)=\frac{1}{3}x^3-\frac{7}{2}x^2+10x-1 \enspace , \enspace x\in\left\langle1 , 7\right\rangle$\quad \dotfill\; ???\;\dotfill \quad $30$
\end{enumerate}
\end{varwidth}
\end{center}
\end{minipage}
\begin{minipage}{0.20\linewidth}
\begin{center}
{\Huge\bfseries 4.} \\[2mm]
\includegraphics[height=37mm]{images/braille.png}
{\small Braillovo písmeno}
\end{center}
\end{minipage}
\end{center}
\end{minipage}
%
\end{tabular}
\newpage
\thispagestyle{empty}
\begin{tabular}{c:c}
\begin{minipage}[c][104.5mm][t]{0.5\linewidth}
\begin{center}
\vspace{7mm}
{\huge Monotónnosť a extrémy, skupina \textit{Theta $\theta$} -\romannumeral1}\\[5mm]
\begin{minipage}{0.95\linewidth}
\begin{center}
{\small \begin{multicols}{2} \raggedright
\textbf{(a)} Určte monotónnosť funkcie \(f(x)\) v okolí bodu \(x_0\). \\[1mm]
\textbf{(b)} Určte intervaly, na ktorých je funkcia \(g(x)\) klesajúca. \\[1mm]
\textbf{(c)} Určte počet stacionárnych bodov funkcie \(h(t)\). \par \columnbreak
\textbf{(d)} Nadobúda funkcia \(v(x)\) lokálny extrém? \\[1mm]
\textbf{(e)} Určte rozdiel hodnôt lokálneho maxima a lokálneho minima funkcie \(\varphi(x)\). \\[1mm]
\textbf{(f)} Určte súčin \(x\)-ových súradníc bodov globálneho maxima a globálneho minima funkcie \(\omega(x)\) na zadanom intervale. \\[1mm]
\end{multicols}
}
\end{center}
\end{minipage}
\\[4mm]
\begin{minipage}{0.79\linewidth}
\begin{center}
\begin{varwidth}{\linewidth}
\begin{enumerate}
\normalsize
\item $f(x)=\frac{x}{7}+\frac{7}{x} \enspace , \enspace x_0=1$\quad \dotfill\; ???\;\dotfill \quad \text{klesajúca}
\item $g(x)=x^2\, e^{-5x+3}$\quad \dotfill\; ???\;\dotfill \quad $\left(-\infty , 0\right\rangle\enspace , \enspace \left\langle\frac{2}{5} , \infty\right)$
\item $h(t)=\cfrac{\ln t}{3t^3}$\quad \dotfill\; ???\;\dotfill \quad 1
\item $v(x)=1-\sqrt[7]{(-7x+6)^2}$\quad \dotfill\; ???\;\dotfill \quad $\text{nie}$
\item $\varphi(x)=\cfrac{x^2+x+47}{x+2}$\quad \dotfill\; ???\;\dotfill \quad $-28$
\item $\omega(x)=-\frac{5}{3}x^3+\frac{11}{2}x^2+12x-4 \enspace , \enspace x\in\left\langle-3 , 4\right\rangle$\quad \dotfill\; ???\;\dotfill \quad $-\frac{12}{25}$
\end{enumerate}
\end{varwidth}
\end{center}
\end{minipage}
\begin{minipage}{0.20\linewidth}
\begin{center}
{\Huge\bfseries 1.} \\[2mm]
\includegraphics[height=37mm]{images/braille.png}
{\small Braillovo písmeno}
\end{center}
\end{minipage}
\end{center}
\end{minipage}
&
\begin{minipage}[c][104.5mm][t]{0.5\linewidth}
\begin{center}
\vspace{7mm}
{\huge Monotónnosť a extrémy, skupina \textit{Theta $\theta$} -\romannumeral2}\\[5mm]
\begin{minipage}{0.95\linewidth}
\begin{center}
{\small \begin{multicols}{2} \raggedright
\textbf{(a)} Určte monotónnosť funkcie \(f(x)\) v okolí bodu \(x_0\). \\[1mm]
\textbf{(b)} Určte intervaly, na ktorých je funkcia \(g(x)\) klesajúca. \\[1mm]
\textbf{(c)} Určte počet stacionárnych bodov funkcie \(h(t)\). \par \columnbreak
\textbf{(d)} Nadobúda funkcia \(v(x)\) lokálny extrém? \\[1mm]
\textbf{(e)} Určte rozdiel hodnôt lokálneho maxima a lokálneho minima funkcie \(\varphi(x)\). \\[1mm]
\textbf{(f)} Určte súčin \(x\)-ových súradníc bodov globálneho maxima a globálneho minima funkcie \(\omega(x)\) na zadanom intervale. \\[1mm]
\end{multicols}
}
\end{center}
\end{minipage}
\\[4mm]
\begin{minipage}{0.79\linewidth}
\begin{center}
\begin{varwidth}{\linewidth}
\begin{enumerate}
\normalsize
\item $f(x)=\frac{x}{4}+\frac{4}{x} \enspace , \enspace x_0=-5$\quad \dotfill\; ???\;\dotfill \quad \text{rastúca}
\item $g(x)=x^2\, e^{2x-3}$\quad \dotfill\; ???\;\dotfill \quad $\left(-\infty , 0\right\rangle\enspace , \enspace \left\langle-1 , \infty\right)$
\item $h(t)=\cfrac{\ln t}{4t^3}$\quad \dotfill\; ???\;\dotfill \quad 0
\item $v(x)=-2-\sqrt[3]{(4x+8)^2}$\quad \dotfill\; ???\;\dotfill \quad $\text{nie}$
\item $\varphi(x)=\cfrac{x^2+x-1}{x+2}$\quad \dotfill\; ???\;\dotfill \quad $-4$
\item $\omega(x)=-3x^3+\frac{19}{2}x^2+24x-8 \enspace , \enspace x\in\left\langle-3 , 4\right\rangle$\quad \dotfill\; ???\;\dotfill \quad $-\frac{56}{81}$
\end{enumerate}
\end{varwidth}
\end{center}
\end{minipage}
\begin{minipage}{0.20\linewidth}
\begin{center}
{\Huge\bfseries 2.} \\[2mm]
\includegraphics[height=37mm]{images/braille.png}
{\small Braillovo písmeno}
\end{center}
\end{minipage}
\end{center}
\end{minipage}
\\ \hdashline
\begin{minipage}[c][104.5mm][t]{0.5\linewidth}
\begin{center}
\vspace{7mm}
{\huge Monotónnosť a extrémy, skupina \textit{Theta $\theta$} -\romannumeral3}\\[5mm]
\begin{minipage}{0.95\linewidth}
\begin{center}
{\small \begin{multicols}{2} \raggedright
\textbf{(a)} Určte monotónnosť funkcie \(f(x)\) v okolí bodu \(x_0\). \\[1mm]
\textbf{(b)} Určte intervaly, na ktorých je funkcia \(g(x)\) klesajúca. \\[1mm]
\textbf{(c)} Určte počet stacionárnych bodov funkcie \(h(t)\). \par \columnbreak
\textbf{(d)} Nadobúda funkcia \(v(x)\) lokálny extrém? \\[1mm]
\textbf{(e)} Určte rozdiel hodnôt lokálneho maxima a lokálneho minima funkcie \(\varphi(x)\). \\[1mm]
\textbf{(f)} Určte súčin \(x\)-ových súradníc bodov globálneho maxima a globálneho minima funkcie \(\omega(x)\) na zadanom intervale. \\[1mm]
\end{multicols}
}
\end{center}
\end{minipage}
\\[4mm]
\begin{minipage}{0.79\linewidth}
\begin{center}
\begin{varwidth}{\linewidth}
\begin{enumerate}
\normalsize
\item $f(x)=\frac{x}{6}+\frac{6}{x} \enspace , \enspace x_0=4$\quad \dotfill\; ???\;\dotfill \quad \text{klesajúca}
\item $g(x)=x^2\, e^{-6x+2}$\quad \dotfill\; ???\;\dotfill \quad $\left(-\infty , 0\right\rangle\enspace , \enspace \left\langle\frac{1}{3} , \infty\right)$
\item $h(t)=\cfrac{\ln t}{6t^4}$\quad \dotfill\; ???\;\dotfill \quad 1
\item $v(x)=-2-\sqrt[3]{(7x+5)^2}$\quad \dotfill\; ???\;\dotfill \quad $\text{nie}$
\item $\varphi(x)=\cfrac{x^2+6x-54}{x-5}$\quad \dotfill\; ???\;\dotfill \quad $-6$
\item $\omega(x)=-\frac{1}{3}x^3-\frac{7}{2}x^2-6x-1 \enspace , \enspace x\in\left\langle-9 , 0\right\rangle$\quad \dotfill\; ???\;\dotfill \quad $54$
\end{enumerate}
\end{varwidth}
\end{center}
\end{minipage}
\begin{minipage}{0.20\linewidth}
\begin{center}
{\Huge\bfseries 3.} \\[2mm]
\includegraphics[height=37mm]{images/braille.png}
{\small Braillovo písmeno}
\end{center}
\end{minipage}
\end{center}
\end{minipage}
&
\begin{minipage}[c][104.5mm][t]{0.5\linewidth}
\begin{center}
\vspace{7mm}
{\huge Monotónnosť a extrémy, skupina \textit{Theta $\theta$} -\romannumeral4}\\[5mm]
\begin{minipage}{0.95\linewidth}
\begin{center}
{\small \begin{multicols}{2} \raggedright
\textbf{(a)} Určte monotónnosť funkcie \(f(x)\) v okolí bodu \(x_0\). \\[1mm]
\textbf{(b)} Určte intervaly, na ktorých je funkcia \(g(x)\) klesajúca. \\[1mm]
\textbf{(c)} Určte počet stacionárnych bodov funkcie \(h(t)\). \par \columnbreak
\textbf{(d)} Nadobúda funkcia \(v(x)\) lokálny extrém? \\[1mm]
\textbf{(e)} Určte rozdiel hodnôt lokálneho maxima a lokálneho minima funkcie \(\varphi(x)\). \\[1mm]
\textbf{(f)} Určte súčin \(x\)-ových súradníc bodov globálneho maxima a globálneho minima funkcie \(\omega(x)\) na zadanom intervale. \\[1mm]
\end{multicols}
}
\end{center}
\end{minipage}
\\[4mm]
\begin{minipage}{0.79\linewidth}
\begin{center}
\begin{varwidth}{\linewidth}
\begin{enumerate}
\normalsize
\item $f(x)=\frac{x}{7}+\frac{7}{x} \enspace , \enspace x_0=-15$\quad \dotfill\; ???\;\dotfill \quad \text{rastúca}
\item $g(x)=x^2\, e^{-6x+7}$\quad \dotfill\; ???\;\dotfill \quad $\left\langle\frac{1}{3} , 0\right\rangle$
\item $h(t)=\cfrac{\ln t}{5t^3}$\quad \dotfill\; ???\;\dotfill \quad 2
\item $v(x)=7-\sqrt[3]{(9x+6)^2}$\quad \dotfill\; ???\;\dotfill \quad $\text{nie}$
\item $\varphi(x)=\cfrac{x^2+5x+7}{x+2}$\quad \dotfill\; ???\;\dotfill \quad $-6$
\item $\omega(x)=x^3-\frac{19}{2}x^2+28x-6 \enspace , \enspace x\in\left\langle2 , 6\right\rangle$\quad \dotfill\; ???\;\dotfill \quad $20$
\end{enumerate}
\end{varwidth}
\end{center}
\end{minipage}
\begin{minipage}{0.20\linewidth}
\begin{center}
{\Huge\bfseries 4.} \\[2mm]
\includegraphics[height=37mm]{images/braille.png}
{\small Braillovo písmeno}
\end{center}
\end{minipage}
\end{center}
\end{minipage}
%
\end{tabular}
\newpage
\thispagestyle{empty}
\begin{tabular}{c:c}
\begin{minipage}[c][104.5mm][t]{0.5\linewidth}
\begin{center}
\vspace{7mm}
{\huge Monotónnosť a extrémy, skupina \textit{Iota $\iota$} -\romannumeral1}\\[5mm]
\begin{minipage}{0.95\linewidth}
\begin{center}
{\small \begin{multicols}{2} \raggedright
\textbf{(a)} Určte monotónnosť funkcie \(f(x)\) v okolí bodu \(x_0\). \\[1mm]
\textbf{(b)} Určte intervaly, na ktorých je funkcia \(g(x)\) klesajúca. \\[1mm]
\textbf{(c)} Určte počet stacionárnych bodov funkcie \(h(t)\). \par \columnbreak
\textbf{(d)} Nadobúda funkcia \(v(x)\) lokálny extrém? \\[1mm]
\textbf{(e)} Určte rozdiel hodnôt lokálneho maxima a lokálneho minima funkcie \(\varphi(x)\). \\[1mm]
\textbf{(f)} Určte súčin \(x\)-ových súradníc bodov globálneho maxima a globálneho minima funkcie \(\omega(x)\) na zadanom intervale. \\[1mm]
\end{multicols}
}
\end{center}
\end{minipage}
\\[4mm]
\begin{minipage}{0.79\linewidth}
\begin{center}
\begin{varwidth}{\linewidth}
\begin{enumerate}
\normalsize
\item $f(x)=\frac{x}{6}+\frac{6}{x} \enspace , \enspace x_0=4$\quad \dotfill\; ???\;\dotfill \quad \text{rastúca}
\item $g(x)=x^2\, e^{-6x+1}$\quad \dotfill\; ???\;\dotfill \quad $\left(-\infty , 0\right\rangle\enspace , \enspace \left\langle\frac{1}{3} , \infty\right)$
\item $h(t)=\cfrac{\ln t}{9t^5}$\quad \dotfill\; ???\;\dotfill \quad 0
\item $v(x)=-1-\sqrt[7]{(2x+1)^2}$\quad \dotfill\; ???\;\dotfill \quad $\text{nie}$
\item $\varphi(x)=\cfrac{x^2-2x+13}{x+1}$\quad \dotfill\; ???\;\dotfill \quad $16$
\item $\omega(x)=\frac{2}{3}x^3-\frac{5}{2}x^2-3x+1 \enspace , \enspace x\in\left\langle-1 , 5\right\rangle$\quad \dotfill\; ???\;\dotfill \quad $15$
\end{enumerate}
\end{varwidth}
\end{center}
\end{minipage}
\begin{minipage}{0.20\linewidth}
\begin{center}
{\Huge\bfseries 1.} \\[2mm]
\includegraphics[height=37mm]{images/braille.png}
{\small Braillovo písmeno}
\end{center}
\end{minipage}
\end{center}
\end{minipage}
&
\begin{minipage}[c][104.5mm][t]{0.5\linewidth}
\begin{center}
\vspace{7mm}
{\huge Monotónnosť a extrémy, skupina \textit{Iota $\iota$} -\romannumeral2}\\[5mm]
\begin{minipage}{0.95\linewidth}
\begin{center}
{\small \begin{multicols}{2} \raggedright
\textbf{(a)} Určte monotónnosť funkcie \(f(x)\) v okolí bodu \(x_0\). \\[1mm]
\textbf{(b)} Určte intervaly, na ktorých je funkcia \(g(x)\) klesajúca. \\[1mm]
\textbf{(c)} Určte počet stacionárnych bodov funkcie \(h(t)\). \par \columnbreak
\textbf{(d)} Nadobúda funkcia \(v(x)\) lokálny extrém? \\[1mm]
\textbf{(e)} Určte rozdiel hodnôt lokálneho maxima a lokálneho minima funkcie \(\varphi(x)\). \\[1mm]
\textbf{(f)} Určte súčin \(x\)-ových súradníc bodov globálneho maxima a globálneho minima funkcie \(\omega(x)\) na zadanom intervale. \\[1mm]
\end{multicols}
}
\end{center}
\end{minipage}
\\[4mm]
\begin{minipage}{0.79\linewidth}
\begin{center}
\begin{varwidth}{\linewidth}
\begin{enumerate}
\normalsize
\item $f(x)=\frac{x}{2}+\frac{2}{x} \enspace , \enspace x_0=-3$\quad \dotfill\; ???\;\dotfill \quad \text{rastúca}
\item $g(x)=x^2\, e^{-7x-2}$\quad \dotfill\; ???\;\dotfill \quad $\left(-\infty , 0\right\rangle\enspace , \enspace \left\langle\frac{2}{7} , \infty\right)$
\item $h(t)=\cfrac{\ln t}{6t^4}$\quad \dotfill\; ???\;\dotfill \quad 0
\item $v(x)=-6-\sqrt[7]{(-8x-6)^2}$\quad \dotfill\; ???\;\dotfill \quad $\text{nie}$
\item $\varphi(x)=\cfrac{x^2+x+3}{x-2}$\quad \dotfill\; ???\;\dotfill \quad $-12$
\item $\omega(x)=-x^3-x^2+5x-7 \enspace , \enspace x\in\left\langle-4 , 2\right\rangle$\quad \dotfill\; ???\;\dotfill \quad $\frac{20}{3}$
\end{enumerate}
\end{varwidth}
\end{center}
\end{minipage}
\begin{minipage}{0.20\linewidth}
\begin{center}
{\Huge\bfseries 2.} \\[2mm]
\includegraphics[height=37mm]{images/braille.png}
{\small Braillovo písmeno}
\end{center}
\end{minipage}
\end{center}
\end{minipage}
\\ \hdashline
\begin{minipage}[c][104.5mm][t]{0.5\linewidth}
\begin{center}
\vspace{7mm}
{\huge Monotónnosť a extrémy, skupina \textit{Iota $\iota$} -\romannumeral3}\\[5mm]
\begin{minipage}{0.95\linewidth}
\begin{center}
{\small \begin{multicols}{2} \raggedright
\textbf{(a)} Určte monotónnosť funkcie \(f(x)\) v okolí bodu \(x_0\). \\[1mm]
\textbf{(b)} Určte intervaly, na ktorých je funkcia \(g(x)\) klesajúca. \\[1mm]
\textbf{(c)} Určte počet stacionárnych bodov funkcie \(h(t)\). \par \columnbreak
\textbf{(d)} Nadobúda funkcia \(v(x)\) lokálny extrém? \\[1mm]
\textbf{(e)} Určte rozdiel hodnôt lokálneho maxima a lokálneho minima funkcie \(\varphi(x)\). \\[1mm]
\textbf{(f)} Určte súčin \(x\)-ových súradníc bodov globálneho maxima a globálneho minima funkcie \(\omega(x)\) na zadanom intervale. \\[1mm]
\end{multicols}
}
\end{center}
\end{minipage}
\\[4mm]
\begin{minipage}{0.79\linewidth}
\begin{center}
\begin{varwidth}{\linewidth}
\begin{enumerate}
\normalsize
\item $f(x)=\frac{x}{5}+\frac{5}{x} \enspace , \enspace x_0=-11$\quad \dotfill\; ???\;\dotfill \quad \text{rastúca}
\item $g(x)=x^2\, e^{-4x-6}$\quad \dotfill\; ???\;\dotfill \quad $\left\langle\frac{1}{2} , 0\right\rangle$
\item $h(t)=\cfrac{\ln t}{4t^4}$\quad \dotfill\; ???\;\dotfill \quad 1
\item $v(x)=-3-\sqrt[3]{(-4x+3)^2}$\quad \dotfill\; ???\;\dotfill \quad $\text{nie}$
\item $\varphi(x)=\cfrac{x^2-2x-14}{x+3}$\quad \dotfill\; ???\;\dotfill \quad $-4$
\item $\omega(x)=-\frac{2}{3}x^3-\frac{1}{2}x^2+3x+1 \enspace , \enspace x\in\left\langle-4 , 2\right\rangle$\quad \dotfill\; ???\;\dotfill \quad $-\frac{3}{2}$
\end{enumerate}
\end{varwidth}
\end{center}
\end{minipage}
\begin{minipage}{0.20\linewidth}
\begin{center}
{\Huge\bfseries 3.} \\[2mm]
\includegraphics[height=37mm]{images/braille.png}
{\small Braillovo písmeno}
\end{center}
\end{minipage}
\end{center}
\end{minipage}
&
\begin{minipage}[c][104.5mm][t]{0.5\linewidth}
\begin{center}
\vspace{7mm}
{\huge Monotónnosť a extrémy, skupina \textit{Iota $\iota$} -\romannumeral4}\\[5mm]
\begin{minipage}{0.95\linewidth}
\begin{center}
{\small \begin{multicols}{2} \raggedright
\textbf{(a)} Určte monotónnosť funkcie \(f(x)\) v okolí bodu \(x_0\). \\[1mm]
\textbf{(b)} Určte intervaly, na ktorých je funkcia \(g(x)\) klesajúca. \\[1mm]
\textbf{(c)} Určte počet stacionárnych bodov funkcie \(h(t)\). \par \columnbreak
\textbf{(d)} Nadobúda funkcia \(v(x)\) lokálny extrém? \\[1mm]
\textbf{(e)} Určte rozdiel hodnôt lokálneho maxima a lokálneho minima funkcie \(\varphi(x)\). \\[1mm]
\textbf{(f)} Určte súčin \(x\)-ových súradníc bodov globálneho maxima a globálneho minima funkcie \(\omega(x)\) na zadanom intervale. \\[1mm]
\end{multicols}
}
\end{center}
\end{minipage}
\\[4mm]
\begin{minipage}{0.79\linewidth}
\begin{center}
\begin{varwidth}{\linewidth}
\begin{enumerate}
\normalsize
\item $f(x)=\frac{x}{9}+\frac{9}{x} \enspace , \enspace x_0=3$\quad \dotfill\; ???\;\dotfill \quad \text{klesajúca}
\item $g(x)=x^2\, e^{8x-2}$\quad \dotfill\; ???\;\dotfill \quad $\left\langle-\frac{1}{4} , 0\right\rangle$
\item $h(t)=\cfrac{\ln t}{4t^5}$\quad \dotfill\; ???\;\dotfill \quad 1
\item $v(x)=-8-\sqrt[7]{(5x-1)^2}$\quad \dotfill\; ???\;\dotfill \quad $\text{nie}$
\item $\varphi(x)=\cfrac{x^2-x-5}{x-3}$\quad \dotfill\; ???\;\dotfill \quad $-4$
\item $\omega(x)=-2x^3+\frac{37}{2}x^2-35x+5 \enspace , \enspace x\in\left\langle-1 , 6\right\rangle$\quad \dotfill\; ???\;\dotfill \quad $\frac{35}{18}$
\end{enumerate}
\end{varwidth}
\end{center}
\end{minipage}
\begin{minipage}{0.20\linewidth}
\begin{center}
{\Huge\bfseries 4.} \\[2mm]
\includegraphics[height=37mm]{images/braille.png}
{\small Braillovo písmeno}
\end{center}
\end{minipage}
\end{center}
\end{minipage}
%
\end{tabular}
\newpage
\thispagestyle{empty}
\begin{tabular}{c:c}
\begin{minipage}[c][104.5mm][t]{0.5\linewidth}
\begin{center}
\vspace{7mm}
{\huge Monotónnosť a extrémy, skupina \textit{Kappa $\kappa$} -\romannumeral1}\\[5mm]
\begin{minipage}{0.95\linewidth}
\begin{center}
{\small \begin{multicols}{2} \raggedright
\textbf{(a)} Určte monotónnosť funkcie \(f(x)\) v okolí bodu \(x_0\). \\[1mm]
\textbf{(b)} Určte intervaly, na ktorých je funkcia \(g(x)\) klesajúca. \\[1mm]
\textbf{(c)} Určte počet stacionárnych bodov funkcie \(h(t)\). \par \columnbreak
\textbf{(d)} Nadobúda funkcia \(v(x)\) lokálny extrém? \\[1mm]
\textbf{(e)} Určte rozdiel hodnôt lokálneho maxima a lokálneho minima funkcie \(\varphi(x)\). \\[1mm]
\textbf{(f)} Určte súčin \(x\)-ových súradníc bodov globálneho maxima a globálneho minima funkcie \(\omega(x)\) na zadanom intervale. \\[1mm]
\end{multicols}
}
\end{center}
\end{minipage}
\\[4mm]
\begin{minipage}{0.79\linewidth}
\begin{center}
\begin{varwidth}{\linewidth}
\begin{enumerate}
\normalsize
\item $f(x)=\frac{x}{4}+\frac{4}{x} \enspace , \enspace x_0=-5$\quad \dotfill\; ???\;\dotfill \quad \text{rastúca}
\item $g(x)=x^2\, e^{-4x-2}$\quad \dotfill\; ???\;\dotfill \quad $\left\langle\frac{1}{2} , 0\right\rangle$
\item $h(t)=\cfrac{\ln t}{3t^4}$\quad \dotfill\; ???\;\dotfill \quad 2
\item $v(x)=2-\sqrt[5]{(8x-2)^2}$\quad \dotfill\; ???\;\dotfill \quad $\text{nie}$
\item $\varphi(x)=\cfrac{x^2+x-2}{x+3}$\quad \dotfill\; ???\;\dotfill \quad $-6$
\item $\omega(x)=\frac{1}{3}x^3+\frac{11}{2}x^2+28x-3 \enspace , \enspace x\in\left\langle-8 , -2\right\rangle$\quad \dotfill\; ???\;\dotfill \quad $16$
\end{enumerate}
\end{varwidth}
\end{center}
\end{minipage}
\begin{minipage}{0.20\linewidth}
\begin{center}
{\Huge\bfseries 1.} \\[2mm]
\includegraphics[height=37mm]{images/braille.png}
{\small Braillovo písmeno}
\end{center}
\end{minipage}
\end{center}
\end{minipage}
&
\begin{minipage}[c][104.5mm][t]{0.5\linewidth}
\begin{center}
\vspace{7mm}
{\huge Monotónnosť a extrémy, skupina \textit{Kappa $\kappa$} -\romannumeral2}\\[5mm]
\begin{minipage}{0.95\linewidth}
\begin{center}
{\small \begin{multicols}{2} \raggedright
\textbf{(a)} Určte monotónnosť funkcie \(f(x)\) v okolí bodu \(x_0\). \\[1mm]
\textbf{(b)} Určte intervaly, na ktorých je funkcia \(g(x)\) klesajúca. \\[1mm]
\textbf{(c)} Určte počet stacionárnych bodov funkcie \(h(t)\). \par \columnbreak
\textbf{(d)} Nadobúda funkcia \(v(x)\) lokálny extrém? \\[1mm]
\textbf{(e)} Určte rozdiel hodnôt lokálneho maxima a lokálneho minima funkcie \(\varphi(x)\). \\[1mm]
\textbf{(f)} Určte súčin \(x\)-ových súradníc bodov globálneho maxima a globálneho minima funkcie \(\omega(x)\) na zadanom intervale. \\[1mm]
\end{multicols}
}
\end{center}
\end{minipage}
\\[4mm]
\begin{minipage}{0.79\linewidth}
\begin{center}
\begin{varwidth}{\linewidth}
\begin{enumerate}
\normalsize
\item $f(x)=\frac{x}{4}+\frac{4}{x} \enspace , \enspace x_0=-12$\quad \dotfill\; ???\;\dotfill \quad \text{rastúca}
\item $g(x)=x^2\, e^{-2x+3}$\quad \dotfill\; ???\;\dotfill \quad $\left\langle1 , 0\right\rangle$
\item $h(t)=\cfrac{\ln t}{4t^4}$\quad \dotfill\; ???\;\dotfill \quad 1
\item $v(x)=-2-\sqrt[3]{(2x+4)^2}$\quad \dotfill\; ???\;\dotfill \quad $\text{áno}$
\item $\varphi(x)=\cfrac{x^2-2x+1}{x-5}$\quad \dotfill\; ???\;\dotfill \quad $16$
\item $\omega(x)=-\frac{4}{3}x^3+\frac{5}{2}x^2+6x+7 \enspace , \enspace x\in\left\langle-3 , 3\right\rangle$\quad \dotfill\; ???\;\dotfill \quad $-\frac{3}{8}$
\end{enumerate}
\end{varwidth}
\end{center}
\end{minipage}
\begin{minipage}{0.20\linewidth}
\begin{center}
{\Huge\bfseries 2.} \\[2mm]
\includegraphics[height=37mm]{images/braille.png}
{\small Braillovo písmeno}
\end{center}
\end{minipage}
\end{center}
\end{minipage}
\\ \hdashline
\begin{minipage}[c][104.5mm][t]{0.5\linewidth}
\begin{center}
\vspace{7mm}
{\huge Monotónnosť a extrémy, skupina \textit{Kappa $\kappa$} -\romannumeral3}\\[5mm]
\begin{minipage}{0.95\linewidth}
\begin{center}
{\small \begin{multicols}{2} \raggedright
\textbf{(a)} Určte monotónnosť funkcie \(f(x)\) v okolí bodu \(x_0\). \\[1mm]
\textbf{(b)} Určte intervaly, na ktorých je funkcia \(g(x)\) klesajúca. \\[1mm]
\textbf{(c)} Určte počet stacionárnych bodov funkcie \(h(t)\). \par \columnbreak
\textbf{(d)} Nadobúda funkcia \(v(x)\) lokálny extrém? \\[1mm]
\textbf{(e)} Určte rozdiel hodnôt lokálneho maxima a lokálneho minima funkcie \(\varphi(x)\). \\[1mm]
\textbf{(f)} Určte súčin \(x\)-ových súradníc bodov globálneho maxima a globálneho minima funkcie \(\omega(x)\) na zadanom intervale. \\[1mm]
\end{multicols}
}
\end{center}
\end{minipage}
\\[4mm]
\begin{minipage}{0.79\linewidth}
\begin{center}
\begin{varwidth}{\linewidth}
\begin{enumerate}
\normalsize
\item $f(x)=\frac{x}{7}+\frac{7}{x} \enspace , \enspace x_0=4$\quad \dotfill\; ???\;\dotfill \quad \text{rastúca}
\item $g(x)=x^2\, e^{-2x+6}$\quad \dotfill\; ???\;\dotfill \quad $\left(-\infty , 0\right\rangle\enspace , \enspace \left\langle1 , \infty\right)$
\item $h(t)=\cfrac{\ln t}{5t^5}$\quad \dotfill\; ???\;\dotfill \quad 1
\item $v(x)=-1-\sqrt[5]{(3x+1)^2}$\quad \dotfill\; ???\;\dotfill \quad $\text{nie}$
\item $\varphi(x)=\cfrac{x^2-5x-20}{x+3}$\quad \dotfill\; ???\;\dotfill \quad $-8$
\item $\omega(x)=x^3-5x^2+3x+1 \enspace , \enspace x\in\left\langle0 , 5\right\rangle$\quad \dotfill\; ???\;\dotfill \quad $15$
\end{enumerate}
\end{varwidth}
\end{center}
\end{minipage}
\begin{minipage}{0.20\linewidth}
\begin{center}
{\Huge\bfseries 3.} \\[2mm]
\includegraphics[height=37mm]{images/braille.png}
{\small Braillovo písmeno}
\end{center}
\end{minipage}
\end{center}
\end{minipage}
&
\begin{minipage}[c][104.5mm][t]{0.5\linewidth}
\begin{center}
\vspace{7mm}
{\huge Monotónnosť a extrémy, skupina \textit{Kappa $\kappa$} -\romannumeral4}\\[5mm]
\begin{minipage}{0.95\linewidth}
\begin{center}
{\small \begin{multicols}{2} \raggedright
\textbf{(a)} Určte monotónnosť funkcie \(f(x)\) v okolí bodu \(x_0\). \\[1mm]
\textbf{(b)} Určte intervaly, na ktorých je funkcia \(g(x)\) klesajúca. \\[1mm]
\textbf{(c)} Určte počet stacionárnych bodov funkcie \(h(t)\). \par \columnbreak
\textbf{(d)} Nadobúda funkcia \(v(x)\) lokálny extrém? \\[1mm]
\textbf{(e)} Určte rozdiel hodnôt lokálneho maxima a lokálneho minima funkcie \(\varphi(x)\). \\[1mm]
\textbf{(f)} Určte súčin \(x\)-ových súradníc bodov globálneho maxima a globálneho minima funkcie \(\omega(x)\) na zadanom intervale. \\[1mm]
\end{multicols}
}
\end{center}
\end{minipage}
\\[4mm]
\begin{minipage}{0.79\linewidth}
\begin{center}
\begin{varwidth}{\linewidth}
\begin{enumerate}
\normalsize
\item $f(x)=\frac{x}{2}+\frac{2}{x} \enspace , \enspace x_0=-5$\quad \dotfill\; ???\;\dotfill \quad \text{rastúca}
\item $g(x)=x^2\, e^{5x-2}$\quad \dotfill\; ???\;\dotfill \quad $\left(-\infty , 0\right\rangle\enspace , \enspace \left\langle-\frac{2}{5} , \infty\right)$
\item $h(t)=\cfrac{\ln t}{4t^5}$\quad \dotfill\; ???\;\dotfill \quad 1
\item $v(x)=2-\sqrt[7]{(-x-7)^2}$\quad \dotfill\; ???\;\dotfill \quad $\text{nie}$
\item $\varphi(x)=\cfrac{x^2-5x+19}{x+1}$\quad \dotfill\; ???\;\dotfill \quad $-20$
\item $\omega(x)=2x^3-\frac{7}{2}x^2+x-1 \enspace , \enspace x\in\left\langle0 , 3\right\rangle$\quad \dotfill\; ???\;\dotfill \quad $2$
\end{enumerate}
\end{varwidth}
\end{center}
\end{minipage}
\begin{minipage}{0.20\linewidth}
\begin{center}
{\Huge\bfseries 4.} \\[2mm]
\includegraphics[height=37mm]{images/braille.png}
{\small Braillovo písmeno}
\end{center}
\end{minipage}
\end{center}
\end{minipage}
%
\end{tabular}
\newpage
\thispagestyle{empty}
\begin{tabular}{c:c}
\begin{minipage}[c][104.5mm][t]{0.5\linewidth}
\begin{center}
\vspace{7mm}
{\huge Monotónnosť a extrémy, skupina \textit{Lambda $\lambda$} -\romannumeral1}\\[5mm]
\begin{minipage}{0.95\linewidth}
\begin{center}
{\small \begin{multicols}{2} \raggedright
\textbf{(a)} Určte monotónnosť funkcie \(f(x)\) v okolí bodu \(x_0\). \\[1mm]
\textbf{(b)} Určte intervaly, na ktorých je funkcia \(g(x)\) klesajúca. \\[1mm]
\textbf{(c)} Určte počet stacionárnych bodov funkcie \(h(t)\). \par \columnbreak
\textbf{(d)} Nadobúda funkcia \(v(x)\) lokálny extrém? \\[1mm]
\textbf{(e)} Určte rozdiel hodnôt lokálneho maxima a lokálneho minima funkcie \(\varphi(x)\). \\[1mm]
\textbf{(f)} Určte súčin \(x\)-ových súradníc bodov globálneho maxima a globálneho minima funkcie \(\omega(x)\) na zadanom intervale. \\[1mm]
\end{multicols}
}
\end{center}
\end{minipage}
\\[4mm]
\begin{minipage}{0.79\linewidth}
\begin{center}
\begin{varwidth}{\linewidth}
\begin{enumerate}
\normalsize
\item $f(x)=\frac{x}{3}+\frac{3}{x} \enspace , \enspace x_0=-5$\quad \dotfill\; ???\;\dotfill \quad \text{klesajúca}
\item $g(x)=x^2\, e^{4x-9}$\quad \dotfill\; ???\;\dotfill \quad $\left\langle-\frac{1}{2} , 0\right\rangle$
\item $h(t)=\cfrac{\ln t}{t^4}$\quad \dotfill\; ???\;\dotfill \quad 1
\item $v(x)=4-\sqrt[3]{(-3x+3)^2}$\quad \dotfill\; ???\;\dotfill \quad $\text{nie}$
\item $\varphi(x)=\cfrac{x^2-x+10}{x-3}$\quad \dotfill\; ???\;\dotfill \quad $-16$
\item $\omega(x)=\frac{5}{3}x^3-\frac{1}{2}x^2-4x-1 \enspace , \enspace x\in\left\langle-1 , 3\right\rangle$\quad \dotfill\; ???\;\dotfill \quad $3$
\end{enumerate}
\end{varwidth}
\end{center}
\end{minipage}
\begin{minipage}{0.20\linewidth}
\begin{center}
{\Huge\bfseries 1.} \\[2mm]
\includegraphics[height=37mm]{images/braille.png}
{\small Braillovo písmeno}
\end{center}
\end{minipage}
\end{center}
\end{minipage}
&
\begin{minipage}[c][104.5mm][t]{0.5\linewidth}
\begin{center}
\vspace{7mm}
{\huge Monotónnosť a extrémy, skupina \textit{Lambda $\lambda$} -\romannumeral2}\\[5mm]
\begin{minipage}{0.95\linewidth}
\begin{center}
{\small \begin{multicols}{2} \raggedright
\textbf{(a)} Určte monotónnosť funkcie \(f(x)\) v okolí bodu \(x_0\). \\[1mm]
\textbf{(b)} Určte intervaly, na ktorých je funkcia \(g(x)\) klesajúca. \\[1mm]
\textbf{(c)} Určte počet stacionárnych bodov funkcie \(h(t)\). \par \columnbreak
\textbf{(d)} Nadobúda funkcia \(v(x)\) lokálny extrém? \\[1mm]
\textbf{(e)} Určte rozdiel hodnôt lokálneho maxima a lokálneho minima funkcie \(\varphi(x)\). \\[1mm]
\textbf{(f)} Určte súčin \(x\)-ových súradníc bodov globálneho maxima a globálneho minima funkcie \(\omega(x)\) na zadanom intervale. \\[1mm]
\end{multicols}
}
\end{center}
\end{minipage}
\\[4mm]
\begin{minipage}{0.79\linewidth}
\begin{center}
\begin{varwidth}{\linewidth}
\begin{enumerate}
\normalsize
\item $f(x)=\frac{x}{2}+\frac{2}{x} \enspace , \enspace x_0=-6$\quad \dotfill\; ???\;\dotfill \quad \text{rastúca}
\item $g(x)=x^2\, e^{-7x+3}$\quad \dotfill\; ???\;\dotfill \quad $\left(-\infty , 0\right\rangle\enspace , \enspace \left\langle\frac{2}{7} , \infty\right)$
\item $h(t)=\cfrac{\ln t}{5t^3}$\quad \dotfill\; ???\;\dotfill \quad 1
\item $v(x)=7-\sqrt[5]{(-6x-1)^2}$\quad \dotfill\; ???\;\dotfill \quad $\text{nie}$
\item $\varphi(x)=\cfrac{x^2+2x+1}{x-4}$\quad \dotfill\; ???\;\dotfill \quad $-18$
\item $\omega(x)=-\frac{2}{3}x^3-\frac{3}{2}x^2+5x-2 \enspace , \enspace x\in\left\langle-5 , 2\right\rangle$\quad \dotfill\; ???\;\dotfill \quad $\frac{5}{2}$
\end{enumerate}
\end{varwidth}
\end{center}
\end{minipage}
\begin{minipage}{0.20\linewidth}
\begin{center}
{\Huge\bfseries 2.} \\[2mm]
\includegraphics[height=37mm]{images/braille.png}
{\small Braillovo písmeno}
\end{center}
\end{minipage}
\end{center}
\end{minipage}
\\ \hdashline
\begin{minipage}[c][104.5mm][t]{0.5\linewidth}
\begin{center}
\vspace{7mm}
{\huge Monotónnosť a extrémy, skupina \textit{Lambda $\lambda$} -\romannumeral3}\\[5mm]
\begin{minipage}{0.95\linewidth}
\begin{center}
{\small \begin{multicols}{2} \raggedright
\textbf{(a)} Určte monotónnosť funkcie \(f(x)\) v okolí bodu \(x_0\). \\[1mm]
\textbf{(b)} Určte intervaly, na ktorých je funkcia \(g(x)\) klesajúca. \\[1mm]
\textbf{(c)} Určte počet stacionárnych bodov funkcie \(h(t)\). \par \columnbreak
\textbf{(d)} Nadobúda funkcia \(v(x)\) lokálny extrém? \\[1mm]
\textbf{(e)} Určte rozdiel hodnôt lokálneho maxima a lokálneho minima funkcie \(\varphi(x)\). \\[1mm]
\textbf{(f)} Určte súčin \(x\)-ových súradníc bodov globálneho maxima a globálneho minima funkcie \(\omega(x)\) na zadanom intervale. \\[1mm]
\end{multicols}
}
\end{center}
\end{minipage}
\\[4mm]
\begin{minipage}{0.79\linewidth}
\begin{center}
\begin{varwidth}{\linewidth}
\begin{enumerate}
\normalsize
\item $f(x)=\frac{x}{5}+\frac{5}{x} \enspace , \enspace x_0=-7$\quad \dotfill\; ???\;\dotfill \quad \text{rastúca}
\item $g(x)=x^2\, e^{-2x+2}$\quad \dotfill\; ???\;\dotfill \quad $\left\langle1 , 0\right\rangle$
\item $h(t)=\cfrac{\ln t}{8t^4}$\quad \dotfill\; ???\;\dotfill \quad 2
\item $v(x)=2-\sqrt[3]{(4x+1)^2}$\quad \dotfill\; ???\;\dotfill \quad $\text{nie}$
\item $\varphi(x)=\cfrac{x^2-7x+13}{x-4}$\quad \dotfill\; ???\;\dotfill \quad $-2$
\item $\omega(x)=\frac{1}{3}x^3-\frac{5}{2}x^2+4x-3 \enspace , \enspace x\in\left\langle0 , 6\right\rangle$\quad \dotfill\; ???\;\dotfill \quad $16$
\end{enumerate}
\end{varwidth}
\end{center}
\end{minipage}
\begin{minipage}{0.20\linewidth}
\begin{center}
{\Huge\bfseries 3.} \\[2mm]
\includegraphics[height=37mm]{images/braille.png}
{\small Braillovo písmeno}
\end{center}
\end{minipage}
\end{center}
\end{minipage}
&
\begin{minipage}[c][104.5mm][t]{0.5\linewidth}
\begin{center}
\vspace{7mm}
{\huge Monotónnosť a extrémy, skupina \textit{Lambda $\lambda$} -\romannumeral4}\\[5mm]
\begin{minipage}{0.95\linewidth}
\begin{center}
{\small \begin{multicols}{2} \raggedright
\textbf{(a)} Určte monotónnosť funkcie \(f(x)\) v okolí bodu \(x_0\). \\[1mm]
\textbf{(b)} Určte intervaly, na ktorých je funkcia \(g(x)\) klesajúca. \\[1mm]
\textbf{(c)} Určte počet stacionárnych bodov funkcie \(h(t)\). \par \columnbreak
\textbf{(d)} Nadobúda funkcia \(v(x)\) lokálny extrém? \\[1mm]
\textbf{(e)} Určte rozdiel hodnôt lokálneho maxima a lokálneho minima funkcie \(\varphi(x)\). \\[1mm]
\textbf{(f)} Určte súčin \(x\)-ových súradníc bodov globálneho maxima a globálneho minima funkcie \(\omega(x)\) na zadanom intervale. \\[1mm]
\end{multicols}
}
\end{center}
\end{minipage}
\\[4mm]
\begin{minipage}{0.79\linewidth}
\begin{center}
\begin{varwidth}{\linewidth}
\begin{enumerate}
\normalsize
\item $f(x)=\frac{x}{9}+\frac{9}{x} \enspace , \enspace x_0=-11$\quad \dotfill\; ???\;\dotfill \quad \text{rastúca}
\item $g(x)=x^2\, e^{5x-7}$\quad \dotfill\; ???\;\dotfill \quad $\left(-\infty , 0\right\rangle\enspace , \enspace \left\langle-\frac{2}{5} , \infty\right)$
\item $h(t)=\cfrac{\ln t}{2t^3}$\quad \dotfill\; ???\;\dotfill \quad 1
\item $v(x)=5-\sqrt[3]{(-x-4)^2}$\quad \dotfill\; ???\;\dotfill \quad $\text{nie}$
\item $\varphi(x)=\cfrac{x^2-x+16}{x-1}$\quad \dotfill\; ???\;\dotfill \quad $16$
\item $\omega(x)=2x^3-\frac{31}{2}x^2+28x+5 \enspace , \enspace x\in\left\langle1 , 6\right\rangle$\quad \dotfill\; ???\;\dotfill \quad $16$
\end{enumerate}
\end{varwidth}
\end{center}
\end{minipage}
\begin{minipage}{0.20\linewidth}
\begin{center}
{\Huge\bfseries 4.} \\[2mm]
\includegraphics[height=37mm]{images/braille.png}
{\small Braillovo písmeno}
\end{center}
\end{minipage}
\end{center}
\end{minipage}
%
\end{tabular}
\newpage
\thispagestyle{empty}
\begin{tabular}{c:c}
\begin{minipage}[c][104.5mm][t]{0.5\linewidth}
\begin{center}
\vspace{7mm}
{\huge Monotónnosť a extrémy, skupina \textit{Mu $\mu$} -\romannumeral1}\\[5mm]
\begin{minipage}{0.95\linewidth}
\begin{center}
{\small \begin{multicols}{2} \raggedright
\textbf{(a)} Určte monotónnosť funkcie \(f(x)\) v okolí bodu \(x_0\). \\[1mm]
\textbf{(b)} Určte intervaly, na ktorých je funkcia \(g(x)\) klesajúca. \\[1mm]
\textbf{(c)} Určte počet stacionárnych bodov funkcie \(h(t)\). \par \columnbreak
\textbf{(d)} Nadobúda funkcia \(v(x)\) lokálny extrém? \\[1mm]
\textbf{(e)} Určte rozdiel hodnôt lokálneho maxima a lokálneho minima funkcie \(\varphi(x)\). \\[1mm]
\textbf{(f)} Určte súčin \(x\)-ových súradníc bodov globálneho maxima a globálneho minima funkcie \(\omega(x)\) na zadanom intervale. \\[1mm]
\end{multicols}
}
\end{center}
\end{minipage}
\\[4mm]
\begin{minipage}{0.79\linewidth}
\begin{center}
\begin{varwidth}{\linewidth}
\begin{enumerate}
\normalsize
\item $f(x)=\frac{x}{8}+\frac{8}{x} \enspace , \enspace x_0=4$\quad \dotfill\; ???\;\dotfill \quad \text{klesajúca}
\item $g(x)=x^2\, e^{-3x+1}$\quad \dotfill\; ???\;\dotfill \quad $\left(-\infty , 0\right\rangle\enspace , \enspace \left\langle\frac{2}{3} , \infty\right)$
\item $h(t)=\cfrac{\ln t}{3t^4}$\quad \dotfill\; ???\;\dotfill \quad 1
\item $v(x)=7-\sqrt[3]{(2x-3)^2}$\quad \dotfill\; ???\;\dotfill \quad $\text{nie}$
\item $\varphi(x)=\cfrac{x^2-2x+49}{x-2}$\quad \dotfill\; ???\;\dotfill \quad $-30$
\item $\omega(x)=-\frac{1}{3}x^3+6x^2-32x-3 \enspace , \enspace x\in\left\langle1 , 9\right\rangle$\quad \dotfill\; ???\;\dotfill \quad $4$
\end{enumerate}
\end{varwidth}
\end{center}
\end{minipage}
\begin{minipage}{0.20\linewidth}
\begin{center}
{\Huge\bfseries 1.} \\[2mm]
\includegraphics[height=37mm]{images/braille.png}
{\small Braillovo písmeno}
\end{center}
\end{minipage}
\end{center}
\end{minipage}
&
\begin{minipage}[c][104.5mm][t]{0.5\linewidth}
\begin{center}
\vspace{7mm}
{\huge Monotónnosť a extrémy, skupina \textit{Mu $\mu$} -\romannumeral2}\\[5mm]
\begin{minipage}{0.95\linewidth}
\begin{center}
{\small \begin{multicols}{2} \raggedright
\textbf{(a)} Určte monotónnosť funkcie \(f(x)\) v okolí bodu \(x_0\). \\[1mm]
\textbf{(b)} Určte intervaly, na ktorých je funkcia \(g(x)\) klesajúca. \\[1mm]
\textbf{(c)} Určte počet stacionárnych bodov funkcie \(h(t)\). \par \columnbreak
\textbf{(d)} Nadobúda funkcia \(v(x)\) lokálny extrém? \\[1mm]
\textbf{(e)} Určte rozdiel hodnôt lokálneho maxima a lokálneho minima funkcie \(\varphi(x)\). \\[1mm]
\textbf{(f)} Určte súčin \(x\)-ových súradníc bodov globálneho maxima a globálneho minima funkcie \(\omega(x)\) na zadanom intervale. \\[1mm]
\end{multicols}
}
\end{center}
\end{minipage}
\\[4mm]
\begin{minipage}{0.79\linewidth}
\begin{center}
\begin{varwidth}{\linewidth}
\begin{enumerate}
\normalsize
\item $f(x)=\frac{x}{3}+\frac{3}{x} \enspace , \enspace x_0=-7$\quad \dotfill\; ???\;\dotfill \quad \text{rastúca}
\item $g(x)=x^2\, e^{7x+2}$\quad \dotfill\; ???\;\dotfill \quad $\left(-\infty , 0\right\rangle\enspace , \enspace \left\langle-\frac{2}{7} , \infty\right)$
\item $h(t)=\cfrac{\ln t}{3t^3}$\quad \dotfill\; ???\;\dotfill \quad 0
\item $v(x)=-9-\sqrt[3]{(-7x-3)^2}$\quad \dotfill\; ???\;\dotfill \quad $\text{nie}$
\item $\varphi(x)=\cfrac{x^2-2x-6}{x-5}$\quad \dotfill\; ???\;\dotfill \quad $-12$
\item $\omega(x)=-\frac{2}{3}x^3+\frac{11}{2}x^2-5x-5 \enspace , \enspace x\in\left\langle-2 , 6\right\rangle$\quad \dotfill\; ???\;\dotfill \quad $1$
\end{enumerate}
\end{varwidth}
\end{center}
\end{minipage}
\begin{minipage}{0.20\linewidth}
\begin{center}
{\Huge\bfseries 2.} \\[2mm]
\includegraphics[height=37mm]{images/braille.png}
{\small Braillovo písmeno}
\end{center}
\end{minipage}
\end{center}
\end{minipage}
\\ \hdashline
\begin{minipage}[c][104.5mm][t]{0.5\linewidth}
\begin{center}
\vspace{7mm}
{\huge Monotónnosť a extrémy, skupina \textit{Mu $\mu$} -\romannumeral3}\\[5mm]
\begin{minipage}{0.95\linewidth}
\begin{center}
{\small \begin{multicols}{2} \raggedright
\textbf{(a)} Určte monotónnosť funkcie \(f(x)\) v okolí bodu \(x_0\). \\[1mm]
\textbf{(b)} Určte intervaly, na ktorých je funkcia \(g(x)\) klesajúca. \\[1mm]
\textbf{(c)} Určte počet stacionárnych bodov funkcie \(h(t)\). \par \columnbreak
\textbf{(d)} Nadobúda funkcia \(v(x)\) lokálny extrém? \\[1mm]
\textbf{(e)} Určte rozdiel hodnôt lokálneho maxima a lokálneho minima funkcie \(\varphi(x)\). \\[1mm]
\textbf{(f)} Určte súčin \(x\)-ových súradníc bodov globálneho maxima a globálneho minima funkcie \(\omega(x)\) na zadanom intervale. \\[1mm]
\end{multicols}
}
\end{center}
\end{minipage}
\\[4mm]
\begin{minipage}{0.79\linewidth}
\begin{center}
\begin{varwidth}{\linewidth}
\begin{enumerate}
\normalsize
\item $f(x)=\frac{x}{3}+\frac{3}{x} \enspace , \enspace x_0=2$\quad \dotfill\; ???\;\dotfill \quad \text{klesajúca}
\item $g(x)=x^2\, e^{-6x-6}$\quad \dotfill\; ???\;\dotfill \quad $\left\langle\frac{1}{3} , 0\right\rangle$
\item $h(t)=\cfrac{\ln t}{6t^5}$\quad \dotfill\; ???\;\dotfill \quad 1
\item $v(x)=-1-\sqrt[3]{(2x+6)^2}$\quad \dotfill\; ???\;\dotfill \quad $\text{nie}$
\item $\varphi(x)=\cfrac{x^2+7x+16}{x+3}$\quad \dotfill\; ???\;\dotfill \quad $-8$
\item $\omega(x)=\frac{2}{3}x^3-\frac{9}{2}x^2+9x-6 \enspace , \enspace x\in\left\langle1 , 5\right\rangle$\quad \dotfill\; ???\;\dotfill \quad $15$
\end{enumerate}
\end{varwidth}
\end{center}
\end{minipage}
\begin{minipage}{0.20\linewidth}
\begin{center}
{\Huge\bfseries 3.} \\[2mm]
\includegraphics[height=37mm]{images/braille.png}
{\small Braillovo písmeno}
\end{center}
\end{minipage}
\end{center}
\end{minipage}
&
\begin{minipage}[c][104.5mm][t]{0.5\linewidth}
\begin{center}
\vspace{7mm}
{\huge Monotónnosť a extrémy, skupina \textit{Mu $\mu$} -\romannumeral4}\\[5mm]
\begin{minipage}{0.95\linewidth}
\begin{center}
{\small \begin{multicols}{2} \raggedright
\textbf{(a)} Určte monotónnosť funkcie \(f(x)\) v okolí bodu \(x_0\). \\[1mm]
\textbf{(b)} Určte intervaly, na ktorých je funkcia \(g(x)\) klesajúca. \\[1mm]
\textbf{(c)} Určte počet stacionárnych bodov funkcie \(h(t)\). \par \columnbreak
\textbf{(d)} Nadobúda funkcia \(v(x)\) lokálny extrém? \\[1mm]
\textbf{(e)} Určte rozdiel hodnôt lokálneho maxima a lokálneho minima funkcie \(\varphi(x)\). \\[1mm]
\textbf{(f)} Určte súčin \(x\)-ových súradníc bodov globálneho maxima a globálneho minima funkcie \(\omega(x)\) na zadanom intervale. \\[1mm]
\end{multicols}
}
\end{center}
\end{minipage}
\\[4mm]
\begin{minipage}{0.79\linewidth}
\begin{center}
\begin{varwidth}{\linewidth}
\begin{enumerate}
\normalsize
\item $f(x)=\frac{x}{3}+\frac{3}{x} \enspace , \enspace x_0=-6$\quad \dotfill\; ???\;\dotfill \quad \text{rastúca}
\item $g(x)=x^2\, e^{-5x-5}$\quad \dotfill\; ???\;\dotfill \quad $\left\langle\frac{2}{5} , 0\right\rangle$
\item $h(t)=\cfrac{\ln t}{4t^3}$\quad \dotfill\; ???\;\dotfill \quad 2
\item $v(x)=-4-\sqrt[5]{(-6x+2)^2}$\quad \dotfill\; ???\;\dotfill \quad $\text{nie}$
\item $\varphi(x)=\cfrac{x^2-9x-35}{x+3}$\quad \dotfill\; ???\;\dotfill \quad $-2$
\item $\omega(x)=-x^3-\frac{5}{2}x^2+8x+1 \enspace , \enspace x\in\left\langle-5 , 2\right\rangle$\quad \dotfill\; ???\;\dotfill \quad $\frac{40}{9}$
\end{enumerate}
\end{varwidth}
\end{center}
\end{minipage}
\begin{minipage}{0.20\linewidth}
\begin{center}
{\Huge\bfseries 4.} \\[2mm]
\includegraphics[height=37mm]{images/braille.png}
{\small Braillovo písmeno}
\end{center}
\end{minipage}
\end{center}
\end{minipage}
%
\end{tabular}
\newpage
\thispagestyle{empty}
\begin{tabular}{c:c}
\begin{minipage}[c][104.5mm][t]{0.5\linewidth}
\begin{center}
\vspace{7mm}
{\huge Monotónnosť a extrémy, skupina \textit{Nu $\nu$} -\romannumeral1}\\[5mm]
\begin{minipage}{0.95\linewidth}
\begin{center}
{\small \begin{multicols}{2} \raggedright
\textbf{(a)} Určte monotónnosť funkcie \(f(x)\) v okolí bodu \(x_0\). \\[1mm]
\textbf{(b)} Určte intervaly, na ktorých je funkcia \(g(x)\) klesajúca. \\[1mm]
\textbf{(c)} Určte počet stacionárnych bodov funkcie \(h(t)\). \par \columnbreak
\textbf{(d)} Nadobúda funkcia \(v(x)\) lokálny extrém? \\[1mm]
\textbf{(e)} Určte rozdiel hodnôt lokálneho maxima a lokálneho minima funkcie \(\varphi(x)\). \\[1mm]
\textbf{(f)} Určte súčin \(x\)-ových súradníc bodov globálneho maxima a globálneho minima funkcie \(\omega(x)\) na zadanom intervale. \\[1mm]
\end{multicols}
}
\end{center}
\end{minipage}
\\[4mm]
\begin{minipage}{0.79\linewidth}
\begin{center}
\begin{varwidth}{\linewidth}
\begin{enumerate}
\normalsize
\item $f(x)=\frac{x}{5}+\frac{5}{x} \enspace , \enspace x_0=-6$\quad \dotfill\; ???\;\dotfill \quad \text{rastúca}
\item $g(x)=x^2\, e^{-5x-3}$\quad \dotfill\; ???\;\dotfill \quad $\left\langle\frac{2}{5} , 0\right\rangle$
\item $h(t)=\cfrac{\ln t}{t^5}$\quad \dotfill\; ???\;\dotfill \quad 1
\item $v(x)=-2-\sqrt[5]{(8x-6)^2}$\quad \dotfill\; ???\;\dotfill \quad $\text{nie}$
\item $\varphi(x)=\cfrac{x^2+3x+16}{x+3}$\quad \dotfill\; ???\;\dotfill \quad $-16$
\item $\omega(x)=\frac{2}{3}x^3-\frac{13}{2}x^2+15x+2 \enspace , \enspace x\in\left\langle1 , 7\right\rangle$\quad \dotfill\; ???\;\dotfill \quad $25$
\end{enumerate}
\end{varwidth}
\end{center}
\end{minipage}
\begin{minipage}{0.20\linewidth}
\begin{center}
{\Huge\bfseries 1.} \\[2mm]
\includegraphics[height=37mm]{images/braille.png}
{\small Braillovo písmeno}
\end{center}
\end{minipage}
\end{center}
\end{minipage}
&
\begin{minipage}[c][104.5mm][t]{0.5\linewidth}
\begin{center}
\vspace{7mm}
{\huge Monotónnosť a extrémy, skupina \textit{Nu $\nu$} -\romannumeral2}\\[5mm]
\begin{minipage}{0.95\linewidth}
\begin{center}
{\small \begin{multicols}{2} \raggedright
\textbf{(a)} Určte monotónnosť funkcie \(f(x)\) v okolí bodu \(x_0\). \\[1mm]
\textbf{(b)} Určte intervaly, na ktorých je funkcia \(g(x)\) klesajúca. \\[1mm]
\textbf{(c)} Určte počet stacionárnych bodov funkcie \(h(t)\). \par \columnbreak
\textbf{(d)} Nadobúda funkcia \(v(x)\) lokálny extrém? \\[1mm]
\textbf{(e)} Určte rozdiel hodnôt lokálneho maxima a lokálneho minima funkcie \(\varphi(x)\). \\[1mm]
\textbf{(f)} Určte súčin \(x\)-ových súradníc bodov globálneho maxima a globálneho minima funkcie \(\omega(x)\) na zadanom intervale. \\[1mm]
\end{multicols}
}
\end{center}
\end{minipage}
\\[4mm]
\begin{minipage}{0.79\linewidth}
\begin{center}
\begin{varwidth}{\linewidth}
\begin{enumerate}
\normalsize
\item $f(x)=\frac{x}{8}+\frac{8}{x} \enspace , \enspace x_0=2$\quad \dotfill\; ???\;\dotfill \quad \text{klesajúca}
\item $g(x)=x^2\, e^{8x-1}$\quad \dotfill\; ???\;\dotfill \quad $\left\langle-\frac{1}{4} , 0\right\rangle$
\item $h(t)=\cfrac{\ln t}{7t^5}$\quad \dotfill\; ???\;\dotfill \quad 1
\item $v(x)=-5-\sqrt[7]{(-5x+8)^2}$\quad \dotfill\; ???\;\dotfill \quad $\text{nie}$
\item $\varphi(x)=\cfrac{x^2-2x+1}{x-3}$\quad \dotfill\; ???\;\dotfill \quad $8$
\item $\omega(x)=-2x^3+\frac{19}{2}x^2-3x-4 \enspace , \enspace x\in\left\langle-2 , 4\right\rangle$\quad \dotfill\; ???\;\dotfill \quad $-\frac{1}{18}$
\end{enumerate}
\end{varwidth}
\end{center}
\end{minipage}
\begin{minipage}{0.20\linewidth}
\begin{center}
{\Huge\bfseries 2.} \\[2mm]
\includegraphics[height=37mm]{images/braille.png}
{\small Braillovo písmeno}
\end{center}
\end{minipage}
\end{center}
\end{minipage}
\\ \hdashline
\begin{minipage}[c][104.5mm][t]{0.5\linewidth}
\begin{center}
\vspace{7mm}
{\huge Monotónnosť a extrémy, skupina \textit{Nu $\nu$} -\romannumeral3}\\[5mm]
\begin{minipage}{0.95\linewidth}
\begin{center}
{\small \begin{multicols}{2} \raggedright
\textbf{(a)} Určte monotónnosť funkcie \(f(x)\) v okolí bodu \(x_0\). \\[1mm]
\textbf{(b)} Určte intervaly, na ktorých je funkcia \(g(x)\) klesajúca. \\[1mm]
\textbf{(c)} Určte počet stacionárnych bodov funkcie \(h(t)\). \par \columnbreak
\textbf{(d)} Nadobúda funkcia \(v(x)\) lokálny extrém? \\[1mm]
\textbf{(e)} Určte rozdiel hodnôt lokálneho maxima a lokálneho minima funkcie \(\varphi(x)\). \\[1mm]
\textbf{(f)} Určte súčin \(x\)-ových súradníc bodov globálneho maxima a globálneho minima funkcie \(\omega(x)\) na zadanom intervale. \\[1mm]
\end{multicols}
}
\end{center}
\end{minipage}
\\[4mm]
\begin{minipage}{0.79\linewidth}
\begin{center}
\begin{varwidth}{\linewidth}
\begin{enumerate}
\normalsize
\item $f(x)=\frac{x}{2}+\frac{2}{x} \enspace , \enspace x_0=-4$\quad \dotfill\; ???\;\dotfill \quad \text{rastúca}
\item $g(x)=x^2\, e^{2x+5}$\quad \dotfill\; ???\;\dotfill \quad $\left(-\infty , 0\right\rangle\enspace , \enspace \left\langle-1 , \infty\right)$
\item $h(t)=\cfrac{\ln t}{2t^4}$\quad \dotfill\; ???\;\dotfill \quad 0
\item $v(x)=7-\sqrt[7]{(-5x-7)^2}$\quad \dotfill\; ???\;\dotfill \quad $\text{nie}$
\item $\varphi(x)=\cfrac{x^2+2x-44}{x-6}$\quad \dotfill\; ???\;\dotfill \quad $-8$
\item $\omega(x)=3x^3-\frac{37}{2}x^2+4x+5 \enspace , \enspace x\in\left\langle0 , 6\right\rangle$\quad \dotfill\; ???\;\dotfill \quad $20$
\end{enumerate}
\end{varwidth}
\end{center}
\end{minipage}
\begin{minipage}{0.20\linewidth}
\begin{center}
{\Huge\bfseries 3.} \\[2mm]
\includegraphics[height=37mm]{images/braille.png}
{\small Braillovo písmeno}
\end{center}
\end{minipage}
\end{center}
\end{minipage}
&
\begin{minipage}[c][104.5mm][t]{0.5\linewidth}
\begin{center}
\vspace{7mm}
{\huge Monotónnosť a extrémy, skupina \textit{Nu $\nu$} -\romannumeral4}\\[5mm]
\begin{minipage}{0.95\linewidth}
\begin{center}
{\small \begin{multicols}{2} \raggedright
\textbf{(a)} Určte monotónnosť funkcie \(f(x)\) v okolí bodu \(x_0\). \\[1mm]
\textbf{(b)} Určte intervaly, na ktorých je funkcia \(g(x)\) klesajúca. \\[1mm]
\textbf{(c)} Určte počet stacionárnych bodov funkcie \(h(t)\). \par \columnbreak
\textbf{(d)} Nadobúda funkcia \(v(x)\) lokálny extrém? \\[1mm]
\textbf{(e)} Určte rozdiel hodnôt lokálneho maxima a lokálneho minima funkcie \(\varphi(x)\). \\[1mm]
\textbf{(f)} Určte súčin \(x\)-ových súradníc bodov globálneho maxima a globálneho minima funkcie \(\omega(x)\) na zadanom intervale. \\[1mm]
\end{multicols}
}
\end{center}
\end{minipage}
\\[4mm]
\begin{minipage}{0.79\linewidth}
\begin{center}
\begin{varwidth}{\linewidth}
\begin{enumerate}
\normalsize
\item $f(x)=\frac{x}{4}+\frac{4}{x} \enspace , \enspace x_0=1$\quad \dotfill\; ???\;\dotfill \quad \text{rastúca}
\item $g(x)=x^2\, e^{8x+1}$\quad \dotfill\; ???\;\dotfill \quad $\left\langle-\frac{1}{4} , 0\right\rangle$
\item $h(t)=\cfrac{\ln t}{9t^5}$\quad \dotfill\; ???\;\dotfill \quad 0
\item $v(x)=6-\sqrt[5]{(x+1)^2}$\quad \dotfill\; ???\;\dotfill \quad $\text{nie}$
\item $\varphi(x)=\cfrac{x^2+3x-9}{x-2}$\quad \dotfill\; ???\;\dotfill \quad $-4$
\item $\omega(x)=\frac{1}{3}x^3-\frac{1}{2}x^2-2x-1 \enspace , \enspace x\in\left\langle-2 , 4\right\rangle$\quad \dotfill\; ???\;\dotfill \quad $8$
\end{enumerate}
\end{varwidth}
\end{center}
\end{minipage}
\begin{minipage}{0.20\linewidth}
\begin{center}
{\Huge\bfseries 4.} \\[2mm]
\includegraphics[height=37mm]{images/braille.png}
{\small Braillovo písmeno}
\end{center}
\end{minipage}
\end{center}
\end{minipage}
%
\end{tabular}
\newpage
\thispagestyle{empty}
\begin{tabular}{c:c}
\begin{minipage}[c][104.5mm][t]{0.5\linewidth}
\begin{center}
\vspace{7mm}
{\huge Monotónnosť a extrémy, skupina \textit{Xi $\xi$} -\romannumeral1}\\[5mm]
\begin{minipage}{0.95\linewidth}
\begin{center}
{\small \begin{multicols}{2} \raggedright
\textbf{(a)} Určte monotónnosť funkcie \(f(x)\) v okolí bodu \(x_0\). \\[1mm]
\textbf{(b)} Určte intervaly, na ktorých je funkcia \(g(x)\) klesajúca. \\[1mm]
\textbf{(c)} Určte počet stacionárnych bodov funkcie \(h(t)\). \par \columnbreak
\textbf{(d)} Nadobúda funkcia \(v(x)\) lokálny extrém? \\[1mm]
\textbf{(e)} Určte rozdiel hodnôt lokálneho maxima a lokálneho minima funkcie \(\varphi(x)\). \\[1mm]
\textbf{(f)} Určte súčin \(x\)-ových súradníc bodov globálneho maxima a globálneho minima funkcie \(\omega(x)\) na zadanom intervale. \\[1mm]
\end{multicols}
}
\end{center}
\end{minipage}
\\[4mm]
\begin{minipage}{0.79\linewidth}
\begin{center}
\begin{varwidth}{\linewidth}
\begin{enumerate}
\normalsize
\item $f(x)=\frac{x}{9}+\frac{9}{x} \enspace , \enspace x_0=-13$\quad \dotfill\; ???\;\dotfill \quad \text{klesajúca}
\item $g(x)=x^2\, e^{-x-3}$\quad \dotfill\; ???\;\dotfill \quad $\left(-\infty , 0\right\rangle\enspace , \enspace \left\langle2 , \infty\right)$
\item $h(t)=\cfrac{\ln t}{8t^5}$\quad \dotfill\; ???\;\dotfill \quad 1
\item $v(x)=-2-\sqrt[3]{(x-5)^2}$\quad \dotfill\; ???\;\dotfill \quad $\text{nie}$
\item $\varphi(x)=\cfrac{x^2+8x-19}{x-2}$\quad \dotfill\; ???\;\dotfill \quad $4$
\item $\omega(x)=-\frac{2}{3}x^3+\frac{7}{2}x^2+4x+4 \enspace , \enspace x\in\left\langle-3 , 5\right\rangle$\quad \dotfill\; ???\;\dotfill \quad $\frac{3}{2}$
\end{enumerate}
\end{varwidth}
\end{center}
\end{minipage}
\begin{minipage}{0.20\linewidth}
\begin{center}
{\Huge\bfseries 1.} \\[2mm]
\includegraphics[height=37mm]{images/braille.png}
{\small Braillovo písmeno}
\end{center}
\end{minipage}
\end{center}
\end{minipage}
&
\begin{minipage}[c][104.5mm][t]{0.5\linewidth}
\begin{center}
\vspace{7mm}
{\huge Monotónnosť a extrémy, skupina \textit{Xi $\xi$} -\romannumeral2}\\[5mm]
\begin{minipage}{0.95\linewidth}
\begin{center}
{\small \begin{multicols}{2} \raggedright
\textbf{(a)} Určte monotónnosť funkcie \(f(x)\) v okolí bodu \(x_0\). \\[1mm]
\textbf{(b)} Určte intervaly, na ktorých je funkcia \(g(x)\) klesajúca. \\[1mm]
\textbf{(c)} Určte počet stacionárnych bodov funkcie \(h(t)\). \par \columnbreak
\textbf{(d)} Nadobúda funkcia \(v(x)\) lokálny extrém? \\[1mm]
\textbf{(e)} Určte rozdiel hodnôt lokálneho maxima a lokálneho minima funkcie \(\varphi(x)\). \\[1mm]
\textbf{(f)} Určte súčin \(x\)-ových súradníc bodov globálneho maxima a globálneho minima funkcie \(\omega(x)\) na zadanom intervale. \\[1mm]
\end{multicols}
}
\end{center}
\end{minipage}
\\[4mm]
\begin{minipage}{0.79\linewidth}
\begin{center}
\begin{varwidth}{\linewidth}
\begin{enumerate}
\normalsize
\item $f(x)=\frac{x}{4}+\frac{4}{x} \enspace , \enspace x_0=-12$\quad \dotfill\; ???\;\dotfill \quad \text{rastúca}
\item $g(x)=x^2\, e^{4x+8}$\quad \dotfill\; ???\;\dotfill \quad $\left(-\infty , 0\right\rangle\enspace , \enspace \left\langle-\frac{1}{2} , \infty\right)$
\item $h(t)=\cfrac{\ln t}{3t^4}$\quad \dotfill\; ???\;\dotfill \quad 1
\item $v(x)=4-\sqrt[3]{(-6x+1)^2}$\quad \dotfill\; ???\;\dotfill \quad $\text{nie}$
\item $\varphi(x)=\cfrac{x^2-3x+1}{x-3}$\quad \dotfill\; ???\;\dotfill \quad $-4$
\item $\omega(x)=-x^3-x^2+5x-1 \enspace , \enspace x\in\left\langle-4 , 2\right\rangle$\quad \dotfill\; ???\;\dotfill \quad $-\frac{20}{9}$
\end{enumerate}
\end{varwidth}
\end{center}
\end{minipage}
\begin{minipage}{0.20\linewidth}
\begin{center}
{\Huge\bfseries 2.} \\[2mm]
\includegraphics[height=37mm]{images/braille.png}
{\small Braillovo písmeno}
\end{center}
\end{minipage}
\end{center}
\end{minipage}
\\ \hdashline
\begin{minipage}[c][104.5mm][t]{0.5\linewidth}
\begin{center}
\vspace{7mm}
{\huge Monotónnosť a extrémy, skupina \textit{Xi $\xi$} -\romannumeral3}\\[5mm]
\begin{minipage}{0.95\linewidth}
\begin{center}
{\small \begin{multicols}{2} \raggedright
\textbf{(a)} Určte monotónnosť funkcie \(f(x)\) v okolí bodu \(x_0\). \\[1mm]
\textbf{(b)} Určte intervaly, na ktorých je funkcia \(g(x)\) klesajúca. \\[1mm]
\textbf{(c)} Určte počet stacionárnych bodov funkcie \(h(t)\). \par \columnbreak
\textbf{(d)} Nadobúda funkcia \(v(x)\) lokálny extrém? \\[1mm]
\textbf{(e)} Určte rozdiel hodnôt lokálneho maxima a lokálneho minima funkcie \(\varphi(x)\). \\[1mm]
\textbf{(f)} Určte súčin \(x\)-ových súradníc bodov globálneho maxima a globálneho minima funkcie \(\omega(x)\) na zadanom intervale. \\[1mm]
\end{multicols}
}
\end{center}
\end{minipage}
\\[4mm]
\begin{minipage}{0.79\linewidth}
\begin{center}
\begin{varwidth}{\linewidth}
\begin{enumerate}
\normalsize
\item $f(x)=\frac{x}{10}+\frac{10}{x} \enspace , \enspace x_0=8$\quad \dotfill\; ???\;\dotfill \quad \text{klesajúca}
\item $g(x)=x^2\, e^{x+2}$\quad \dotfill\; ???\;\dotfill \quad $\left\langle-2 , 0\right\rangle$
\item $h(t)=\cfrac{\ln t}{4t^3}$\quad \dotfill\; ???\;\dotfill \quad 1
\item $v(x)=4-\sqrt[3]{(2x-3)^2}$\quad \dotfill\; ???\;\dotfill \quad $\text{áno}$
\item $\varphi(x)=\cfrac{x^2+3x-1}{x-2}$\quad \dotfill\; ???\;\dotfill \quad $-14$
\item $\omega(x)=-\frac{1}{3}x^3-\frac{3}{2}x^2+4x-7 \enspace , \enspace x\in\left\langle-7 , 2\right\rangle$\quad \dotfill\; ???\;\dotfill \quad $-12$
\end{enumerate}
\end{varwidth}
\end{center}
\end{minipage}
\begin{minipage}{0.20\linewidth}
\begin{center}
{\Huge\bfseries 3.} \\[2mm]
\includegraphics[height=37mm]{images/braille.png}
{\small Braillovo písmeno}
\end{center}
\end{minipage}
\end{center}
\end{minipage}
&
\begin{minipage}[c][104.5mm][t]{0.5\linewidth}
\begin{center}
\vspace{7mm}
{\huge Monotónnosť a extrémy, skupina \textit{Xi $\xi$} -\romannumeral4}\\[5mm]
\begin{minipage}{0.95\linewidth}
\begin{center}
{\small \begin{multicols}{2} \raggedright
\textbf{(a)} Určte monotónnosť funkcie \(f(x)\) v okolí bodu \(x_0\). \\[1mm]
\textbf{(b)} Určte intervaly, na ktorých je funkcia \(g(x)\) klesajúca. \\[1mm]
\textbf{(c)} Určte počet stacionárnych bodov funkcie \(h(t)\). \par \columnbreak
\textbf{(d)} Nadobúda funkcia \(v(x)\) lokálny extrém? \\[1mm]
\textbf{(e)} Určte rozdiel hodnôt lokálneho maxima a lokálneho minima funkcie \(\varphi(x)\). \\[1mm]
\textbf{(f)} Určte súčin \(x\)-ových súradníc bodov globálneho maxima a globálneho minima funkcie \(\omega(x)\) na zadanom intervale. \\[1mm]
\end{multicols}
}
\end{center}
\end{minipage}
\\[4mm]
\begin{minipage}{0.79\linewidth}
\begin{center}
\begin{varwidth}{\linewidth}
\begin{enumerate}
\normalsize
\item $f(x)=\frac{x}{10}+\frac{10}{x} \enspace , \enspace x_0=8$\quad \dotfill\; ???\;\dotfill \quad \text{klesajúca}
\item $g(x)=x^2\, e^{-5x+9}$\quad \dotfill\; ???\;\dotfill \quad $\left\langle\frac{2}{5} , 0\right\rangle$
\item $h(t)=\cfrac{\ln t}{8t^3}$\quad \dotfill\; ???\;\dotfill \quad 0
\item $v(x)=2-\sqrt[5]{(-6x+1)^2}$\quad \dotfill\; ???\;\dotfill \quad $\text{nie}$
\item $\varphi(x)=\cfrac{x^2-6x+12}{x-4}$\quad \dotfill\; ???\;\dotfill \quad $-10$
\item $\omega(x)=\frac{8}{3}x^3-\frac{25}{2}x^2+3x+3 \enspace , \enspace x\in\left\langle0 , 5\right\rangle$\quad \dotfill\; ???\;\dotfill \quad $12$
\end{enumerate}
\end{varwidth}
\end{center}
\end{minipage}
\begin{minipage}{0.20\linewidth}
\begin{center}
{\Huge\bfseries 4.} \\[2mm]
\includegraphics[height=37mm]{images/braille.png}
{\small Braillovo písmeno}
\end{center}
\end{minipage}
\end{center}
\end{minipage}
%
\end{tabular}
\newpage
\thispagestyle{empty}
\begin{tabular}{c:c}
\begin{minipage}[c][104.5mm][t]{0.5\linewidth}
\begin{center}
\vspace{7mm}
{\huge Monotónnosť a extrémy, skupina \textit{Omicron $\omicron$} -\romannumeral1}\\[5mm]
\begin{minipage}{0.95\linewidth}
\begin{center}
{\small \begin{multicols}{2} \raggedright
\textbf{(a)} Určte monotónnosť funkcie \(f(x)\) v okolí bodu \(x_0\). \\[1mm]
\textbf{(b)} Určte intervaly, na ktorých je funkcia \(g(x)\) klesajúca. \\[1mm]
\textbf{(c)} Určte počet stacionárnych bodov funkcie \(h(t)\). \par \columnbreak
\textbf{(d)} Nadobúda funkcia \(v(x)\) lokálny extrém? \\[1mm]
\textbf{(e)} Určte rozdiel hodnôt lokálneho maxima a lokálneho minima funkcie \(\varphi(x)\). \\[1mm]
\textbf{(f)} Určte súčin \(x\)-ových súradníc bodov globálneho maxima a globálneho minima funkcie \(\omega(x)\) na zadanom intervale. \\[1mm]
\end{multicols}
}
\end{center}
\end{minipage}
\\[4mm]
\begin{minipage}{0.79\linewidth}
\begin{center}
\begin{varwidth}{\linewidth}
\begin{enumerate}
\normalsize
\item $f(x)=\frac{x}{7}+\frac{7}{x} \enspace , \enspace x_0=3$\quad \dotfill\; ???\;\dotfill \quad \text{klesajúca}
\item $g(x)=x^2\, e^{x-1}$\quad \dotfill\; ???\;\dotfill \quad $\left(-\infty , 0\right\rangle\enspace , \enspace \left\langle-2 , \infty\right)$
\item $h(t)=\cfrac{\ln t}{9t^5}$\quad \dotfill\; ???\;\dotfill \quad 2
\item $v(x)=3-\sqrt[3]{(2x+4)^2}$\quad \dotfill\; ???\;\dotfill \quad $\text{nie}$
\item $\varphi(x)=\cfrac{x^2-2x+1}{x+3}$\quad \dotfill\; ???\;\dotfill \quad $16$
\item $\omega(x)=\frac{2}{3}x^3+\frac{13}{2}x^2+18x-7 \enspace , \enspace x\in\left\langle-5 , 0\right\rangle$\quad \dotfill\; ???\;\dotfill \quad $0$
\end{enumerate}
\end{varwidth}
\end{center}
\end{minipage}
\begin{minipage}{0.20\linewidth}
\begin{center}
{\Huge\bfseries 1.} \\[2mm]
\includegraphics[height=37mm]{images/braille.png}
{\small Braillovo písmeno}
\end{center}
\end{minipage}
\end{center}
\end{minipage}
&
\begin{minipage}[c][104.5mm][t]{0.5\linewidth}
\begin{center}
\vspace{7mm}
{\huge Monotónnosť a extrémy, skupina \textit{Omicron $\omicron$} -\romannumeral2}\\[5mm]
\begin{minipage}{0.95\linewidth}
\begin{center}
{\small \begin{multicols}{2} \raggedright
\textbf{(a)} Určte monotónnosť funkcie \(f(x)\) v okolí bodu \(x_0\). \\[1mm]
\textbf{(b)} Určte intervaly, na ktorých je funkcia \(g(x)\) klesajúca. \\[1mm]
\textbf{(c)} Určte počet stacionárnych bodov funkcie \(h(t)\). \par \columnbreak
\textbf{(d)} Nadobúda funkcia \(v(x)\) lokálny extrém? \\[1mm]
\textbf{(e)} Určte rozdiel hodnôt lokálneho maxima a lokálneho minima funkcie \(\varphi(x)\). \\[1mm]
\textbf{(f)} Určte súčin \(x\)-ových súradníc bodov globálneho maxima a globálneho minima funkcie \(\omega(x)\) na zadanom intervale. \\[1mm]
\end{multicols}
}
\end{center}
\end{minipage}
\\[4mm]
\begin{minipage}{0.79\linewidth}
\begin{center}
\begin{varwidth}{\linewidth}
\begin{enumerate}
\normalsize
\item $f(x)=\frac{x}{8}+\frac{8}{x} \enspace , \enspace x_0=4$\quad \dotfill\; ???\;\dotfill \quad \text{klesajúca}
\item $g(x)=x^2\, e^{x+3}$\quad \dotfill\; ???\;\dotfill \quad $\left(-\infty , 0\right\rangle\enspace , \enspace \left\langle-2 , \infty\right)$
\item $h(t)=\cfrac{\ln t}{2t^3}$\quad \dotfill\; ???\;\dotfill \quad 0
\item $v(x)=-3-\sqrt[3]{(-7x+1)^2}$\quad \dotfill\; ???\;\dotfill \quad $\text{nie}$
\item $\varphi(x)=\cfrac{x^2-x-1}{x-2}$\quad \dotfill\; ???\;\dotfill \quad $-4$
\item $\omega(x)=-\frac{1}{3}x^3+4x^2-12x+7 \enspace , \enspace x\in\left\langle-1 , 7\right\rangle$\quad \dotfill\; ???\;\dotfill \quad $-6$
\end{enumerate}
\end{varwidth}
\end{center}
\end{minipage}
\begin{minipage}{0.20\linewidth}
\begin{center}
{\Huge\bfseries 2.} \\[2mm]
\includegraphics[height=37mm]{images/braille.png}
{\small Braillovo písmeno}
\end{center}
\end{minipage}
\end{center}
\end{minipage}
\\ \hdashline
\begin{minipage}[c][104.5mm][t]{0.5\linewidth}
\begin{center}
\vspace{7mm}
{\huge Monotónnosť a extrémy, skupina \textit{Omicron $\omicron$} -\romannumeral3}\\[5mm]
\begin{minipage}{0.95\linewidth}
\begin{center}
{\small \begin{multicols}{2} \raggedright
\textbf{(a)} Určte monotónnosť funkcie \(f(x)\) v okolí bodu \(x_0\). \\[1mm]
\textbf{(b)} Určte intervaly, na ktorých je funkcia \(g(x)\) klesajúca. \\[1mm]
\textbf{(c)} Určte počet stacionárnych bodov funkcie \(h(t)\). \par \columnbreak
\textbf{(d)} Nadobúda funkcia \(v(x)\) lokálny extrém? \\[1mm]
\textbf{(e)} Určte rozdiel hodnôt lokálneho maxima a lokálneho minima funkcie \(\varphi(x)\). \\[1mm]
\textbf{(f)} Určte súčin \(x\)-ových súradníc bodov globálneho maxima a globálneho minima funkcie \(\omega(x)\) na zadanom intervale. \\[1mm]
\end{multicols}
}
\end{center}
\end{minipage}
\\[4mm]
\begin{minipage}{0.79\linewidth}
\begin{center}
\begin{varwidth}{\linewidth}
\begin{enumerate}
\normalsize
\item $f(x)=\frac{x}{5}+\frac{5}{x} \enspace , \enspace x_0=-12$\quad \dotfill\; ???\;\dotfill \quad \text{rastúca}
\item $g(x)=x^2\, e^{x+8}$\quad \dotfill\; ???\;\dotfill \quad $\left\langle-2 , 0\right\rangle$
\item $h(t)=\cfrac{\ln t}{4t^5}$\quad \dotfill\; ???\;\dotfill \quad 1
\item $v(x)=-4-\sqrt[7]{(-3x+1)^2}$\quad \dotfill\; ???\;\dotfill \quad $\text{nie}$
\item $\varphi(x)=\cfrac{x^2+5x+10}{x+3}$\quad \dotfill\; ???\;\dotfill \quad $8$
\item $\omega(x)=-\frac{2}{3}x^3+\frac{11}{2}x^2-5x+3 \enspace , \enspace x\in\left\langle-2 , 6\right\rangle$\quad \dotfill\; ???\;\dotfill \quad $-\frac{1}{2}$
\end{enumerate}
\end{varwidth}
\end{center}
\end{minipage}
\begin{minipage}{0.20\linewidth}
\begin{center}
{\Huge\bfseries 3.} \\[2mm]
\includegraphics[height=37mm]{images/braille.png}
{\small Braillovo písmeno}
\end{center}
\end{minipage}
\end{center}
\end{minipage}
&
\begin{minipage}[c][104.5mm][t]{0.5\linewidth}
\begin{center}
\vspace{7mm}
{\huge Monotónnosť a extrémy, skupina \textit{Omicron $\omicron$} -\romannumeral4}\\[5mm]
\begin{minipage}{0.95\linewidth}
\begin{center}
{\small \begin{multicols}{2} \raggedright
\textbf{(a)} Určte monotónnosť funkcie \(f(x)\) v okolí bodu \(x_0\). \\[1mm]
\textbf{(b)} Určte intervaly, na ktorých je funkcia \(g(x)\) klesajúca. \\[1mm]
\textbf{(c)} Určte počet stacionárnych bodov funkcie \(h(t)\). \par \columnbreak
\textbf{(d)} Nadobúda funkcia \(v(x)\) lokálny extrém? \\[1mm]
\textbf{(e)} Určte rozdiel hodnôt lokálneho maxima a lokálneho minima funkcie \(\varphi(x)\). \\[1mm]
\textbf{(f)} Určte súčin \(x\)-ových súradníc bodov globálneho maxima a globálneho minima funkcie \(\omega(x)\) na zadanom intervale. \\[1mm]
\end{multicols}
}
\end{center}
\end{minipage}
\\[4mm]
\begin{minipage}{0.79\linewidth}
\begin{center}
\begin{varwidth}{\linewidth}
\begin{enumerate}
\normalsize
\item $f(x)=\frac{x}{5}+\frac{5}{x} \enspace , \enspace x_0=4$\quad \dotfill\; ???\;\dotfill \quad \text{klesajúca}
\item $g(x)=x^2\, e^{-3x-2}$\quad \dotfill\; ???\;\dotfill \quad $\left\langle\frac{2}{3} , 0\right\rangle$
\item $h(t)=\cfrac{\ln t}{7t^3}$\quad \dotfill\; ???\;\dotfill \quad 0
\item $v(x)=-8-\sqrt[7]{(6x-3)^2}$\quad \dotfill\; ???\;\dotfill \quad $\text{nie}$
\item $\varphi(x)=\cfrac{x^2-x-8}{x+3}$\quad \dotfill\; ???\;\dotfill \quad $8$
\item $\omega(x)=-\frac{4}{3}x^3+\frac{33}{2}x^2-54x-1 \enspace , \enspace x\in\left\langle0 , 7\right\rangle$\quad \dotfill\; ???\;\dotfill \quad $\frac{27}{4}$
\end{enumerate}
\end{varwidth}
\end{center}
\end{minipage}
\begin{minipage}{0.20\linewidth}
\begin{center}
{\Huge\bfseries 4.} \\[2mm]
\includegraphics[height=37mm]{images/braille.png}
{\small Braillovo písmeno}
\end{center}
\end{minipage}
\end{center}
\end{minipage}
%
\end{tabular}
\newpage
\thispagestyle{empty}
\begin{tabular}{c:c}
\begin{minipage}[c][104.5mm][t]{0.5\linewidth}
\begin{center}
\vspace{7mm}
{\huge Monotónnosť a extrémy, skupina \textit{Pi $\pi$} -\romannumeral1}\\[5mm]
\begin{minipage}{0.95\linewidth}
\begin{center}
{\small \begin{multicols}{2} \raggedright
\textbf{(a)} Určte monotónnosť funkcie \(f(x)\) v okolí bodu \(x_0\). \\[1mm]
\textbf{(b)} Určte intervaly, na ktorých je funkcia \(g(x)\) klesajúca. \\[1mm]
\textbf{(c)} Určte počet stacionárnych bodov funkcie \(h(t)\). \par \columnbreak
\textbf{(d)} Nadobúda funkcia \(v(x)\) lokálny extrém? \\[1mm]
\textbf{(e)} Určte rozdiel hodnôt lokálneho maxima a lokálneho minima funkcie \(\varphi(x)\). \\[1mm]
\textbf{(f)} Určte súčin \(x\)-ových súradníc bodov globálneho maxima a globálneho minima funkcie \(\omega(x)\) na zadanom intervale. \\[1mm]
\end{multicols}
}
\end{center}
\end{minipage}
\\[4mm]
\begin{minipage}{0.79\linewidth}
\begin{center}
\begin{varwidth}{\linewidth}
\begin{enumerate}
\normalsize
\item $f(x)=\frac{x}{9}+\frac{9}{x} \enspace , \enspace x_0=1$\quad \dotfill\; ???\;\dotfill \quad \text{rastúca}
\item $g(x)=x^2\, e^{-x-5}$\quad \dotfill\; ???\;\dotfill \quad $\left(-\infty , 0\right\rangle\enspace , \enspace \left\langle2 , \infty\right)$
\item $h(t)=\cfrac{\ln t}{7t^3}$\quad \dotfill\; ???\;\dotfill \quad 1
\item $v(x)=4-\sqrt[5]{(2x-7)^2}$\quad \dotfill\; ???\;\dotfill \quad $\text{nie}$
\item $\varphi(x)=\cfrac{x^2+4x+4}{x-2}$\quad \dotfill\; ???\;\dotfill \quad $16$
\item $\omega(x)=-2x^3+\frac{37}{2}x^2-35x+6 \enspace , \enspace x\in\left\langle-1 , 6\right\rangle$\quad \dotfill\; ???\;\dotfill \quad $-\frac{7}{6}$
\end{enumerate}
\end{varwidth}
\end{center}
\end{minipage}
\begin{minipage}{0.20\linewidth}
\begin{center}
{\Huge\bfseries 1.} \\[2mm]
\includegraphics[height=37mm]{images/braille.png}
{\small Braillovo písmeno}
\end{center}
\end{minipage}
\end{center}
\end{minipage}
&
\begin{minipage}[c][104.5mm][t]{0.5\linewidth}
\begin{center}
\vspace{7mm}
{\huge Monotónnosť a extrémy, skupina \textit{Pi $\pi$} -\romannumeral2}\\[5mm]
\begin{minipage}{0.95\linewidth}
\begin{center}
{\small \begin{multicols}{2} \raggedright
\textbf{(a)} Určte monotónnosť funkcie \(f(x)\) v okolí bodu \(x_0\). \\[1mm]
\textbf{(b)} Určte intervaly, na ktorých je funkcia \(g(x)\) klesajúca. \\[1mm]
\textbf{(c)} Určte počet stacionárnych bodov funkcie \(h(t)\). \par \columnbreak
\textbf{(d)} Nadobúda funkcia \(v(x)\) lokálny extrém? \\[1mm]
\textbf{(e)} Určte rozdiel hodnôt lokálneho maxima a lokálneho minima funkcie \(\varphi(x)\). \\[1mm]
\textbf{(f)} Určte súčin \(x\)-ových súradníc bodov globálneho maxima a globálneho minima funkcie \(\omega(x)\) na zadanom intervale. \\[1mm]
\end{multicols}
}
\end{center}
\end{minipage}
\\[4mm]
\begin{minipage}{0.79\linewidth}
\begin{center}
\begin{varwidth}{\linewidth}
\begin{enumerate}
\normalsize
\item $f(x)=\frac{x}{8}+\frac{8}{x} \enspace , \enspace x_0=-9$\quad \dotfill\; ???\;\dotfill \quad \text{rastúca}
\item $g(x)=x^2\, e^{7x+5}$\quad \dotfill\; ???\;\dotfill \quad $\left\langle-\frac{2}{7} , 0\right\rangle$
\item $h(t)=\cfrac{\ln t}{7t^3}$\quad \dotfill\; ???\;\dotfill \quad 1
\item $v(x)=3-\sqrt[3]{(4x-3)^2}$\quad \dotfill\; ???\;\dotfill \quad $\text{nie}$
\item $\varphi(x)=\cfrac{x^2-2x-20}{x+4}$\quad \dotfill\; ???\;\dotfill \quad $-6$
\item $\omega(x)=-\frac{1}{3}x^3-\frac{9}{2}x^2-18x+1 \enspace , \enspace x\in\left\langle-9 , -2\right\rangle$\quad \dotfill\; ???\;\dotfill \quad $-30$
\end{enumerate}
\end{varwidth}
\end{center}
\end{minipage}
\begin{minipage}{0.20\linewidth}
\begin{center}
{\Huge\bfseries 2.} \\[2mm]
\includegraphics[height=37mm]{images/braille.png}
{\small Braillovo písmeno}
\end{center}
\end{minipage}
\end{center}
\end{minipage}
\\ \hdashline
\begin{minipage}[c][104.5mm][t]{0.5\linewidth}
\begin{center}
\vspace{7mm}
{\huge Monotónnosť a extrémy, skupina \textit{Pi $\pi$} -\romannumeral3}\\[5mm]
\begin{minipage}{0.95\linewidth}
\begin{center}
{\small \begin{multicols}{2} \raggedright
\textbf{(a)} Určte monotónnosť funkcie \(f(x)\) v okolí bodu \(x_0\). \\[1mm]
\textbf{(b)} Určte intervaly, na ktorých je funkcia \(g(x)\) klesajúca. \\[1mm]
\textbf{(c)} Určte počet stacionárnych bodov funkcie \(h(t)\). \par \columnbreak
\textbf{(d)} Nadobúda funkcia \(v(x)\) lokálny extrém? \\[1mm]
\textbf{(e)} Určte rozdiel hodnôt lokálneho maxima a lokálneho minima funkcie \(\varphi(x)\). \\[1mm]
\textbf{(f)} Určte súčin \(x\)-ových súradníc bodov globálneho maxima a globálneho minima funkcie \(\omega(x)\) na zadanom intervale. \\[1mm]
\end{multicols}
}
\end{center}
\end{minipage}
\\[4mm]
\begin{minipage}{0.79\linewidth}
\begin{center}
\begin{varwidth}{\linewidth}
\begin{enumerate}
\normalsize
\item $f(x)=\frac{x}{7}+\frac{7}{x} \enspace , \enspace x_0=-14$\quad \dotfill\; ???\;\dotfill \quad \text{rastúca}
\item $g(x)=x^2\, e^{5x-2}$\quad \dotfill\; ???\;\dotfill \quad $\left(-\infty , 0\right\rangle\enspace , \enspace \left\langle-\frac{2}{5} , \infty\right)$
\item $h(t)=\cfrac{\ln t}{5t^4}$\quad \dotfill\; ???\;\dotfill \quad 1
\item $v(x)=3-\sqrt[5]{(-2x-5)^2}$\quad \dotfill\; ???\;\dotfill \quad $\text{nie}$
\item $\varphi(x)=\cfrac{x^2-8x+23}{x-1}$\quad \dotfill\; ???\;\dotfill \quad $-16$
\item $\omega(x)=-\frac{1}{3}x^3+3x^2-8x+6 \enspace , \enspace x\in\left\langle-1 , 5\right\rangle$\quad \dotfill\; ???\;\dotfill \quad $10$
\end{enumerate}
\end{varwidth}
\end{center}
\end{minipage}
\begin{minipage}{0.20\linewidth}
\begin{center}
{\Huge\bfseries 3.} \\[2mm]
\includegraphics[height=37mm]{images/braille.png}
{\small Braillovo písmeno}
\end{center}
\end{minipage}
\end{center}
\end{minipage}
&
\begin{minipage}[c][104.5mm][t]{0.5\linewidth}
\begin{center}
\vspace{7mm}
{\huge Monotónnosť a extrémy, skupina \textit{Pi $\pi$} -\romannumeral4}\\[5mm]
\begin{minipage}{0.95\linewidth}
\begin{center}
{\small \begin{multicols}{2} \raggedright
\textbf{(a)} Určte monotónnosť funkcie \(f(x)\) v okolí bodu \(x_0\). \\[1mm]
\textbf{(b)} Určte intervaly, na ktorých je funkcia \(g(x)\) klesajúca. \\[1mm]
\textbf{(c)} Určte počet stacionárnych bodov funkcie \(h(t)\). \par \columnbreak
\textbf{(d)} Nadobúda funkcia \(v(x)\) lokálny extrém? \\[1mm]
\textbf{(e)} Určte rozdiel hodnôt lokálneho maxima a lokálneho minima funkcie \(\varphi(x)\). \\[1mm]
\textbf{(f)} Určte súčin \(x\)-ových súradníc bodov globálneho maxima a globálneho minima funkcie \(\omega(x)\) na zadanom intervale. \\[1mm]
\end{multicols}
}
\end{center}
\end{minipage}
\\[4mm]
\begin{minipage}{0.79\linewidth}
\begin{center}
\begin{varwidth}{\linewidth}
\begin{enumerate}
\normalsize
\item $f(x)=\frac{x}{9}+\frac{9}{x} \enspace , \enspace x_0=6$\quad \dotfill\; ???\;\dotfill \quad \text{klesajúca}
\item $g(x)=x^2\, e^{-6x-4}$\quad \dotfill\; ???\;\dotfill \quad $\left\langle\frac{1}{3} , 0\right\rangle$
\item $h(t)=\cfrac{\ln t}{t^3}$\quad \dotfill\; ???\;\dotfill \quad 1
\item $v(x)=1-\sqrt[5]{(-x-4)^2}$\quad \dotfill\; ???\;\dotfill \quad $\text{nie}$
\item $\varphi(x)=\cfrac{x^2-2x+8}{x+2}$\quad \dotfill\; ???\;\dotfill \quad $-16$
\item $\omega(x)=\frac{5}{3}x^3-\frac{21}{2}x^2+4x+9 \enspace , \enspace x\in\left\langle0 , 6\right\rangle$\quad \dotfill\; ???\;\dotfill \quad $24$
\end{enumerate}
\end{varwidth}
\end{center}
\end{minipage}
\begin{minipage}{0.20\linewidth}
\begin{center}
{\Huge\bfseries 4.} \\[2mm]
\includegraphics[height=37mm]{images/braille.png}
{\small Braillovo písmeno}
\end{center}
\end{minipage}
\end{center}
\end{minipage}
%
\end{tabular}
\newpage
\thispagestyle{empty}
\begin{tabular}{c:c}
\begin{minipage}[c][104.5mm][t]{0.5\linewidth}
\begin{center}
\vspace{7mm}
{\huge Monotónnosť a extrémy, skupina \textit{Rho $\rho$} -\romannumeral1}\\[5mm]
\begin{minipage}{0.95\linewidth}
\begin{center}
{\small \begin{multicols}{2} \raggedright
\textbf{(a)} Určte monotónnosť funkcie \(f(x)\) v okolí bodu \(x_0\). \\[1mm]
\textbf{(b)} Určte intervaly, na ktorých je funkcia \(g(x)\) klesajúca. \\[1mm]
\textbf{(c)} Určte počet stacionárnych bodov funkcie \(h(t)\). \par \columnbreak
\textbf{(d)} Nadobúda funkcia \(v(x)\) lokálny extrém? \\[1mm]
\textbf{(e)} Určte rozdiel hodnôt lokálneho maxima a lokálneho minima funkcie \(\varphi(x)\). \\[1mm]
\textbf{(f)} Určte súčin \(x\)-ových súradníc bodov globálneho maxima a globálneho minima funkcie \(\omega(x)\) na zadanom intervale. \\[1mm]
\end{multicols}
}
\end{center}
\end{minipage}
\\[4mm]
\begin{minipage}{0.79\linewidth}
\begin{center}
\begin{varwidth}{\linewidth}
\begin{enumerate}
\normalsize
\item $f(x)=\frac{x}{9}+\frac{9}{x} \enspace , \enspace x_0=6$\quad \dotfill\; ???\;\dotfill \quad \text{klesajúca}
\item $g(x)=x^2\, e^{-x+4}$\quad \dotfill\; ???\;\dotfill \quad $\left(-\infty , 0\right\rangle\enspace , \enspace \left\langle2 , \infty\right)$
\item $h(t)=\cfrac{\ln t}{2t^4}$\quad \dotfill\; ???\;\dotfill \quad 1
\item $v(x)=5-\sqrt[5]{(3x+3)^2}$\quad \dotfill\; ???\;\dotfill \quad $\text{nie}$
\item $\varphi(x)=\cfrac{x^2-4x-3}{x+2}$\quad \dotfill\; ???\;\dotfill \quad $-14$
\item $\omega(x)=\frac{1}{3}x^3-7x^2+48x+3 \enspace , \enspace x\in\left\langle5 , 10\right\rangle$\quad \dotfill\; ???\;\dotfill \quad $80$
\end{enumerate}
\end{varwidth}
\end{center}
\end{minipage}
\begin{minipage}{0.20\linewidth}
\begin{center}
{\Huge\bfseries 1.} \\[2mm]
\includegraphics[height=37mm]{images/braille.png}
{\small Braillovo písmeno}
\end{center}
\end{minipage}
\end{center}
\end{minipage}
&
\begin{minipage}[c][104.5mm][t]{0.5\linewidth}
\begin{center}
\vspace{7mm}
{\huge Monotónnosť a extrémy, skupina \textit{Rho $\rho$} -\romannumeral2}\\[5mm]
\begin{minipage}{0.95\linewidth}
\begin{center}
{\small \begin{multicols}{2} \raggedright
\textbf{(a)} Určte monotónnosť funkcie \(f(x)\) v okolí bodu \(x_0\). \\[1mm]
\textbf{(b)} Určte intervaly, na ktorých je funkcia \(g(x)\) klesajúca. \\[1mm]
\textbf{(c)} Určte počet stacionárnych bodov funkcie \(h(t)\). \par \columnbreak
\textbf{(d)} Nadobúda funkcia \(v(x)\) lokálny extrém? \\[1mm]
\textbf{(e)} Určte rozdiel hodnôt lokálneho maxima a lokálneho minima funkcie \(\varphi(x)\). \\[1mm]
\textbf{(f)} Určte súčin \(x\)-ových súradníc bodov globálneho maxima a globálneho minima funkcie \(\omega(x)\) na zadanom intervale. \\[1mm]
\end{multicols}
}
\end{center}
\end{minipage}
\\[4mm]
\begin{minipage}{0.79\linewidth}
\begin{center}
\begin{varwidth}{\linewidth}
\begin{enumerate}
\normalsize
\item $f(x)=\frac{x}{2}+\frac{2}{x} \enspace , \enspace x_0=1$\quad \dotfill\; ???\;\dotfill \quad \text{klesajúca}
\item $g(x)=x^2\, e^{7x+5}$\quad \dotfill\; ???\;\dotfill \quad $\left\langle-\frac{2}{7} , 0\right\rangle$
\item $h(t)=\cfrac{\ln t}{t^3}$\quad \dotfill\; ???\;\dotfill \quad 1
\item $v(x)=5-\sqrt[3]{(4x-7)^2}$\quad \dotfill\; ???\;\dotfill \quad $\text{nie}$
\item $\varphi(x)=\cfrac{x^2+2x+1}{x+2}$\quad \dotfill\; ???\;\dotfill \quad $-4$
\item $\omega(x)=\frac{4}{3}x^3-\frac{25}{2}x^2+36x+3 \enspace , \enspace x\in\left\langle2 , 6\right\rangle$\quad \dotfill\; ???\;\dotfill \quad $20$
\end{enumerate}
\end{varwidth}
\end{center}
\end{minipage}
\begin{minipage}{0.20\linewidth}
\begin{center}
{\Huge\bfseries 2.} \\[2mm]
\includegraphics[height=37mm]{images/braille.png}
{\small Braillovo písmeno}
\end{center}
\end{minipage}
\end{center}
\end{minipage}
\\ \hdashline
\begin{minipage}[c][104.5mm][t]{0.5\linewidth}
\begin{center}
\vspace{7mm}
{\huge Monotónnosť a extrémy, skupina \textit{Rho $\rho$} -\romannumeral3}\\[5mm]
\begin{minipage}{0.95\linewidth}
\begin{center}
{\small \begin{multicols}{2} \raggedright
\textbf{(a)} Určte monotónnosť funkcie \(f(x)\) v okolí bodu \(x_0\). \\[1mm]
\textbf{(b)} Určte intervaly, na ktorých je funkcia \(g(x)\) klesajúca. \\[1mm]
\textbf{(c)} Určte počet stacionárnych bodov funkcie \(h(t)\). \par \columnbreak
\textbf{(d)} Nadobúda funkcia \(v(x)\) lokálny extrém? \\[1mm]
\textbf{(e)} Určte rozdiel hodnôt lokálneho maxima a lokálneho minima funkcie \(\varphi(x)\). \\[1mm]
\textbf{(f)} Určte súčin \(x\)-ových súradníc bodov globálneho maxima a globálneho minima funkcie \(\omega(x)\) na zadanom intervale. \\[1mm]
\end{multicols}
}
\end{center}
\end{minipage}
\\[4mm]
\begin{minipage}{0.79\linewidth}
\begin{center}
\begin{varwidth}{\linewidth}
\begin{enumerate}
\normalsize
\item $f(x)=\frac{x}{2}+\frac{2}{x} \enspace , \enspace x_0=-4$\quad \dotfill\; ???\;\dotfill \quad \text{klesajúca}
\item $g(x)=x^2\, e^{x-7}$\quad \dotfill\; ???\;\dotfill \quad $\left\langle-2 , 0\right\rangle$
\item $h(t)=\cfrac{\ln t}{7t^3}$\quad \dotfill\; ???\;\dotfill \quad 1
\item $v(x)=4-\sqrt[3]{(-x+1)^2}$\quad \dotfill\; ???\;\dotfill \quad $\text{nie}$
\item $\varphi(x)=\cfrac{x^2+x-29}{x-5}$\quad \dotfill\; ???\;\dotfill \quad $-2$
\item $\omega(x)=-\frac{2}{3}x^3-\frac{1}{2}x^2+10x+1 \enspace , \enspace x\in\left\langle-5 , 3\right\rangle$\quad \dotfill\; ???\;\dotfill \quad $\frac{25}{2}$
\end{enumerate}
\end{varwidth}
\end{center}
\end{minipage}
\begin{minipage}{0.20\linewidth}
\begin{center}
{\Huge\bfseries 3.} \\[2mm]
\includegraphics[height=37mm]{images/braille.png}
{\small Braillovo písmeno}
\end{center}
\end{minipage}
\end{center}
\end{minipage}
&
\begin{minipage}[c][104.5mm][t]{0.5\linewidth}
\begin{center}
\vspace{7mm}
{\huge Monotónnosť a extrémy, skupina \textit{Rho $\rho$} -\romannumeral4}\\[5mm]
\begin{minipage}{0.95\linewidth}
\begin{center}
{\small \begin{multicols}{2} \raggedright
\textbf{(a)} Určte monotónnosť funkcie \(f(x)\) v okolí bodu \(x_0\). \\[1mm]
\textbf{(b)} Určte intervaly, na ktorých je funkcia \(g(x)\) klesajúca. \\[1mm]
\textbf{(c)} Určte počet stacionárnych bodov funkcie \(h(t)\). \par \columnbreak
\textbf{(d)} Nadobúda funkcia \(v(x)\) lokálny extrém? \\[1mm]
\textbf{(e)} Určte rozdiel hodnôt lokálneho maxima a lokálneho minima funkcie \(\varphi(x)\). \\[1mm]
\textbf{(f)} Určte súčin \(x\)-ových súradníc bodov globálneho maxima a globálneho minima funkcie \(\omega(x)\) na zadanom intervale. \\[1mm]
\end{multicols}
}
\end{center}
\end{minipage}
\\[4mm]
\begin{minipage}{0.79\linewidth}
\begin{center}
\begin{varwidth}{\linewidth}
\begin{enumerate}
\normalsize
\item $f(x)=\frac{x}{6}+\frac{6}{x} \enspace , \enspace x_0=-9$\quad \dotfill\; ???\;\dotfill \quad \text{klesajúca}
\item $g(x)=x^2\, e^{-x+9}$\quad \dotfill\; ???\;\dotfill \quad $\left(-\infty , 0\right\rangle\enspace , \enspace \left\langle2 , \infty\right)$
\item $h(t)=\cfrac{\ln t}{2t^5}$\quad \dotfill\; ???\;\dotfill \quad 1
\item $v(x)=-6-\sqrt[7]{(-3x-2)^2}$\quad \dotfill\; ???\;\dotfill \quad $\text{nie}$
\item $\varphi(x)=\cfrac{x^2-6x-46}{x+5}$\quad \dotfill\; ???\;\dotfill \quad $-12$
\item $\omega(x)=-\frac{2}{3}x^3-\frac{5}{2}x^2+7x+5 \enspace , \enspace x\in\left\langle-6 , 2\right\rangle$\quad \dotfill\; ???\;\dotfill \quad $21$
\end{enumerate}
\end{varwidth}
\end{center}
\end{minipage}
\begin{minipage}{0.20\linewidth}
\begin{center}
{\Huge\bfseries 4.} \\[2mm]
\includegraphics[height=37mm]{images/braille.png}
{\small Braillovo písmeno}
\end{center}
\end{minipage}
\end{center}
\end{minipage}
%
\end{tabular}
\newpage
\thispagestyle{empty}
\begin{tabular}{c:c}
\begin{minipage}[c][104.5mm][t]{0.5\linewidth}
\begin{center}
\vspace{7mm}
{\huge Monotónnosť a extrémy, skupina \textit{Sigma $\sigma$} -\romannumeral1}\\[5mm]
\begin{minipage}{0.95\linewidth}
\begin{center}
{\small \begin{multicols}{2} \raggedright
\textbf{(a)} Určte monotónnosť funkcie \(f(x)\) v okolí bodu \(x_0\). \\[1mm]
\textbf{(b)} Určte intervaly, na ktorých je funkcia \(g(x)\) klesajúca. \\[1mm]
\textbf{(c)} Určte počet stacionárnych bodov funkcie \(h(t)\). \par \columnbreak
\textbf{(d)} Nadobúda funkcia \(v(x)\) lokálny extrém? \\[1mm]
\textbf{(e)} Určte rozdiel hodnôt lokálneho maxima a lokálneho minima funkcie \(\varphi(x)\). \\[1mm]
\textbf{(f)} Určte súčin \(x\)-ových súradníc bodov globálneho maxima a globálneho minima funkcie \(\omega(x)\) na zadanom intervale. \\[1mm]
\end{multicols}
}
\end{center}
\end{minipage}
\\[4mm]
\begin{minipage}{0.79\linewidth}
\begin{center}
\begin{varwidth}{\linewidth}
\begin{enumerate}
\normalsize
\item $f(x)=\frac{x}{2}+\frac{2}{x} \enspace , \enspace x_0=-6$\quad \dotfill\; ???\;\dotfill \quad \text{rastúca}
\item $g(x)=x^2\, e^{-8x-3}$\quad \dotfill\; ???\;\dotfill \quad $\left(-\infty , 0\right\rangle\enspace , \enspace \left\langle\frac{1}{4} , \infty\right)$
\item $h(t)=\cfrac{\ln t}{3t^4}$\quad \dotfill\; ???\;\dotfill \quad 2
\item $v(x)=-8-\sqrt[3]{(-5x-1)^2}$\quad \dotfill\; ???\;\dotfill \quad $\text{nie}$
\item $\varphi(x)=\cfrac{x^2+x+3}{x-2}$\quad \dotfill\; ???\;\dotfill \quad $-14$
\item $\omega(x)=-x^3+2x^2+4x-3 \enspace , \enspace x\in\left\langle-3 , 3\right\rangle$\quad \dotfill\; ???\;\dotfill \quad $2$
\end{enumerate}
\end{varwidth}
\end{center}
\end{minipage}
\begin{minipage}{0.20\linewidth}
\begin{center}
{\Huge\bfseries 1.} \\[2mm]
\includegraphics[height=37mm]{images/braille.png}
{\small Braillovo písmeno}
\end{center}
\end{minipage}
\end{center}
\end{minipage}
&
\begin{minipage}[c][104.5mm][t]{0.5\linewidth}
\begin{center}
\vspace{7mm}
{\huge Monotónnosť a extrémy, skupina \textit{Sigma $\sigma$} -\romannumeral2}\\[5mm]
\begin{minipage}{0.95\linewidth}
\begin{center}
{\small \begin{multicols}{2} \raggedright
\textbf{(a)} Určte monotónnosť funkcie \(f(x)\) v okolí bodu \(x_0\). \\[1mm]
\textbf{(b)} Určte intervaly, na ktorých je funkcia \(g(x)\) klesajúca. \\[1mm]
\textbf{(c)} Určte počet stacionárnych bodov funkcie \(h(t)\). \par \columnbreak
\textbf{(d)} Nadobúda funkcia \(v(x)\) lokálny extrém? \\[1mm]
\textbf{(e)} Určte rozdiel hodnôt lokálneho maxima a lokálneho minima funkcie \(\varphi(x)\). \\[1mm]
\textbf{(f)} Určte súčin \(x\)-ových súradníc bodov globálneho maxima a globálneho minima funkcie \(\omega(x)\) na zadanom intervale. \\[1mm]
\end{multicols}
}
\end{center}
\end{minipage}
\\[4mm]
\begin{minipage}{0.79\linewidth}
\begin{center}
\begin{varwidth}{\linewidth}
\begin{enumerate}
\normalsize
\item $f(x)=\frac{x}{5}+\frac{5}{x} \enspace , \enspace x_0=-6$\quad \dotfill\; ???\;\dotfill \quad \text{klesajúca}
\item $g(x)=x^2\, e^{2x-2}$\quad \dotfill\; ???\;\dotfill \quad $\left(-\infty , 0\right\rangle\enspace , \enspace \left\langle-1 , \infty\right)$
\item $h(t)=\cfrac{\ln t}{4t^3}$\quad \dotfill\; ???\;\dotfill \quad 0
\item $v(x)=3-\sqrt[7]{(4x-2)^2}$\quad \dotfill\; ???\;\dotfill \quad $\text{áno}$
\item $\varphi(x)=\cfrac{x^2+x+36}{x+1}$\quad \dotfill\; ???\;\dotfill \quad $-24$
\item $\omega(x)=-\frac{1}{3}x^3-\frac{11}{2}x^2-28x+2 \enspace , \enspace x\in\left\langle-10 , -3\right\rangle$\quad \dotfill\; ???\;\dotfill \quad $70$
\end{enumerate}
\end{varwidth}
\end{center}
\end{minipage}
\begin{minipage}{0.20\linewidth}
\begin{center}
{\Huge\bfseries 2.} \\[2mm]
\includegraphics[height=37mm]{images/braille.png}
{\small Braillovo písmeno}
\end{center}
\end{minipage}
\end{center}
\end{minipage}
\\ \hdashline
\begin{minipage}[c][104.5mm][t]{0.5\linewidth}
\begin{center}
\vspace{7mm}
{\huge Monotónnosť a extrémy, skupina \textit{Sigma $\sigma$} -\romannumeral3}\\[5mm]
\begin{minipage}{0.95\linewidth}
\begin{center}
{\small \begin{multicols}{2} \raggedright
\textbf{(a)} Určte monotónnosť funkcie \(f(x)\) v okolí bodu \(x_0\). \\[1mm]
\textbf{(b)} Určte intervaly, na ktorých je funkcia \(g(x)\) klesajúca. \\[1mm]
\textbf{(c)} Určte počet stacionárnych bodov funkcie \(h(t)\). \par \columnbreak
\textbf{(d)} Nadobúda funkcia \(v(x)\) lokálny extrém? \\[1mm]
\textbf{(e)} Určte rozdiel hodnôt lokálneho maxima a lokálneho minima funkcie \(\varphi(x)\). \\[1mm]
\textbf{(f)} Určte súčin \(x\)-ových súradníc bodov globálneho maxima a globálneho minima funkcie \(\omega(x)\) na zadanom intervale. \\[1mm]
\end{multicols}
}
\end{center}
\end{minipage}
\\[4mm]
\begin{minipage}{0.79\linewidth}
\begin{center}
\begin{varwidth}{\linewidth}
\begin{enumerate}
\normalsize
\item $f(x)=\frac{x}{9}+\frac{9}{x} \enspace , \enspace x_0=5$\quad \dotfill\; ???\;\dotfill \quad \text{klesajúca}
\item $g(x)=x^2\, e^{x-3}$\quad \dotfill\; ???\;\dotfill \quad $\left(-\infty , 0\right\rangle\enspace , \enspace \left\langle-2 , \infty\right)$
\item $h(t)=\cfrac{\ln t}{2t^5}$\quad \dotfill\; ???\;\dotfill \quad 2
\item $v(x)=4-\sqrt[5]{(-8x+1)^2}$\quad \dotfill\; ???\;\dotfill \quad $\text{nie}$
\item $\varphi(x)=\cfrac{x^2-x+7}{x+1}$\quad \dotfill\; ???\;\dotfill \quad $-10$
\item $\omega(x)=3x^3-\frac{37}{2}x^2+4x-1 \enspace , \enspace x\in\left\langle0 , 6\right\rangle$\quad \dotfill\; ???\;\dotfill \quad $20$
\end{enumerate}
\end{varwidth}
\end{center}
\end{minipage}
\begin{minipage}{0.20\linewidth}
\begin{center}
{\Huge\bfseries 3.} \\[2mm]
\includegraphics[height=37mm]{images/braille.png}
{\small Braillovo písmeno}
\end{center}
\end{minipage}
\end{center}
\end{minipage}
&
\begin{minipage}[c][104.5mm][t]{0.5\linewidth}
\begin{center}
\vspace{7mm}
{\huge Monotónnosť a extrémy, skupina \textit{Sigma $\sigma$} -\romannumeral4}\\[5mm]
\begin{minipage}{0.95\linewidth}
\begin{center}
{\small \begin{multicols}{2} \raggedright
\textbf{(a)} Určte monotónnosť funkcie \(f(x)\) v okolí bodu \(x_0\). \\[1mm]
\textbf{(b)} Určte intervaly, na ktorých je funkcia \(g(x)\) klesajúca. \\[1mm]
\textbf{(c)} Určte počet stacionárnych bodov funkcie \(h(t)\). \par \columnbreak
\textbf{(d)} Nadobúda funkcia \(v(x)\) lokálny extrém? \\[1mm]
\textbf{(e)} Určte rozdiel hodnôt lokálneho maxima a lokálneho minima funkcie \(\varphi(x)\). \\[1mm]
\textbf{(f)} Určte súčin \(x\)-ových súradníc bodov globálneho maxima a globálneho minima funkcie \(\omega(x)\) na zadanom intervale. \\[1mm]
\end{multicols}
}
\end{center}
\end{minipage}
\\[4mm]
\begin{minipage}{0.79\linewidth}
\begin{center}
\begin{varwidth}{\linewidth}
\begin{enumerate}
\normalsize
\item $f(x)=\frac{x}{7}+\frac{7}{x} \enspace , \enspace x_0=6$\quad \dotfill\; ???\;\dotfill \quad \text{klesajúca}
\item $g(x)=x^2\, e^{7x+7}$\quad \dotfill\; ???\;\dotfill \quad $\left(-\infty , 0\right\rangle\enspace , \enspace \left\langle-\frac{2}{7} , \infty\right)$
\item $h(t)=\cfrac{\ln t}{2t^5}$\quad \dotfill\; ???\;\dotfill \quad 1
\item $v(x)=-5-\sqrt[5]{(3x+1)^2}$\quad \dotfill\; ???\;\dotfill \quad $\text{nie}$
\item $\varphi(x)=\cfrac{x^2+2x+6}{x-1}$\quad \dotfill\; ???\;\dotfill \quad $12$
\item $\omega(x)=-\frac{4}{3}x^3+\frac{29}{2}x^2-45x-1 \enspace , \enspace x\in\left\langle0 , 6\right\rangle$\quad \dotfill\; ???\;\dotfill \quad $-\frac{45}{8}$
\end{enumerate}
\end{varwidth}
\end{center}
\end{minipage}
\begin{minipage}{0.20\linewidth}
\begin{center}
{\Huge\bfseries 4.} \\[2mm]
\includegraphics[height=37mm]{images/braille.png}
{\small Braillovo písmeno}
\end{center}
\end{minipage}
\end{center}
\end{minipage}
%
\end{tabular}
\newpage
\thispagestyle{empty}
\begin{tabular}{c:c}
\begin{minipage}[c][104.5mm][t]{0.5\linewidth}
\begin{center}
\vspace{7mm}
{\huge Monotónnosť a extrémy, skupina \textit{Tau $\tau$} -\romannumeral1}\\[5mm]
\begin{minipage}{0.95\linewidth}
\begin{center}
{\small \begin{multicols}{2} \raggedright
\textbf{(a)} Určte monotónnosť funkcie \(f(x)\) v okolí bodu \(x_0\). \\[1mm]
\textbf{(b)} Určte intervaly, na ktorých je funkcia \(g(x)\) klesajúca. \\[1mm]
\textbf{(c)} Určte počet stacionárnych bodov funkcie \(h(t)\). \par \columnbreak
\textbf{(d)} Nadobúda funkcia \(v(x)\) lokálny extrém? \\[1mm]
\textbf{(e)} Určte rozdiel hodnôt lokálneho maxima a lokálneho minima funkcie \(\varphi(x)\). \\[1mm]
\textbf{(f)} Určte súčin \(x\)-ových súradníc bodov globálneho maxima a globálneho minima funkcie \(\omega(x)\) na zadanom intervale. \\[1mm]
\end{multicols}
}
\end{center}
\end{minipage}
\\[4mm]
\begin{minipage}{0.79\linewidth}
\begin{center}
\begin{varwidth}{\linewidth}
\begin{enumerate}
\normalsize
\item $f(x)=\frac{x}{3}+\frac{3}{x} \enspace , \enspace x_0=-4$\quad \dotfill\; ???\;\dotfill \quad \text{rastúca}
\item $g(x)=x^2\, e^{6x-7}$\quad \dotfill\; ???\;\dotfill \quad $\left(-\infty , 0\right\rangle\enspace , \enspace \left\langle-\frac{1}{3} , \infty\right)$
\item $h(t)=\cfrac{\ln t}{t^4}$\quad \dotfill\; ???\;\dotfill \quad 1
\item $v(x)=6-\sqrt[3]{(4x+1)^2}$\quad \dotfill\; ???\;\dotfill \quad $\text{nie}$
\item $\varphi(x)=\cfrac{x^2+5x+13}{x+1}$\quad \dotfill\; ???\;\dotfill \quad $-12$
\item $\omega(x)=\frac{4}{3}x^3+\frac{3}{2}x^2-7x+3 \enspace , \enspace x\in\left\langle-2 , 3\right\rangle$\quad \dotfill\; ???\;\dotfill \quad $1$
\end{enumerate}
\end{varwidth}
\end{center}
\end{minipage}
\begin{minipage}{0.20\linewidth}
\begin{center}
{\Huge\bfseries 1.} \\[2mm]
\includegraphics[height=37mm]{images/braille.png}
{\small Braillovo písmeno}
\end{center}
\end{minipage}
\end{center}
\end{minipage}
&
\begin{minipage}[c][104.5mm][t]{0.5\linewidth}
\begin{center}
\vspace{7mm}
{\huge Monotónnosť a extrémy, skupina \textit{Tau $\tau$} -\romannumeral2}\\[5mm]
\begin{minipage}{0.95\linewidth}
\begin{center}
{\small \begin{multicols}{2} \raggedright
\textbf{(a)} Určte monotónnosť funkcie \(f(x)\) v okolí bodu \(x_0\). \\[1mm]
\textbf{(b)} Určte intervaly, na ktorých je funkcia \(g(x)\) klesajúca. \\[1mm]
\textbf{(c)} Určte počet stacionárnych bodov funkcie \(h(t)\). \par \columnbreak
\textbf{(d)} Nadobúda funkcia \(v(x)\) lokálny extrém? \\[1mm]
\textbf{(e)} Určte rozdiel hodnôt lokálneho maxima a lokálneho minima funkcie \(\varphi(x)\). \\[1mm]
\textbf{(f)} Určte súčin \(x\)-ových súradníc bodov globálneho maxima a globálneho minima funkcie \(\omega(x)\) na zadanom intervale. \\[1mm]
\end{multicols}
}
\end{center}
\end{minipage}
\\[4mm]
\begin{minipage}{0.79\linewidth}
\begin{center}
\begin{varwidth}{\linewidth}
\begin{enumerate}
\normalsize
\item $f(x)=\frac{x}{3}+\frac{3}{x} \enspace , \enspace x_0=1$\quad \dotfill\; ???\;\dotfill \quad \text{rastúca}
\item $g(x)=x^2\, e^{6x-1}$\quad \dotfill\; ???\;\dotfill \quad $\left\langle-\frac{1}{3} , 0\right\rangle$
\item $h(t)=\cfrac{\ln t}{t^3}$\quad \dotfill\; ???\;\dotfill \quad 1
\item $v(x)=1-\sqrt[5]{(-3x+1)^2}$\quad \dotfill\; ???\;\dotfill \quad $\text{nie}$
\item $\varphi(x)=\cfrac{x^2+3x+16}{x+3}$\quad \dotfill\; ???\;\dotfill \quad $16$
\item $\omega(x)=-\frac{2}{3}x^3+\frac{15}{2}x^2-25x-3 \enspace , \enspace x\in\left\langle0 , 6\right\rangle$\quad \dotfill\; ???\;\dotfill \quad $0$
\end{enumerate}
\end{varwidth}
\end{center}
\end{minipage}
\begin{minipage}{0.20\linewidth}
\begin{center}
{\Huge\bfseries 2.} \\[2mm]
\includegraphics[height=37mm]{images/braille.png}
{\small Braillovo písmeno}
\end{center}
\end{minipage}
\end{center}
\end{minipage}
\\ \hdashline
\begin{minipage}[c][104.5mm][t]{0.5\linewidth}
\begin{center}
\vspace{7mm}
{\huge Monotónnosť a extrémy, skupina \textit{Tau $\tau$} -\romannumeral3}\\[5mm]
\begin{minipage}{0.95\linewidth}
\begin{center}
{\small \begin{multicols}{2} \raggedright
\textbf{(a)} Určte monotónnosť funkcie \(f(x)\) v okolí bodu \(x_0\). \\[1mm]
\textbf{(b)} Určte intervaly, na ktorých je funkcia \(g(x)\) klesajúca. \\[1mm]
\textbf{(c)} Určte počet stacionárnych bodov funkcie \(h(t)\). \par \columnbreak
\textbf{(d)} Nadobúda funkcia \(v(x)\) lokálny extrém? \\[1mm]
\textbf{(e)} Určte rozdiel hodnôt lokálneho maxima a lokálneho minima funkcie \(\varphi(x)\). \\[1mm]
\textbf{(f)} Určte súčin \(x\)-ových súradníc bodov globálneho maxima a globálneho minima funkcie \(\omega(x)\) na zadanom intervale. \\[1mm]
\end{multicols}
}
\end{center}
\end{minipage}
\\[4mm]
\begin{minipage}{0.79\linewidth}
\begin{center}
\begin{varwidth}{\linewidth}
\begin{enumerate}
\normalsize
\item $f(x)=\frac{x}{6}+\frac{6}{x} \enspace , \enspace x_0=1$\quad \dotfill\; ???\;\dotfill \quad \text{klesajúca}
\item $g(x)=x^2\, e^{6x+1}$\quad \dotfill\; ???\;\dotfill \quad $\left(-\infty , 0\right\rangle\enspace , \enspace \left\langle-\frac{1}{3} , \infty\right)$
\item $h(t)=\cfrac{\ln t}{t^4}$\quad \dotfill\; ???\;\dotfill \quad 1
\item $v(x)=-3-\sqrt[5]{(7x-1)^2}$\quad \dotfill\; ???\;\dotfill \quad $\text{nie}$
\item $\varphi(x)=\cfrac{x^2-x+7}{x-2}$\quad \dotfill\; ???\;\dotfill \quad $-12$
\item $\omega(x)=\frac{5}{3}x^3-\frac{11}{2}x^2+2x+2 \enspace , \enspace x\in\left\langle0 , 4\right\rangle$\quad \dotfill\; ???\;\dotfill \quad $4$
\end{enumerate}
\end{varwidth}
\end{center}
\end{minipage}
\begin{minipage}{0.20\linewidth}
\begin{center}
{\Huge\bfseries 3.} \\[2mm]
\includegraphics[height=37mm]{images/braille.png}
{\small Braillovo písmeno}
\end{center}
\end{minipage}
\end{center}
\end{minipage}
&
\begin{minipage}[c][104.5mm][t]{0.5\linewidth}
\begin{center}
\vspace{7mm}
{\huge Monotónnosť a extrémy, skupina \textit{Tau $\tau$} -\romannumeral4}\\[5mm]
\begin{minipage}{0.95\linewidth}
\begin{center}
{\small \begin{multicols}{2} \raggedright
\textbf{(a)} Určte monotónnosť funkcie \(f(x)\) v okolí bodu \(x_0\). \\[1mm]
\textbf{(b)} Určte intervaly, na ktorých je funkcia \(g(x)\) klesajúca. \\[1mm]
\textbf{(c)} Určte počet stacionárnych bodov funkcie \(h(t)\). \par \columnbreak
\textbf{(d)} Nadobúda funkcia \(v(x)\) lokálny extrém? \\[1mm]
\textbf{(e)} Určte rozdiel hodnôt lokálneho maxima a lokálneho minima funkcie \(\varphi(x)\). \\[1mm]
\textbf{(f)} Určte súčin \(x\)-ových súradníc bodov globálneho maxima a globálneho minima funkcie \(\omega(x)\) na zadanom intervale. \\[1mm]
\end{multicols}
}
\end{center}
\end{minipage}
\\[4mm]
\begin{minipage}{0.79\linewidth}
\begin{center}
\begin{varwidth}{\linewidth}
\begin{enumerate}
\normalsize
\item $f(x)=\frac{x}{5}+\frac{5}{x} \enspace , \enspace x_0=-10$\quad \dotfill\; ???\;\dotfill \quad \text{rastúca}
\item $g(x)=x^2\, e^{-x+1}$\quad \dotfill\; ???\;\dotfill \quad $\left(-\infty , 0\right\rangle\enspace , \enspace \left\langle2 , \infty\right)$
\item $h(t)=\cfrac{\ln t}{3t^4}$\quad \dotfill\; ???\;\dotfill \quad 1
\item $v(x)=-1-\sqrt[5]{(2x+7)^2}$\quad \dotfill\; ???\;\dotfill \quad $\text{nie}$
\item $\varphi(x)=\cfrac{x^2-3x+3}{x-2}$\quad \dotfill\; ???\;\dotfill \quad $-2$
\item $\omega(x)=\frac{4}{3}x^3-\frac{13}{2}x^2+3x+1 \enspace , \enspace x\in\left\langle0 , 5\right\rangle$\quad \dotfill\; ???\;\dotfill \quad $12$
\end{enumerate}
\end{varwidth}
\end{center}
\end{minipage}
\begin{minipage}{0.20\linewidth}
\begin{center}
{\Huge\bfseries 4.} \\[2mm]
\includegraphics[height=37mm]{images/braille.png}
{\small Braillovo písmeno}
\end{center}
\end{minipage}
\end{center}
\end{minipage}
%
\end{tabular}
\newpage
\thispagestyle{empty}
\begin{tabular}{c:c}
\begin{minipage}[c][104.5mm][t]{0.5\linewidth}
\begin{center}
\vspace{7mm}
{\huge Monotónnosť a extrémy, skupina \textit{Upsilon $\upsilon$} -\romannumeral1}\\[5mm]
\begin{minipage}{0.95\linewidth}
\begin{center}
{\small \begin{multicols}{2} \raggedright
\textbf{(a)} Určte monotónnosť funkcie \(f(x)\) v okolí bodu \(x_0\). \\[1mm]
\textbf{(b)} Určte intervaly, na ktorých je funkcia \(g(x)\) klesajúca. \\[1mm]
\textbf{(c)} Určte počet stacionárnych bodov funkcie \(h(t)\). \par \columnbreak
\textbf{(d)} Nadobúda funkcia \(v(x)\) lokálny extrém? \\[1mm]
\textbf{(e)} Určte rozdiel hodnôt lokálneho maxima a lokálneho minima funkcie \(\varphi(x)\). \\[1mm]
\textbf{(f)} Určte súčin \(x\)-ových súradníc bodov globálneho maxima a globálneho minima funkcie \(\omega(x)\) na zadanom intervale. \\[1mm]
\end{multicols}
}
\end{center}
\end{minipage}
\\[4mm]
\begin{minipage}{0.79\linewidth}
\begin{center}
\begin{varwidth}{\linewidth}
\begin{enumerate}
\normalsize
\item $f(x)=\frac{x}{4}+\frac{4}{x} \enspace , \enspace x_0=3$\quad \dotfill\; ???\;\dotfill \quad \text{rastúca}
\item $g(x)=x^2\, e^{4x-2}$\quad \dotfill\; ???\;\dotfill \quad $\left\langle-\frac{1}{2} , 0\right\rangle$
\item $h(t)=\cfrac{\ln t}{3t^3}$\quad \dotfill\; ???\;\dotfill \quad 1
\item $v(x)=-6-\sqrt[3]{(2x-1)^2}$\quad \dotfill\; ???\;\dotfill \quad $\text{áno}$
\item $\varphi(x)=\cfrac{x^2+x-16}{x-4}$\quad \dotfill\; ???\;\dotfill \quad $-6$
\item $\omega(x)=\frac{4}{3}x^3-\frac{17}{2}x^2+4x+8 \enspace , \enspace x\in\left\langle0 , 6\right\rangle$\quad \dotfill\; ???\;\dotfill \quad $24$
\end{enumerate}
\end{varwidth}
\end{center}
\end{minipage}
\begin{minipage}{0.20\linewidth}
\begin{center}
{\Huge\bfseries 1.} \\[2mm]
\includegraphics[height=37mm]{images/braille.png}
{\small Braillovo písmeno}
\end{center}
\end{minipage}
\end{center}
\end{minipage}
&
\begin{minipage}[c][104.5mm][t]{0.5\linewidth}
\begin{center}
\vspace{7mm}
{\huge Monotónnosť a extrémy, skupina \textit{Upsilon $\upsilon$} -\romannumeral2}\\[5mm]
\begin{minipage}{0.95\linewidth}
\begin{center}
{\small \begin{multicols}{2} \raggedright
\textbf{(a)} Určte monotónnosť funkcie \(f(x)\) v okolí bodu \(x_0\). \\[1mm]
\textbf{(b)} Určte intervaly, na ktorých je funkcia \(g(x)\) klesajúca. \\[1mm]
\textbf{(c)} Určte počet stacionárnych bodov funkcie \(h(t)\). \par \columnbreak
\textbf{(d)} Nadobúda funkcia \(v(x)\) lokálny extrém? \\[1mm]
\textbf{(e)} Určte rozdiel hodnôt lokálneho maxima a lokálneho minima funkcie \(\varphi(x)\). \\[1mm]
\textbf{(f)} Určte súčin \(x\)-ových súradníc bodov globálneho maxima a globálneho minima funkcie \(\omega(x)\) na zadanom intervale. \\[1mm]
\end{multicols}
}
\end{center}
\end{minipage}
\\[4mm]
\begin{minipage}{0.79\linewidth}
\begin{center}
\begin{varwidth}{\linewidth}
\begin{enumerate}
\normalsize
\item $f(x)=\frac{x}{9}+\frac{9}{x} \enspace , \enspace x_0=6$\quad \dotfill\; ???\;\dotfill \quad \text{klesajúca}
\item $g(x)=x^2\, e^{-6x+1}$\quad \dotfill\; ???\;\dotfill \quad $\left\langle\frac{1}{3} , 0\right\rangle$
\item $h(t)=\cfrac{\ln t}{t^4}$\quad \dotfill\; ???\;\dotfill \quad 2
\item $v(x)=-5-\sqrt[3]{(4x-4)^2}$\quad \dotfill\; ???\;\dotfill \quad $\text{nie}$
\item $\varphi(x)=\cfrac{x^2+9x+24}{x+4}$\quad \dotfill\; ???\;\dotfill \quad $8$
\item $\omega(x)=-\frac{1}{3}x^3+\frac{11}{2}x^2-28x+4 \enspace , \enspace x\in\left\langle1 , 8\right\rangle$\quad \dotfill\; ???\;\dotfill \quad $-20$
\end{enumerate}
\end{varwidth}
\end{center}
\end{minipage}
\begin{minipage}{0.20\linewidth}
\begin{center}
{\Huge\bfseries 2.} \\[2mm]
\includegraphics[height=37mm]{images/braille.png}
{\small Braillovo písmeno}
\end{center}
\end{minipage}
\end{center}
\end{minipage}
\\ \hdashline
\begin{minipage}[c][104.5mm][t]{0.5\linewidth}
\begin{center}
\vspace{7mm}
{\huge Monotónnosť a extrémy, skupina \textit{Upsilon $\upsilon$} -\romannumeral3}\\[5mm]
\begin{minipage}{0.95\linewidth}
\begin{center}
{\small \begin{multicols}{2} \raggedright
\textbf{(a)} Určte monotónnosť funkcie \(f(x)\) v okolí bodu \(x_0\). \\[1mm]
\textbf{(b)} Určte intervaly, na ktorých je funkcia \(g(x)\) klesajúca. \\[1mm]
\textbf{(c)} Určte počet stacionárnych bodov funkcie \(h(t)\). \par \columnbreak
\textbf{(d)} Nadobúda funkcia \(v(x)\) lokálny extrém? \\[1mm]
\textbf{(e)} Určte rozdiel hodnôt lokálneho maxima a lokálneho minima funkcie \(\varphi(x)\). \\[1mm]
\textbf{(f)} Určte súčin \(x\)-ových súradníc bodov globálneho maxima a globálneho minima funkcie \(\omega(x)\) na zadanom intervale. \\[1mm]
\end{multicols}
}
\end{center}
\end{minipage}
\\[4mm]
\begin{minipage}{0.79\linewidth}
\begin{center}
\begin{varwidth}{\linewidth}
\begin{enumerate}
\normalsize
\item $f(x)=\frac{x}{5}+\frac{5}{x} \enspace , \enspace x_0=-7$\quad \dotfill\; ???\;\dotfill \quad \text{rastúca}
\item $g(x)=x^2\, e^{5x+8}$\quad \dotfill\; ???\;\dotfill \quad $\left\langle-\frac{2}{5} , 0\right\rangle$
\item $h(t)=\cfrac{\ln t}{5t^3}$\quad \dotfill\; ???\;\dotfill \quad 0
\item $v(x)=6-\sqrt[7]{(6x+3)^2}$\quad \dotfill\; ???\;\dotfill \quad $\text{nie}$
\item $\varphi(x)=\cfrac{x^2+x+43}{x-2}$\quad \dotfill\; ???\;\dotfill \quad $28$
\item $\omega(x)=-\frac{1}{3}x^3-\frac{11}{2}x^2-28x-2 \enspace , \enspace x\in\left\langle-10 , -3\right\rangle$\quad \dotfill\; ???\;\dotfill \quad $28$
\end{enumerate}
\end{varwidth}
\end{center}
\end{minipage}
\begin{minipage}{0.20\linewidth}
\begin{center}
{\Huge\bfseries 3.} \\[2mm]
\includegraphics[height=37mm]{images/braille.png}
{\small Braillovo písmeno}
\end{center}
\end{minipage}
\end{center}
\end{minipage}
&
\begin{minipage}[c][104.5mm][t]{0.5\linewidth}
\begin{center}
\vspace{7mm}
{\huge Monotónnosť a extrémy, skupina \textit{Upsilon $\upsilon$} -\romannumeral4}\\[5mm]
\begin{minipage}{0.95\linewidth}
\begin{center}
{\small \begin{multicols}{2} \raggedright
\textbf{(a)} Určte monotónnosť funkcie \(f(x)\) v okolí bodu \(x_0\). \\[1mm]
\textbf{(b)} Určte intervaly, na ktorých je funkcia \(g(x)\) klesajúca. \\[1mm]
\textbf{(c)} Určte počet stacionárnych bodov funkcie \(h(t)\). \par \columnbreak
\textbf{(d)} Nadobúda funkcia \(v(x)\) lokálny extrém? \\[1mm]
\textbf{(e)} Určte rozdiel hodnôt lokálneho maxima a lokálneho minima funkcie \(\varphi(x)\). \\[1mm]
\textbf{(f)} Určte súčin \(x\)-ových súradníc bodov globálneho maxima a globálneho minima funkcie \(\omega(x)\) na zadanom intervale. \\[1mm]
\end{multicols}
}
\end{center}
\end{minipage}
\\[4mm]
\begin{minipage}{0.79\linewidth}
\begin{center}
\begin{varwidth}{\linewidth}
\begin{enumerate}
\normalsize
\item $f(x)=\frac{x}{5}+\frac{5}{x} \enspace , \enspace x_0=4$\quad \dotfill\; ???\;\dotfill \quad \text{klesajúca}
\item $g(x)=x^2\, e^{-x+2}$\quad \dotfill\; ???\;\dotfill \quad $\left\langle2 , 0\right\rangle$
\item $h(t)=\cfrac{\ln t}{4t^3}$\quad \dotfill\; ???\;\dotfill \quad 0
\item $v(x)=4-\sqrt[5]{(-7x-4)^2}$\quad \dotfill\; ???\;\dotfill \quad $\text{nie}$
\item $\varphi(x)=\cfrac{x^2+8x+13}{x+6}$\quad \dotfill\; ???\;\dotfill \quad $-2$
\item $\omega(x)=\frac{2}{3}x^3+\frac{1}{2}x^2-3x+2 \enspace , \enspace x\in\left\langle-2 , 3\right\rangle$\quad \dotfill\; ???\;\dotfill \quad $2$
\end{enumerate}
\end{varwidth}
\end{center}
\end{minipage}
\begin{minipage}{0.20\linewidth}
\begin{center}
{\Huge\bfseries 4.} \\[2mm]
\includegraphics[height=37mm]{images/braille.png}
{\small Braillovo písmeno}
\end{center}
\end{minipage}
\end{center}
\end{minipage}
%
\end{tabular}
\newpage
\thispagestyle{empty}
\begin{tabular}{c:c}
\begin{minipage}[c][104.5mm][t]{0.5\linewidth}
\begin{center}
\vspace{7mm}
{\huge Monotónnosť a extrémy, skupina \textit{Phi $\phi$} -\romannumeral1}\\[5mm]
\begin{minipage}{0.95\linewidth}
\begin{center}
{\small \begin{multicols}{2} \raggedright
\textbf{(a)} Určte monotónnosť funkcie \(f(x)\) v okolí bodu \(x_0\). \\[1mm]
\textbf{(b)} Určte intervaly, na ktorých je funkcia \(g(x)\) klesajúca. \\[1mm]
\textbf{(c)} Určte počet stacionárnych bodov funkcie \(h(t)\). \par \columnbreak
\textbf{(d)} Nadobúda funkcia \(v(x)\) lokálny extrém? \\[1mm]
\textbf{(e)} Určte rozdiel hodnôt lokálneho maxima a lokálneho minima funkcie \(\varphi(x)\). \\[1mm]
\textbf{(f)} Určte súčin \(x\)-ových súradníc bodov globálneho maxima a globálneho minima funkcie \(\omega(x)\) na zadanom intervale. \\[1mm]
\end{multicols}
}
\end{center}
\end{minipage}
\\[4mm]
\begin{minipage}{0.79\linewidth}
\begin{center}
\begin{varwidth}{\linewidth}
\begin{enumerate}
\normalsize
\item $f(x)=\frac{x}{4}+\frac{4}{x} \enspace , \enspace x_0=-10$\quad \dotfill\; ???\;\dotfill \quad \text{klesajúca}
\item $g(x)=x^2\, e^{-5x-6}$\quad \dotfill\; ???\;\dotfill \quad $\left\langle\frac{2}{5} , 0\right\rangle$
\item $h(t)=\cfrac{\ln t}{3t^4}$\quad \dotfill\; ???\;\dotfill \quad 1
\item $v(x)=-2-\sqrt[3]{(-8x-6)^2}$\quad \dotfill\; ???\;\dotfill \quad $\text{áno}$
\item $\varphi(x)=\cfrac{x^2-x+4}{x+3}$\quad \dotfill\; ???\;\dotfill \quad $-18$
\item $\omega(x)=-\frac{4}{3}x^3+\frac{33}{2}x^2-54x+2 \enspace , \enspace x\in\left\langle0 , 7\right\rangle$\quad \dotfill\; ???\;\dotfill \quad $0$
\end{enumerate}
\end{varwidth}
\end{center}
\end{minipage}
\begin{minipage}{0.20\linewidth}
\begin{center}
{\Huge\bfseries 1.} \\[2mm]
\includegraphics[height=37mm]{images/braille.png}
{\small Braillovo písmeno}
\end{center}
\end{minipage}
\end{center}
\end{minipage}
&
\begin{minipage}[c][104.5mm][t]{0.5\linewidth}
\begin{center}
\vspace{7mm}
{\huge Monotónnosť a extrémy, skupina \textit{Phi $\phi$} -\romannumeral2}\\[5mm]
\begin{minipage}{0.95\linewidth}
\begin{center}
{\small \begin{multicols}{2} \raggedright
\textbf{(a)} Určte monotónnosť funkcie \(f(x)\) v okolí bodu \(x_0\). \\[1mm]
\textbf{(b)} Určte intervaly, na ktorých je funkcia \(g(x)\) klesajúca. \\[1mm]
\textbf{(c)} Určte počet stacionárnych bodov funkcie \(h(t)\). \par \columnbreak
\textbf{(d)} Nadobúda funkcia \(v(x)\) lokálny extrém? \\[1mm]
\textbf{(e)} Určte rozdiel hodnôt lokálneho maxima a lokálneho minima funkcie \(\varphi(x)\). \\[1mm]
\textbf{(f)} Určte súčin \(x\)-ových súradníc bodov globálneho maxima a globálneho minima funkcie \(\omega(x)\) na zadanom intervale. \\[1mm]
\end{multicols}
}
\end{center}
\end{minipage}
\\[4mm]
\begin{minipage}{0.79\linewidth}
\begin{center}
\begin{varwidth}{\linewidth}
\begin{enumerate}
\normalsize
\item $f(x)=\frac{x}{10}+\frac{10}{x} \enspace , \enspace x_0=9$\quad \dotfill\; ???\;\dotfill \quad \text{klesajúca}
\item $g(x)=x^2\, e^{4x-4}$\quad \dotfill\; ???\;\dotfill \quad $\left\langle-\frac{1}{2} , 0\right\rangle$
\item $h(t)=\cfrac{\ln t}{2t^4}$\quad \dotfill\; ???\;\dotfill \quad 1
\item $v(x)=5-\sqrt[3]{(8x-9)^2}$\quad \dotfill\; ???\;\dotfill \quad $\text{nie}$
\item $\varphi(x)=\cfrac{x^2+4x-1}{x+5}$\quad \dotfill\; ???\;\dotfill \quad $-8$
\item $\omega(x)=\frac{1}{3}x^3-\frac{1}{2}x^2-2x-5 \enspace , \enspace x\in\left\langle-2 , 4\right\rangle$\quad \dotfill\; ???\;\dotfill \quad $4$
\end{enumerate}
\end{varwidth}
\end{center}
\end{minipage}
\begin{minipage}{0.20\linewidth}
\begin{center}
{\Huge\bfseries 2.} \\[2mm]
\includegraphics[height=37mm]{images/braille.png}
{\small Braillovo písmeno}
\end{center}
\end{minipage}
\end{center}
\end{minipage}
\\ \hdashline
\begin{minipage}[c][104.5mm][t]{0.5\linewidth}
\begin{center}
\vspace{7mm}
{\huge Monotónnosť a extrémy, skupina \textit{Phi $\phi$} -\romannumeral3}\\[5mm]
\begin{minipage}{0.95\linewidth}
\begin{center}
{\small \begin{multicols}{2} \raggedright
\textbf{(a)} Určte monotónnosť funkcie \(f(x)\) v okolí bodu \(x_0\). \\[1mm]
\textbf{(b)} Určte intervaly, na ktorých je funkcia \(g(x)\) klesajúca. \\[1mm]
\textbf{(c)} Určte počet stacionárnych bodov funkcie \(h(t)\). \par \columnbreak
\textbf{(d)} Nadobúda funkcia \(v(x)\) lokálny extrém? \\[1mm]
\textbf{(e)} Určte rozdiel hodnôt lokálneho maxima a lokálneho minima funkcie \(\varphi(x)\). \\[1mm]
\textbf{(f)} Určte súčin \(x\)-ových súradníc bodov globálneho maxima a globálneho minima funkcie \(\omega(x)\) na zadanom intervale. \\[1mm]
\end{multicols}
}
\end{center}
\end{minipage}
\\[4mm]
\begin{minipage}{0.79\linewidth}
\begin{center}
\begin{varwidth}{\linewidth}
\begin{enumerate}
\normalsize
\item $f(x)=\frac{x}{8}+\frac{8}{x} \enspace , \enspace x_0=3$\quad \dotfill\; ???\;\dotfill \quad \text{klesajúca}
\item $g(x)=x^2\, e^{-x-3}$\quad \dotfill\; ???\;\dotfill \quad $\left\langle2 , 0\right\rangle$
\item $h(t)=\cfrac{\ln t}{6t^3}$\quad \dotfill\; ???\;\dotfill \quad 0
\item $v(x)=-2-\sqrt[5]{(9x+5)^2}$\quad \dotfill\; ???\;\dotfill \quad $\text{nie}$
\item $\varphi(x)=\cfrac{x^2+x-5}{x+3}$\quad \dotfill\; ???\;\dotfill \quad $4$
\item $\omega(x)=3x^3-5x^2+x+2 \enspace , \enspace x\in\left\langle0 , 3\right\rangle$\quad \dotfill\; ???\;\dotfill \quad $1$
\end{enumerate}
\end{varwidth}
\end{center}
\end{minipage}
\begin{minipage}{0.20\linewidth}
\begin{center}
{\Huge\bfseries 3.} \\[2mm]
\includegraphics[height=37mm]{images/braille.png}
{\small Braillovo písmeno}
\end{center}
\end{minipage}
\end{center}
\end{minipage}
&
\begin{minipage}[c][104.5mm][t]{0.5\linewidth}
\begin{center}
\vspace{7mm}
{\huge Monotónnosť a extrémy, skupina \textit{Phi $\phi$} -\romannumeral4}\\[5mm]
\begin{minipage}{0.95\linewidth}
\begin{center}
{\small \begin{multicols}{2} \raggedright
\textbf{(a)} Určte monotónnosť funkcie \(f(x)\) v okolí bodu \(x_0\). \\[1mm]
\textbf{(b)} Určte intervaly, na ktorých je funkcia \(g(x)\) klesajúca. \\[1mm]
\textbf{(c)} Určte počet stacionárnych bodov funkcie \(h(t)\). \par \columnbreak
\textbf{(d)} Nadobúda funkcia \(v(x)\) lokálny extrém? \\[1mm]
\textbf{(e)} Určte rozdiel hodnôt lokálneho maxima a lokálneho minima funkcie \(\varphi(x)\). \\[1mm]
\textbf{(f)} Určte súčin \(x\)-ových súradníc bodov globálneho maxima a globálneho minima funkcie \(\omega(x)\) na zadanom intervale. \\[1mm]
\end{multicols}
}
\end{center}
\end{minipage}
\\[4mm]
\begin{minipage}{0.79\linewidth}
\begin{center}
\begin{varwidth}{\linewidth}
\begin{enumerate}
\normalsize
\item $f(x)=\frac{x}{4}+\frac{4}{x} \enspace , \enspace x_0=1$\quad \dotfill\; ???\;\dotfill \quad \text{klesajúca}
\item $g(x)=x^2\, e^{2x-4}$\quad \dotfill\; ???\;\dotfill \quad $\left(-\infty , 0\right\rangle\enspace , \enspace \left\langle-1 , \infty\right)$
\item $h(t)=\cfrac{\ln t}{t^3}$\quad \dotfill\; ???\;\dotfill \quad 1
\item $v(x)=-1-\sqrt[5]{(-2x-3)^2}$\quad \dotfill\; ???\;\dotfill \quad $\text{áno}$
\item $\varphi(x)=\cfrac{x^2-8x-5}{x+1}$\quad \dotfill\; ???\;\dotfill \quad $-8$
\item $\omega(x)=\frac{2}{3}x^3-\frac{13}{2}x^2+20x+6 \enspace , \enspace x\in\left\langle2 , 6\right\rangle$\quad \dotfill\; ???\;\dotfill \quad $20$
\end{enumerate}
\end{varwidth}
\end{center}
\end{minipage}
\begin{minipage}{0.20\linewidth}
\begin{center}
{\Huge\bfseries 4.} \\[2mm]
\includegraphics[height=37mm]{images/braille.png}
{\small Braillovo písmeno}
\end{center}
\end{minipage}
\end{center}
\end{minipage}
%
\end{tabular}
\newpage
\thispagestyle{empty}
\begin{tabular}{c:c}
\begin{minipage}[c][104.5mm][t]{0.5\linewidth}
\begin{center}
\vspace{7mm}
{\huge Monotónnosť a extrémy, skupina \textit{Chi $\chi$} -\romannumeral1}\\[5mm]
\begin{minipage}{0.95\linewidth}
\begin{center}
{\small \begin{multicols}{2} \raggedright
\textbf{(a)} Určte monotónnosť funkcie \(f(x)\) v okolí bodu \(x_0\). \\[1mm]
\textbf{(b)} Určte intervaly, na ktorých je funkcia \(g(x)\) klesajúca. \\[1mm]
\textbf{(c)} Určte počet stacionárnych bodov funkcie \(h(t)\). \par \columnbreak
\textbf{(d)} Nadobúda funkcia \(v(x)\) lokálny extrém? \\[1mm]
\textbf{(e)} Určte rozdiel hodnôt lokálneho maxima a lokálneho minima funkcie \(\varphi(x)\). \\[1mm]
\textbf{(f)} Určte súčin \(x\)-ových súradníc bodov globálneho maxima a globálneho minima funkcie \(\omega(x)\) na zadanom intervale. \\[1mm]
\end{multicols}
}
\end{center}
\end{minipage}
\\[4mm]
\begin{minipage}{0.79\linewidth}
\begin{center}
\begin{varwidth}{\linewidth}
\begin{enumerate}
\normalsize
\item $f(x)=\frac{x}{2}+\frac{2}{x} \enspace , \enspace x_0=-4$\quad \dotfill\; ???\;\dotfill \quad \text{rastúca}
\item $g(x)=x^2\, e^{3x-7}$\quad \dotfill\; ???\;\dotfill \quad $\left(-\infty , 0\right\rangle\enspace , \enspace \left\langle-\frac{2}{3} , \infty\right)$
\item $h(t)=\cfrac{\ln t}{2t^4}$\quad \dotfill\; ???\;\dotfill \quad 1
\item $v(x)=5-\sqrt[3]{(-5x+4)^2}$\quad \dotfill\; ???\;\dotfill \quad $\text{nie}$
\item $\varphi(x)=\cfrac{x^2-4x-11}{x+2}$\quad \dotfill\; ???\;\dotfill \quad $4$
\item $\omega(x)=-\frac{1}{3}x^3+6x^2-32x-1 \enspace , \enspace x\in\left\langle1 , 9\right\rangle$\quad \dotfill\; ???\;\dotfill \quad $4$
\end{enumerate}
\end{varwidth}
\end{center}
\end{minipage}
\begin{minipage}{0.20\linewidth}
\begin{center}
{\Huge\bfseries 1.} \\[2mm]
\includegraphics[height=37mm]{images/braille.png}
{\small Braillovo písmeno}
\end{center}
\end{minipage}
\end{center}
\end{minipage}
&
\begin{minipage}[c][104.5mm][t]{0.5\linewidth}
\begin{center}
\vspace{7mm}
{\huge Monotónnosť a extrémy, skupina \textit{Chi $\chi$} -\romannumeral2}\\[5mm]
\begin{minipage}{0.95\linewidth}
\begin{center}
{\small \begin{multicols}{2} \raggedright
\textbf{(a)} Určte monotónnosť funkcie \(f(x)\) v okolí bodu \(x_0\). \\[1mm]
\textbf{(b)} Určte intervaly, na ktorých je funkcia \(g(x)\) klesajúca. \\[1mm]
\textbf{(c)} Určte počet stacionárnych bodov funkcie \(h(t)\). \par \columnbreak
\textbf{(d)} Nadobúda funkcia \(v(x)\) lokálny extrém? \\[1mm]
\textbf{(e)} Určte rozdiel hodnôt lokálneho maxima a lokálneho minima funkcie \(\varphi(x)\). \\[1mm]
\textbf{(f)} Určte súčin \(x\)-ových súradníc bodov globálneho maxima a globálneho minima funkcie \(\omega(x)\) na zadanom intervale. \\[1mm]
\end{multicols}
}
\end{center}
\end{minipage}
\\[4mm]
\begin{minipage}{0.79\linewidth}
\begin{center}
\begin{varwidth}{\linewidth}
\begin{enumerate}
\normalsize
\item $f(x)=\frac{x}{5}+\frac{5}{x} \enspace , \enspace x_0=4$\quad \dotfill\; ???\;\dotfill \quad \text{klesajúca}
\item $g(x)=x^2\, e^{-4x-8}$\quad \dotfill\; ???\;\dotfill \quad $\left\langle\frac{1}{2} , 0\right\rangle$
\item $h(t)=\cfrac{\ln t}{2t^5}$\quad \dotfill\; ???\;\dotfill \quad 0
\item $v(x)=8-\sqrt[7]{(8x-2)^2}$\quad \dotfill\; ???\;\dotfill \quad $\text{nie}$
\item $\varphi(x)=\cfrac{x^2-x+16}{x-1}$\quad \dotfill\; ???\;\dotfill \quad $16$
\item $\omega(x)=\frac{2}{3}x^3+\frac{13}{2}x^2+18x-5 \enspace , \enspace x\in\left\langle-5 , 0\right\rangle$\quad \dotfill\; ???\;\dotfill \quad $2$
\end{enumerate}
\end{varwidth}
\end{center}
\end{minipage}
\begin{minipage}{0.20\linewidth}
\begin{center}
{\Huge\bfseries 2.} \\[2mm]
\includegraphics[height=37mm]{images/braille.png}
{\small Braillovo písmeno}
\end{center}
\end{minipage}
\end{center}
\end{minipage}
\\ \hdashline
\begin{minipage}[c][104.5mm][t]{0.5\linewidth}
\begin{center}
\vspace{7mm}
{\huge Monotónnosť a extrémy, skupina \textit{Chi $\chi$} -\romannumeral3}\\[5mm]
\begin{minipage}{0.95\linewidth}
\begin{center}
{\small \begin{multicols}{2} \raggedright
\textbf{(a)} Určte monotónnosť funkcie \(f(x)\) v okolí bodu \(x_0\). \\[1mm]
\textbf{(b)} Určte intervaly, na ktorých je funkcia \(g(x)\) klesajúca. \\[1mm]
\textbf{(c)} Určte počet stacionárnych bodov funkcie \(h(t)\). \par \columnbreak
\textbf{(d)} Nadobúda funkcia \(v(x)\) lokálny extrém? \\[1mm]
\textbf{(e)} Určte rozdiel hodnôt lokálneho maxima a lokálneho minima funkcie \(\varphi(x)\). \\[1mm]
\textbf{(f)} Určte súčin \(x\)-ových súradníc bodov globálneho maxima a globálneho minima funkcie \(\omega(x)\) na zadanom intervale. \\[1mm]
\end{multicols}
}
\end{center}
\end{minipage}
\\[4mm]
\begin{minipage}{0.79\linewidth}
\begin{center}
\begin{varwidth}{\linewidth}
\begin{enumerate}
\normalsize
\item $f(x)=\frac{x}{7}+\frac{7}{x} \enspace , \enspace x_0=4$\quad \dotfill\; ???\;\dotfill \quad \text{klesajúca}
\item $g(x)=x^2\, e^{-5x+8}$\quad \dotfill\; ???\;\dotfill \quad $\left(-\infty , 0\right\rangle\enspace , \enspace \left\langle\frac{2}{5} , \infty\right)$
\item $h(t)=\cfrac{\ln t}{3t^4}$\quad \dotfill\; ???\;\dotfill \quad 1
\item $v(x)=3-\sqrt[3]{(x+3)^2}$\quad \dotfill\; ???\;\dotfill \quad $\text{nie}$
\item $\varphi(x)=\cfrac{x^2-8x+37}{x-2}$\quad \dotfill\; ???\;\dotfill \quad $-18$
\item $\omega(x)=-\frac{2}{3}x^3+\frac{5}{2}x^2+3x-3 \enspace , \enspace x\in\left\langle-3 , 4\right\rangle$\quad \dotfill\; ???\;\dotfill \quad $\frac{3}{2}$
\end{enumerate}
\end{varwidth}
\end{center}
\end{minipage}
\begin{minipage}{0.20\linewidth}
\begin{center}
{\Huge\bfseries 3.} \\[2mm]
\includegraphics[height=37mm]{images/braille.png}
{\small Braillovo písmeno}
\end{center}
\end{minipage}
\end{center}
\end{minipage}
&
\begin{minipage}[c][104.5mm][t]{0.5\linewidth}
\begin{center}
\vspace{7mm}
{\huge Monotónnosť a extrémy, skupina \textit{Chi $\chi$} -\romannumeral4}\\[5mm]
\begin{minipage}{0.95\linewidth}
\begin{center}
{\small \begin{multicols}{2} \raggedright
\textbf{(a)} Určte monotónnosť funkcie \(f(x)\) v okolí bodu \(x_0\). \\[1mm]
\textbf{(b)} Určte intervaly, na ktorých je funkcia \(g(x)\) klesajúca. \\[1mm]
\textbf{(c)} Určte počet stacionárnych bodov funkcie \(h(t)\). \par \columnbreak
\textbf{(d)} Nadobúda funkcia \(v(x)\) lokálny extrém? \\[1mm]
\textbf{(e)} Určte rozdiel hodnôt lokálneho maxima a lokálneho minima funkcie \(\varphi(x)\). \\[1mm]
\textbf{(f)} Určte súčin \(x\)-ových súradníc bodov globálneho maxima a globálneho minima funkcie \(\omega(x)\) na zadanom intervale. \\[1mm]
\end{multicols}
}
\end{center}
\end{minipage}
\\[4mm]
\begin{minipage}{0.79\linewidth}
\begin{center}
\begin{varwidth}{\linewidth}
\begin{enumerate}
\normalsize
\item $f(x)=\frac{x}{9}+\frac{9}{x} \enspace , \enspace x_0=6$\quad \dotfill\; ???\;\dotfill \quad \text{klesajúca}
\item $g(x)=x^2\, e^{3x+1}$\quad \dotfill\; ???\;\dotfill \quad $\left(-\infty , 0\right\rangle\enspace , \enspace \left\langle-\frac{2}{3} , \infty\right)$
\item $h(t)=\cfrac{\ln t}{t^5}$\quad \dotfill\; ???\;\dotfill \quad 2
\item $v(x)=5-\sqrt[5]{(3x+4)^2}$\quad \dotfill\; ???\;\dotfill \quad $\text{nie}$
\item $\varphi(x)=\cfrac{x^2-5x+13}{x-4}$\quad \dotfill\; ???\;\dotfill \quad $12$
\item $\omega(x)=\frac{5}{3}x^3-8x^2+3x+2 \enspace , \enspace x\in\left\langle0 , 5\right\rangle$\quad \dotfill\; ???\;\dotfill \quad $12$
\end{enumerate}
\end{varwidth}
\end{center}
\end{minipage}
\begin{minipage}{0.20\linewidth}
\begin{center}
{\Huge\bfseries 4.} \\[2mm]
\includegraphics[height=37mm]{images/braille.png}
{\small Braillovo písmeno}
\end{center}
\end{minipage}
\end{center}
\end{minipage}
%
\end{tabular}
\newpage
\thispagestyle{empty}
\begin{tabular}{c:c}
\begin{minipage}[c][104.5mm][t]{0.5\linewidth}
\begin{center}
\vspace{7mm}
{\huge Monotónnosť a extrémy, skupina \textit{Psi $\psi$} -\romannumeral1}\\[5mm]
\begin{minipage}{0.95\linewidth}
\begin{center}
{\small \begin{multicols}{2} \raggedright
\textbf{(a)} Určte monotónnosť funkcie \(f(x)\) v okolí bodu \(x_0\). \\[1mm]
\textbf{(b)} Určte intervaly, na ktorých je funkcia \(g(x)\) klesajúca. \\[1mm]
\textbf{(c)} Určte počet stacionárnych bodov funkcie \(h(t)\). \par \columnbreak
\textbf{(d)} Nadobúda funkcia \(v(x)\) lokálny extrém? \\[1mm]
\textbf{(e)} Určte rozdiel hodnôt lokálneho maxima a lokálneho minima funkcie \(\varphi(x)\). \\[1mm]
\textbf{(f)} Určte súčin \(x\)-ových súradníc bodov globálneho maxima a globálneho minima funkcie \(\omega(x)\) na zadanom intervale. \\[1mm]
\end{multicols}
}
\end{center}
\end{minipage}
\\[4mm]
\begin{minipage}{0.79\linewidth}
\begin{center}
\begin{varwidth}{\linewidth}
\begin{enumerate}
\normalsize
\item $f(x)=\frac{x}{8}+\frac{8}{x} \enspace , \enspace x_0=7$\quad \dotfill\; ???\;\dotfill \quad \text{klesajúca}
\item $g(x)=x^2\, e^{9x+1}$\quad \dotfill\; ???\;\dotfill \quad $\left(-\infty , 0\right\rangle\enspace , \enspace \left\langle-\frac{2}{9} , \infty\right)$
\item $h(t)=\cfrac{\ln t}{t^3}$\quad \dotfill\; ???\;\dotfill \quad 1
\item $v(x)=-2-\sqrt[5]{(6x+4)^2}$\quad \dotfill\; ???\;\dotfill \quad $\text{áno}$
\item $\varphi(x)=\cfrac{x^2-4x+4}{x-5}$\quad \dotfill\; ???\;\dotfill \quad $12$
\item $\omega(x)=x^3+x^2-5x+3 \enspace , \enspace x\in\left\langle-2 , 3\right\rangle$\quad \dotfill\; ???\;\dotfill \quad $1$
\end{enumerate}
\end{varwidth}
\end{center}
\end{minipage}
\begin{minipage}{0.20\linewidth}
\begin{center}
{\Huge\bfseries 1.} \\[2mm]
\includegraphics[height=37mm]{images/braille.png}
{\small Braillovo písmeno}
\end{center}
\end{minipage}
\end{center}
\end{minipage}
&
\begin{minipage}[c][104.5mm][t]{0.5\linewidth}
\begin{center}
\vspace{7mm}
{\huge Monotónnosť a extrémy, skupina \textit{Psi $\psi$} -\romannumeral2}\\[5mm]
\begin{minipage}{0.95\linewidth}
\begin{center}
{\small \begin{multicols}{2} \raggedright
\textbf{(a)} Určte monotónnosť funkcie \(f(x)\) v okolí bodu \(x_0\). \\[1mm]
\textbf{(b)} Určte intervaly, na ktorých je funkcia \(g(x)\) klesajúca. \\[1mm]
\textbf{(c)} Určte počet stacionárnych bodov funkcie \(h(t)\). \par \columnbreak
\textbf{(d)} Nadobúda funkcia \(v(x)\) lokálny extrém? \\[1mm]
\textbf{(e)} Určte rozdiel hodnôt lokálneho maxima a lokálneho minima funkcie \(\varphi(x)\). \\[1mm]
\textbf{(f)} Určte súčin \(x\)-ových súradníc bodov globálneho maxima a globálneho minima funkcie \(\omega(x)\) na zadanom intervale. \\[1mm]
\end{multicols}
}
\end{center}
\end{minipage}
\\[4mm]
\begin{minipage}{0.79\linewidth}
\begin{center}
\begin{varwidth}{\linewidth}
\begin{enumerate}
\normalsize
\item $f(x)=\frac{x}{9}+\frac{9}{x} \enspace , \enspace x_0=-11$\quad \dotfill\; ???\;\dotfill \quad \text{rastúca}
\item $g(x)=x^2\, e^{9x-6}$\quad \dotfill\; ???\;\dotfill \quad $\left\langle-\frac{2}{9} , 0\right\rangle$
\item $h(t)=\cfrac{\ln t}{5t^3}$\quad \dotfill\; ???\;\dotfill \quad 2
\item $v(x)=-8-\sqrt[7]{(-x-7)^2}$\quad \dotfill\; ???\;\dotfill \quad $\text{nie}$
\item $\varphi(x)=\cfrac{x^2-6x+10}{x-3}$\quad \dotfill\; ???\;\dotfill \quad $-4$
\item $\omega(x)=-\frac{5}{3}x^3+8x^2-3x+3 \enspace , \enspace x\in\left\langle-2 , 4\right\rangle$\quad \dotfill\; ???\;\dotfill \quad $-\frac{2}{25}$
\end{enumerate}
\end{varwidth}
\end{center}
\end{minipage}
\begin{minipage}{0.20\linewidth}
\begin{center}
{\Huge\bfseries 2.} \\[2mm]
\includegraphics[height=37mm]{images/braille.png}
{\small Braillovo písmeno}
\end{center}
\end{minipage}
\end{center}
\end{minipage}
\\ \hdashline
\begin{minipage}[c][104.5mm][t]{0.5\linewidth}
\begin{center}
\vspace{7mm}
{\huge Monotónnosť a extrémy, skupina \textit{Psi $\psi$} -\romannumeral3}\\[5mm]
\begin{minipage}{0.95\linewidth}
\begin{center}
{\small \begin{multicols}{2} \raggedright
\textbf{(a)} Určte monotónnosť funkcie \(f(x)\) v okolí bodu \(x_0\). \\[1mm]
\textbf{(b)} Určte intervaly, na ktorých je funkcia \(g(x)\) klesajúca. \\[1mm]
\textbf{(c)} Určte počet stacionárnych bodov funkcie \(h(t)\). \par \columnbreak
\textbf{(d)} Nadobúda funkcia \(v(x)\) lokálny extrém? \\[1mm]
\textbf{(e)} Určte rozdiel hodnôt lokálneho maxima a lokálneho minima funkcie \(\varphi(x)\). \\[1mm]
\textbf{(f)} Určte súčin \(x\)-ových súradníc bodov globálneho maxima a globálneho minima funkcie \(\omega(x)\) na zadanom intervale. \\[1mm]
\end{multicols}
}
\end{center}
\end{minipage}
\\[4mm]
\begin{minipage}{0.79\linewidth}
\begin{center}
\begin{varwidth}{\linewidth}
\begin{enumerate}
\normalsize
\item $f(x)=\frac{x}{7}+\frac{7}{x} \enspace , \enspace x_0=-10$\quad \dotfill\; ???\;\dotfill \quad \text{rastúca}
\item $g(x)=x^2\, e^{-4x+1}$\quad \dotfill\; ???\;\dotfill \quad $\left\langle\frac{1}{2} , 0\right\rangle$
\item $h(t)=\cfrac{\ln t}{6t^5}$\quad \dotfill\; ???\;\dotfill \quad 1
\item $v(x)=-1-\sqrt[3]{(4x+5)^2}$\quad \dotfill\; ???\;\dotfill \quad $\text{nie}$
\item $\varphi(x)=\cfrac{x^2-5x+43}{x+1}$\quad \dotfill\; ???\;\dotfill \quad $-28$
\item $\omega(x)=-\frac{2}{3}x^3-\frac{5}{2}x^2+7x-6 \enspace , \enspace x\in\left\langle-6 , 2\right\rangle$\quad \dotfill\; ???\;\dotfill \quad $7$
\end{enumerate}
\end{varwidth}
\end{center}
\end{minipage}
\begin{minipage}{0.20\linewidth}
\begin{center}
{\Huge\bfseries 3.} \\[2mm]
\includegraphics[height=37mm]{images/braille.png}
{\small Braillovo písmeno}
\end{center}
\end{minipage}
\end{center}
\end{minipage}
&
\begin{minipage}[c][104.5mm][t]{0.5\linewidth}
\begin{center}
\vspace{7mm}
{\huge Monotónnosť a extrémy, skupina \textit{Psi $\psi$} -\romannumeral4}\\[5mm]
\begin{minipage}{0.95\linewidth}
\begin{center}
{\small \begin{multicols}{2} \raggedright
\textbf{(a)} Určte monotónnosť funkcie \(f(x)\) v okolí bodu \(x_0\). \\[1mm]
\textbf{(b)} Určte intervaly, na ktorých je funkcia \(g(x)\) klesajúca. \\[1mm]
\textbf{(c)} Určte počet stacionárnych bodov funkcie \(h(t)\). \par \columnbreak
\textbf{(d)} Nadobúda funkcia \(v(x)\) lokálny extrém? \\[1mm]
\textbf{(e)} Určte rozdiel hodnôt lokálneho maxima a lokálneho minima funkcie \(\varphi(x)\). \\[1mm]
\textbf{(f)} Určte súčin \(x\)-ových súradníc bodov globálneho maxima a globálneho minima funkcie \(\omega(x)\) na zadanom intervale. \\[1mm]
\end{multicols}
}
\end{center}
\end{minipage}
\\[4mm]
\begin{minipage}{0.79\linewidth}
\begin{center}
\begin{varwidth}{\linewidth}
\begin{enumerate}
\normalsize
\item $f(x)=\frac{x}{7}+\frac{7}{x} \enspace , \enspace x_0=5$\quad \dotfill\; ???\;\dotfill \quad \text{klesajúca}
\item $g(x)=x^2\, e^{-5x-1}$\quad \dotfill\; ???\;\dotfill \quad $\left(-\infty , 0\right\rangle\enspace , \enspace \left\langle\frac{2}{5} , \infty\right)$
\item $h(t)=\cfrac{\ln t}{7t^3}$\quad \dotfill\; ???\;\dotfill \quad 1
\item $v(x)=2-\sqrt[3]{(4x-2)^2}$\quad \dotfill\; ???\;\dotfill \quad $\text{nie}$
\item $\varphi(x)=\cfrac{x^2-x+34}{x+1}$\quad \dotfill\; ???\;\dotfill \quad $-22$
\item $\omega(x)=\frac{1}{3}x^3+3x^2+8x-9 \enspace , \enspace x\in\left\langle-5 , 0\right\rangle$\quad \dotfill\; ???\;\dotfill \quad $4$
\end{enumerate}
\end{varwidth}
\end{center}
\end{minipage}
\begin{minipage}{0.20\linewidth}
\begin{center}
{\Huge\bfseries 4.} \\[2mm]
\includegraphics[height=37mm]{images/braille.png}
{\small Braillovo písmeno}
\end{center}
\end{minipage}
\end{center}
\end{minipage}
%
\end{tabular}
\newpage
\thispagestyle{empty}
\begin{tabular}{c:c}
\begin{minipage}[c][104.5mm][t]{0.5\linewidth}
\begin{center}
\vspace{7mm}
{\huge Monotónnosť a extrémy, skupina \textit{Omega $\omega$} -\romannumeral1}\\[5mm]
\begin{minipage}{0.95\linewidth}
\begin{center}
{\small \begin{multicols}{2} \raggedright
\textbf{(a)} Určte monotónnosť funkcie \(f(x)\) v okolí bodu \(x_0\). \\[1mm]
\textbf{(b)} Určte intervaly, na ktorých je funkcia \(g(x)\) klesajúca. \\[1mm]
\textbf{(c)} Určte počet stacionárnych bodov funkcie \(h(t)\). \par \columnbreak
\textbf{(d)} Nadobúda funkcia \(v(x)\) lokálny extrém? \\[1mm]
\textbf{(e)} Určte rozdiel hodnôt lokálneho maxima a lokálneho minima funkcie \(\varphi(x)\). \\[1mm]
\textbf{(f)} Určte súčin \(x\)-ových súradníc bodov globálneho maxima a globálneho minima funkcie \(\omega(x)\) na zadanom intervale. \\[1mm]
\end{multicols}
}
\end{center}
\end{minipage}
\\[4mm]
\begin{minipage}{0.79\linewidth}
\begin{center}
\begin{varwidth}{\linewidth}
\begin{enumerate}
\normalsize
\item $f(x)=\frac{x}{4}+\frac{4}{x} \enspace , \enspace x_0=-7$\quad \dotfill\; ???\;\dotfill \quad \text{rastúca}
\item $g(x)=x^2\, e^{-5x+3}$\quad \dotfill\; ???\;\dotfill \quad $\left(-\infty , 0\right\rangle\enspace , \enspace \left\langle\frac{2}{5} , \infty\right)$
\item $h(t)=\cfrac{\ln t}{5t^4}$\quad \dotfill\; ???\;\dotfill \quad 1
\item $v(x)=-6-\sqrt[5]{(8x-5)^2}$\quad \dotfill\; ???\;\dotfill \quad $\text{áno}$
\item $\varphi(x)=\cfrac{x^2-5x+8}{x-1}$\quad \dotfill\; ???\;\dotfill \quad $-10$
\item $\omega(x)=-\frac{7}{3}x^3+4x^2+12x-2 \enspace , \enspace x\in\left\langle-3 , 3\right\rangle$\quad \dotfill\; ???\;\dotfill \quad $\frac{18}{49}$
\end{enumerate}
\end{varwidth}
\end{center}
\end{minipage}
\begin{minipage}{0.20\linewidth}
\begin{center}
{\Huge\bfseries 1.} \\[2mm]
\includegraphics[height=37mm]{images/braille.png}
{\small Braillovo písmeno}
\end{center}
\end{minipage}
\end{center}
\end{minipage}
&
\begin{minipage}[c][104.5mm][t]{0.5\linewidth}
\begin{center}
\vspace{7mm}
{\huge Monotónnosť a extrémy, skupina \textit{Omega $\omega$} -\romannumeral2}\\[5mm]
\begin{minipage}{0.95\linewidth}
\begin{center}
{\small \begin{multicols}{2} \raggedright
\textbf{(a)} Určte monotónnosť funkcie \(f(x)\) v okolí bodu \(x_0\). \\[1mm]
\textbf{(b)} Určte intervaly, na ktorých je funkcia \(g(x)\) klesajúca. \\[1mm]
\textbf{(c)} Určte počet stacionárnych bodov funkcie \(h(t)\). \par \columnbreak
\textbf{(d)} Nadobúda funkcia \(v(x)\) lokálny extrém? \\[1mm]
\textbf{(e)} Určte rozdiel hodnôt lokálneho maxima a lokálneho minima funkcie \(\varphi(x)\). \\[1mm]
\textbf{(f)} Určte súčin \(x\)-ových súradníc bodov globálneho maxima a globálneho minima funkcie \(\omega(x)\) na zadanom intervale. \\[1mm]
\end{multicols}
}
\end{center}
\end{minipage}
\\[4mm]
\begin{minipage}{0.79\linewidth}
\begin{center}
\begin{varwidth}{\linewidth}
\begin{enumerate}
\normalsize
\item $f(x)=\frac{x}{6}+\frac{6}{x} \enspace , \enspace x_0=-10$\quad \dotfill\; ???\;\dotfill \quad \text{rastúca}
\item $g(x)=x^2\, e^{x-2}$\quad \dotfill\; ???\;\dotfill \quad $\left\langle-2 , 0\right\rangle$
\item $h(t)=\cfrac{\ln t}{2t^3}$\quad \dotfill\; ???\;\dotfill \quad 1
\item $v(x)=3-\sqrt[3]{(-5x+7)^2}$\quad \dotfill\; ???\;\dotfill \quad $\text{nie}$
\item $\varphi(x)=\cfrac{x^2-8x+13}{x-6}$\quad \dotfill\; ???\;\dotfill \quad $-2$
\item $\omega(x)=-\frac{1}{3}x^3-4x^2-12x-6 \enspace , \enspace x\in\left\langle-9 , -1\right\rangle$\quad \dotfill\; ???\;\dotfill \quad $-30$
\end{enumerate}
\end{varwidth}
\end{center}
\end{minipage}
\begin{minipage}{0.20\linewidth}
\begin{center}
{\Huge\bfseries 2.} \\[2mm]
\includegraphics[height=37mm]{images/braille.png}
{\small Braillovo písmeno}
\end{center}
\end{minipage}
\end{center}
\end{minipage}
\\ \hdashline
\begin{minipage}[c][104.5mm][t]{0.5\linewidth}
\begin{center}
\vspace{7mm}
{\huge Monotónnosť a extrémy, skupina \textit{Omega $\omega$} -\romannumeral3}\\[5mm]
\begin{minipage}{0.95\linewidth}
\begin{center}
{\small \begin{multicols}{2} \raggedright
\textbf{(a)} Určte monotónnosť funkcie \(f(x)\) v okolí bodu \(x_0\). \\[1mm]
\textbf{(b)} Určte intervaly, na ktorých je funkcia \(g(x)\) klesajúca. \\[1mm]
\textbf{(c)} Určte počet stacionárnych bodov funkcie \(h(t)\). \par \columnbreak
\textbf{(d)} Nadobúda funkcia \(v(x)\) lokálny extrém? \\[1mm]
\textbf{(e)} Určte rozdiel hodnôt lokálneho maxima a lokálneho minima funkcie \(\varphi(x)\). \\[1mm]
\textbf{(f)} Určte súčin \(x\)-ových súradníc bodov globálneho maxima a globálneho minima funkcie \(\omega(x)\) na zadanom intervale. \\[1mm]
\end{multicols}
}
\end{center}
\end{minipage}
\\[4mm]
\begin{minipage}{0.79\linewidth}
\begin{center}
\begin{varwidth}{\linewidth}
\begin{enumerate}
\normalsize
\item $f(x)=\frac{x}{5}+\frac{5}{x} \enspace , \enspace x_0=3$\quad \dotfill\; ???\;\dotfill \quad \text{klesajúca}
\item $g(x)=x^2\, e^{4x-4}$\quad \dotfill\; ???\;\dotfill \quad $\left(-\infty , 0\right\rangle\enspace , \enspace \left\langle-\frac{1}{2} , \infty\right)$
\item $h(t)=\cfrac{\ln t}{t^4}$\quad \dotfill\; ???\;\dotfill \quad 2
\item $v(x)=-6-\sqrt[5]{(-4x-1)^2}$\quad \dotfill\; ???\;\dotfill \quad $\text{nie}$
\item $\varphi(x)=\cfrac{x^2+4x-36}{x-5}$\quad \dotfill\; ???\;\dotfill \quad $-12$
\item $\omega(x)=-\frac{1}{3}x^3+\frac{7}{2}x^2-6x-2 \enspace , \enspace x\in\left\langle-2 , 7\right\rangle$\quad \dotfill\; ???\;\dotfill \quad $-2$
\end{enumerate}
\end{varwidth}
\end{center}
\end{minipage}
\begin{minipage}{0.20\linewidth}
\begin{center}
{\Huge\bfseries 3.} \\[2mm]
\includegraphics[height=37mm]{images/braille.png}
{\small Braillovo písmeno}
\end{center}
\end{minipage}
\end{center}
\end{minipage}
&
\begin{minipage}[c][104.5mm][t]{0.5\linewidth}
\begin{center}
\vspace{7mm}
{\huge Monotónnosť a extrémy, skupina \textit{Omega $\omega$} -\romannumeral4}\\[5mm]
\begin{minipage}{0.95\linewidth}
\begin{center}
{\small \begin{multicols}{2} \raggedright
\textbf{(a)} Určte monotónnosť funkcie \(f(x)\) v okolí bodu \(x_0\). \\[1mm]
\textbf{(b)} Určte intervaly, na ktorých je funkcia \(g(x)\) klesajúca. \\[1mm]
\textbf{(c)} Určte počet stacionárnych bodov funkcie \(h(t)\). \par \columnbreak
\textbf{(d)} Nadobúda funkcia \(v(x)\) lokálny extrém? \\[1mm]
\textbf{(e)} Určte rozdiel hodnôt lokálneho maxima a lokálneho minima funkcie \(\varphi(x)\). \\[1mm]
\textbf{(f)} Určte súčin \(x\)-ových súradníc bodov globálneho maxima a globálneho minima funkcie \(\omega(x)\) na zadanom intervale. \\[1mm]
\end{multicols}
}
\end{center}
\end{minipage}
\\[4mm]
\begin{minipage}{0.79\linewidth}
\begin{center}
\begin{varwidth}{\linewidth}
\begin{enumerate}
\normalsize
\item $f(x)=\frac{x}{6}+\frac{6}{x} \enspace , \enspace x_0=4$\quad \dotfill\; ???\;\dotfill \quad \text{klesajúca}
\item $g(x)=x^2\, e^{x+1}$\quad \dotfill\; ???\;\dotfill \quad $\left(-\infty , 0\right\rangle\enspace , \enspace \left\langle-2 , \infty\right)$
\item $h(t)=\cfrac{\ln t}{4t^3}$\quad \dotfill\; ???\;\dotfill \quad 1
\item $v(x)=2-\sqrt[5]{(2x+3)^2}$\quad \dotfill\; ???\;\dotfill \quad $\text{nie}$
\item $\varphi(x)=\cfrac{x^2+5x+43}{x-1}$\quad \dotfill\; ???\;\dotfill \quad $28$
\item $\omega(x)=x^3-\frac{7}{2}x^2-6x-1 \enspace , \enspace x\in\left\langle-1 , 5\right\rangle$\quad \dotfill\; ???\;\dotfill \quad $9$
\end{enumerate}
\end{varwidth}
\end{center}
\end{minipage}
\begin{minipage}{0.20\linewidth}
\begin{center}
{\Huge\bfseries 4.} \\[2mm]
\includegraphics[height=37mm]{images/braille.png}
{\small Braillovo písmeno}
\end{center}
\end{minipage}
\end{center}
\end{minipage}
%
\end{tabular}
\newpage
\begin{landscape}
\newgeometry{total={194mm,285mm}, left=8mm, top=9mm}
\begin{center}
{\huge Monotónnosť a extrémy (riešenia)}\\[4mm]
\begin{varwidth}{\linewidth}
\begin{center}
\footnotesize
\rule[1mm]{\linewidth}{0.5pt}
$\boxed{\bm{\alpha}} \quad \begin{aligned}
\romannumeral1 : \; &\textbf{M} 
 &&\mathrm{\textbf{(a) }} \text{klesajúca}\,\text{\ding{51}}
 &&\mathrm{\textbf{(b) }} \left(-\infty , 0\right\rangle  ,   \left\langle2 , \infty\right)\,\text{\ding{55}}
 &&\mathrm{\textbf{(c) }} 1\,\text{\ding{51}}
 &&\mathrm{\textbf{(d) }} \text{áno}\,\text{\ding{55}}
 &&\mathrm{\textbf{(e) }} -20\,\text{\ding{55}}
 &&\mathrm{\textbf{(f) }} -\nicefrac{3}{2}\,\text{\ding{51}}
\\[-0.09999999999999998mm]
\romannumeral2 : \; &\textbf{L} 
 &&\mathrm{\textbf{(a) }} \text{klesajúca}\,\text{\ding{51}}
 &&\mathrm{\textbf{(b) }} \left\langle-\nicefrac{2}{7} , 0\right\rangle\,\text{\ding{51}}
 &&\mathrm{\textbf{(c) }} 1\,\text{\ding{51}}
 &&\mathrm{\textbf{(d) }} \text{áno}\,\text{\ding{55}}
 &&\mathrm{\textbf{(e) }} -4\,\text{\ding{55}}
 &&\mathrm{\textbf{(f) }} 3\,\text{\ding{55}}
\\[-0.09999999999999998mm]
\romannumeral3 : \; &\textbf{O} 
 &&\mathrm{\textbf{(a) }} \text{klesajúca}\,\text{\ding{51}}
 &&\mathrm{\textbf{(b) }} \left\langle-2 , 0\right\rangle\,\text{\ding{55}}
 &&\mathrm{\textbf{(c) }} 1\,\text{\ding{51}}
 &&\mathrm{\textbf{(d) }} \text{áno}\,\text{\ding{55}}
 &&\mathrm{\textbf{(e) }} -24\,\text{\ding{51}}
 &&\mathrm{\textbf{(f) }} 3\,\text{\ding{55}}
\\[-0.09999999999999998mm]
\romannumeral4 : \; &\textbf{K} 
 &&\mathrm{\textbf{(a) }} \text{klesajúca}\,\text{\ding{51}}
 &&\mathrm{\textbf{(b) }} \left\langle-\nicefrac{1}{3} , 0\right\rangle\,\text{\ding{55}}
 &&\mathrm{\textbf{(c) }} 1\,\text{\ding{51}}
 &&\mathrm{\textbf{(d) }} \text{áno}\,\text{\ding{55}}
 &&\mathrm{\textbf{(e) }} -20\,\text{\ding{55}}
 &&\mathrm{\textbf{(f) }} 3\,\text{\ding{55}}
\end{aligned} $
\\[2.9mm]
\rule[1.45mm]{\linewidth}{0.5pt}
$\boxed{\bm{\beta}} \quad \begin{aligned}
\romannumeral1 : \; &\textbf{V} 
 &&\mathrm{\textbf{(a) }} \text{rastúca}\,\text{\ding{51}}
 &&\mathrm{\textbf{(b) }} \left(-\infty , 0\right\rangle  ,   \left\langle\nicefrac{2}{7} , \infty\right)\,\text{\ding{51}}
 &&\mathrm{\textbf{(c) }} 1\,\text{\ding{51}}
 &&\mathrm{\textbf{(d) }} \text{áno}\,\text{\ding{51}}
 &&\mathrm{\textbf{(e) }} -4\,\text{\ding{55}}
 &&\mathrm{\textbf{(f) }} 3\,\text{\ding{55}}
\\[-0.09999999999999998mm]
\romannumeral2 : \; &\textbf{L} 
 &&\mathrm{\textbf{(a) }} \text{rastúca}\,\text{\ding{51}}
 &&\mathrm{\textbf{(b) }} \left\langle-2 , 0\right\rangle\,\text{\ding{51}}
 &&\mathrm{\textbf{(c) }} 1\,\text{\ding{51}}
 &&\mathrm{\textbf{(d) }} \text{áno}\,\text{\ding{55}}
 &&\mathrm{\textbf{(e) }} -4\,\text{\ding{55}}
 &&\mathrm{\textbf{(f) }} \nicefrac{5}{2}\,\text{\ding{55}}
\\[-0.09999999999999998mm]
\romannumeral3 : \; &\textbf{A} 
 &&\mathrm{\textbf{(a) }} \text{klesajúca}\,\text{\ding{51}}
 &&\mathrm{\textbf{(b) }} \left\langle-\nicefrac{2}{3} , 0\right\rangle\,\text{\ding{55}}
 &&\mathrm{\textbf{(c) }} 1\,\text{\ding{55}}
 &&\mathrm{\textbf{(d) }} \text{áno}\,\text{\ding{55}}
 &&\mathrm{\textbf{(e) }} -12\,\text{\ding{55}}
 &&\mathrm{\textbf{(f) }} 28\,\text{\ding{55}}
\\[-0.09999999999999998mm]
\romannumeral4 : \; &\textbf{K} 
 &&\mathrm{\textbf{(a) }} \text{rastúca}\,\text{\ding{51}}
 &&\mathrm{\textbf{(b) }} \left\langle-\nicefrac{1}{3} , 0\right\rangle\,\text{\ding{55}}
 &&\mathrm{\textbf{(c) }} 1\,\text{\ding{51}}
 &&\mathrm{\textbf{(d) }} \text{áno}\,\text{\ding{55}}
 &&\mathrm{\textbf{(e) }} -8\,\text{\ding{55}}
 &&\mathrm{\textbf{(f) }} -\nicefrac{2}{3}\,\text{\ding{55}}
\end{aligned} $
\\[2.9mm]
\rule[1.45mm]{\linewidth}{0.5pt}
$\boxed{\bm{\gamma}} \quad \begin{aligned}
\romannumeral1 : \; &\textbf{J} 
 &&\mathrm{\textbf{(a) }} \text{rastúca}\,\text{\ding{55}}
 &&\mathrm{\textbf{(b) }} \left(-\infty , 0\right\rangle  ,   \left\langle2 , \infty\right)\,\text{\ding{51}}
 &&\mathrm{\textbf{(c) }} 1\,\text{\ding{55}}
 &&\mathrm{\textbf{(d) }} \text{áno}\,\text{\ding{55}}
 &&\mathrm{\textbf{(e) }} -4\,\text{\ding{51}}
 &&\mathrm{\textbf{(f) }} 4\,\text{\ding{51}}
\\[-0.09999999999999998mm]
\romannumeral2 : \; &\textbf{A} 
 &&\mathrm{\textbf{(a) }} \text{klesajúca}\,\text{\ding{51}}
 &&\mathrm{\textbf{(b) }} \left(-\infty , 0\right\rangle  ,   \left\langle2 , \infty\right)\,\text{\ding{55}}
 &&\mathrm{\textbf{(c) }} 1\,\text{\ding{55}}
 &&\mathrm{\textbf{(d) }} \text{áno}\,\text{\ding{55}}
 &&\mathrm{\textbf{(e) }} -12\,\text{\ding{55}}
 &&\mathrm{\textbf{(f) }} 15\,\text{\ding{55}}
\\[-0.09999999999999998mm]
\romannumeral3 : \; &\textbf{M} 
 &&\mathrm{\textbf{(a) }} \text{rastúca}\,\text{\ding{51}}
 &&\mathrm{\textbf{(b) }} \left\langle-\nicefrac{2}{7} , 0\right\rangle\,\text{\ding{55}}
 &&\mathrm{\textbf{(c) }} 1\,\text{\ding{51}}
 &&\mathrm{\textbf{(d) }} \text{áno}\,\text{\ding{55}}
 &&\mathrm{\textbf{(e) }} -16\,\text{\ding{55}}
 &&\mathrm{\textbf{(f) }} 54\,\text{\ding{51}}
\\[-0.09999999999999998mm]
\romannumeral4 : \; &\textbf{A} 
 &&\mathrm{\textbf{(a) }} \text{klesajúca}\,\text{\ding{51}}
 &&\mathrm{\textbf{(b) }} \left(-\infty , 0\right\rangle  ,   \left\langle\nicefrac{1}{2} , \infty\right)\,\text{\ding{55}}
 &&\mathrm{\textbf{(c) }} 1\,\text{\ding{55}}
 &&\mathrm{\textbf{(d) }} \text{áno}\,\text{\ding{55}}
 &&\mathrm{\textbf{(e) }} -4\,\text{\ding{55}}
 &&\mathrm{\textbf{(f) }} 24\,\text{\ding{55}}
\end{aligned} $
\\[2.9mm]
\rule[1.45mm]{\linewidth}{0.5pt}
$\boxed{\bm{\delta}} \quad \begin{aligned}
\romannumeral1 : \; &\textbf{I} 
 &&\mathrm{\textbf{(a) }} \text{rastúca}\,\text{\ding{55}}
 &&\mathrm{\textbf{(b) }} \left\langle-\nicefrac{1}{3} , 0\right\rangle\,\text{\ding{51}}
 &&\mathrm{\textbf{(c) }} 1\,\text{\ding{55}}
 &&\mathrm{\textbf{(d) }} \text{áno}\,\text{\ding{55}}
 &&\mathrm{\textbf{(e) }} -4\,\text{\ding{55}}
 &&\mathrm{\textbf{(f) }} 3\,\text{\ding{51}}
\\[-0.09999999999999998mm]
\romannumeral2 : \; &\textbf{B} 
 &&\mathrm{\textbf{(a) }} \text{rastúca}\,\text{\ding{51}}
 &&\mathrm{\textbf{(b) }} \left(-\infty , 0\right\rangle  ,   \left\langle2 , \infty\right)\,\text{\ding{51}}
 &&\mathrm{\textbf{(c) }} 1\,\text{\ding{55}}
 &&\mathrm{\textbf{(d) }} \text{áno}\,\text{\ding{55}}
 &&\mathrm{\textbf{(e) }} -4\,\text{\ding{55}}
 &&\mathrm{\textbf{(f) }} 0\,\text{\ding{55}}
\\[-0.09999999999999998mm]
\romannumeral3 : \; &\textbf{I} 
 &&\mathrm{\textbf{(a) }} \text{rastúca}\,\text{\ding{55}}
 &&\mathrm{\textbf{(b) }} \left(-\infty , 0\right\rangle  ,   \left\langle\nicefrac{1}{3} , \infty\right)\,\text{\ding{51}}
 &&\mathrm{\textbf{(c) }} 1\,\text{\ding{55}}
 &&\mathrm{\textbf{(d) }} \text{áno}\,\text{\ding{55}}
 &&\mathrm{\textbf{(e) }} -16\,\text{\ding{55}}
 &&\mathrm{\textbf{(f) }} 15\,\text{\ding{51}}
\\[-0.09999999999999998mm]
\romannumeral4 : \; &\textbf{S} 
 &&\mathrm{\textbf{(a) }} \text{klesajúca}\,\text{\ding{55}}
 &&\mathrm{\textbf{(b) }} \left\langle-1 , 0\right\rangle\,\text{\ding{51}}
 &&\mathrm{\textbf{(c) }} 1\,\text{\ding{51}}
 &&\mathrm{\textbf{(d) }} \text{áno}\,\text{\ding{55}}
 &&\mathrm{\textbf{(e) }} -16\,\text{\ding{55}}
 &&\mathrm{\textbf{(f) }} -1\,\text{\ding{51}}
\end{aligned} $
\\[2.9mm]
\rule[1.45mm]{\linewidth}{0.5pt}
$\boxed{\bm{\epsilon}} \quad \begin{aligned}
\romannumeral1 : \; &\textbf{Č} 
 &&\mathrm{\textbf{(a) }} \text{rastúca}\,\text{\ding{51}}
 &&\mathrm{\textbf{(b) }} \left\langle-\nicefrac{2}{5} , 0\right\rangle\,\text{\ding{55}}
 &&\mathrm{\textbf{(c) }} 1\,\text{\ding{55}}
 &&\mathrm{\textbf{(d) }} \text{áno}\,\text{\ding{51}}
 &&\mathrm{\textbf{(e) }} -16\,\text{\ding{55}}
 &&\mathrm{\textbf{(f) }} 8\,\text{\ding{51}}
\\[-0.09999999999999998mm]
\romannumeral2 : \; &\textbf{E} 
 &&\mathrm{\textbf{(a) }} \text{klesajúca}\,\text{\ding{51}}
 &&\mathrm{\textbf{(b) }} \left\langle-\nicefrac{1}{2} , 0\right\rangle\,\text{\ding{55}}
 &&\mathrm{\textbf{(c) }} 1\,\text{\ding{55}}
 &&\mathrm{\textbf{(d) }} \text{áno}\,\text{\ding{55}}
 &&\mathrm{\textbf{(e) }} -20\,\text{\ding{51}}
 &&\mathrm{\textbf{(f) }} -\nicefrac{2}{5}\,\text{\ding{55}}
\\[-0.09999999999999998mm]
\romannumeral3 : \; &\textbf{S} 
 &&\mathrm{\textbf{(a) }} \text{klesajúca}\,\text{\ding{55}}
 &&\mathrm{\textbf{(b) }} \left\langle-1 , 0\right\rangle\,\text{\ding{51}}
 &&\mathrm{\textbf{(c) }} 1\,\text{\ding{51}}
 &&\mathrm{\textbf{(d) }} \text{áno}\,\text{\ding{55}}
 &&\mathrm{\textbf{(e) }} -16\,\text{\ding{55}}
 &&\mathrm{\textbf{(f) }} 48\,\text{\ding{51}}
\\[-0.09999999999999998mm]
\romannumeral4 : \; &\textbf{Ť} 
 &&\mathrm{\textbf{(a) }} \text{rastúca}\,\text{\ding{51}}
 &&\mathrm{\textbf{(b) }} \left(-\infty , 0\right\rangle  ,   \left\langle2 , \infty\right)\,\text{\ding{51}}
 &&\mathrm{\textbf{(c) }} 1\,\text{\ding{55}}
 &&\mathrm{\textbf{(d) }} \text{áno}\,\text{\ding{51}}
 &&\mathrm{\textbf{(e) }} -12\,\text{\ding{51}}
 &&\mathrm{\textbf{(f) }} 70\,\text{\ding{55}}
\end{aligned} $
\\[2.9mm]
\rule[1.45mm]{\linewidth}{0.5pt}
$\boxed{\bm{\zeta}} \quad \begin{aligned}
\romannumeral1 : \; &\textbf{T} 
 &&\mathrm{\textbf{(a) }} \text{rastúca}\,\text{\ding{55}}
 &&\mathrm{\textbf{(b) }} \left\langle-\nicefrac{2}{3} , 0\right\rangle\,\text{\ding{51}}
 &&\mathrm{\textbf{(c) }} 1\,\text{\ding{51}}
 &&\mathrm{\textbf{(d) }} \text{áno}\,\text{\ding{55}}
 &&\mathrm{\textbf{(e) }} -4\,\text{\ding{51}}
 &&\mathrm{\textbf{(f) }} 8\,\text{\ding{51}}
\\[-0.09999999999999998mm]
\romannumeral2 : \; &\textbf{A} 
 &&\mathrm{\textbf{(a) }} \text{rastúca}\,\text{\ding{51}}
 &&\mathrm{\textbf{(b) }} \left(-\infty , 0\right\rangle  ,   \left\langle\nicefrac{2}{3} , \infty\right)\,\text{\ding{55}}
 &&\mathrm{\textbf{(c) }} 1\,\text{\ding{55}}
 &&\mathrm{\textbf{(d) }} \text{áno}\,\text{\ding{55}}
 &&\mathrm{\textbf{(e) }} -16\,\text{\ding{55}}
 &&\mathrm{\textbf{(f) }} 28\,\text{\ding{55}}
\\[-0.09999999999999998mm]
\romannumeral3 : \; &\textbf{X} 
 &&\mathrm{\textbf{(a) }} \text{rastúca}\,\text{\ding{51}}
 &&\mathrm{\textbf{(b) }} \left\langle-\nicefrac{2}{5} , 0\right\rangle\,\text{\ding{55}}
 &&\mathrm{\textbf{(c) }} 1\,\text{\ding{51}}
 &&\mathrm{\textbf{(d) }} \text{áno}\,\text{\ding{51}}
 &&\mathrm{\textbf{(e) }} -12\,\text{\ding{55}}
 &&\mathrm{\textbf{(f) }} 24\,\text{\ding{51}}
\\[-0.09999999999999998mm]
\romannumeral4 : \; &\textbf{I} 
 &&\mathrm{\textbf{(a) }} \text{klesajúca}\,\text{\ding{55}}
 &&\mathrm{\textbf{(b) }} \left\langle-\nicefrac{1}{2} , 0\right\rangle\,\text{\ding{51}}
 &&\mathrm{\textbf{(c) }} 1\,\text{\ding{55}}
 &&\mathrm{\textbf{(d) }} \text{áno}\,\text{\ding{55}}
 &&\mathrm{\textbf{(e) }} -20\,\text{\ding{55}}
 &&\mathrm{\textbf{(f) }} 2\,\text{\ding{51}}
\end{aligned} $
\\[2.9mm]
\rule[1.45mm]{\linewidth}{0.5pt}
$\boxed{\bm{\eta}} \quad \begin{aligned}
\romannumeral1 : \; &\textbf{M} 
 &&\mathrm{\textbf{(a) }} \text{klesajúca}\,\text{\ding{51}}
 &&\mathrm{\textbf{(b) }} \left\langle-\nicefrac{1}{3} , 0\right\rangle\,\text{\ding{55}}
 &&\mathrm{\textbf{(c) }} 1\,\text{\ding{51}}
 &&\mathrm{\textbf{(d) }} \text{áno}\,\text{\ding{55}}
 &&\mathrm{\textbf{(e) }} -16\,\text{\ding{55}}
 &&\mathrm{\textbf{(f) }} 8\,\text{\ding{51}}
\\[-0.09999999999999998mm]
\romannumeral2 : \; &\textbf{R} 
 &&\mathrm{\textbf{(a) }} \text{klesajúca}\,\text{\ding{51}}
 &&\mathrm{\textbf{(b) }} \left(-\infty , 0\right\rangle  ,   \left\langle1 , \infty\right)\,\text{\ding{51}}
 &&\mathrm{\textbf{(c) }} 1\,\text{\ding{51}}
 &&\mathrm{\textbf{(d) }} \text{áno}\,\text{\ding{55}}
 &&\mathrm{\textbf{(e) }} -20\,\text{\ding{51}}
 &&\mathrm{\textbf{(f) }} 2\,\text{\ding{55}}
\\[-0.09999999999999998mm]
\romannumeral3 : \; &\textbf{A} 
 &&\mathrm{\textbf{(a) }} \text{klesajúca}\,\text{\ding{51}}
 &&\mathrm{\textbf{(b) }} \left(-\infty , 0\right\rangle  ,   \left\langle\nicefrac{1}{4} , \infty\right)\,\text{\ding{55}}
 &&\mathrm{\textbf{(c) }} 1\,\text{\ding{55}}
 &&\mathrm{\textbf{(d) }} \text{áno}\,\text{\ding{55}}
 &&\mathrm{\textbf{(e) }} -16\,\text{\ding{55}}
 &&\mathrm{\textbf{(f) }} \nicefrac{12}{5}\,\text{\ding{55}}
\\[-0.09999999999999998mm]
\romannumeral4 : \; &\textbf{K} 
 &&\mathrm{\textbf{(a) }} \text{klesajúca}\,\text{\ding{51}}
 &&\mathrm{\textbf{(b) }} \left(-\infty , 0\right\rangle  ,   \left\langle\nicefrac{1}{2} , \infty\right)\,\text{\ding{55}}
 &&\mathrm{\textbf{(c) }} 1\,\text{\ding{51}}
 &&\mathrm{\textbf{(d) }} \text{áno}\,\text{\ding{55}}
 &&\mathrm{\textbf{(e) }} -20\,\text{\ding{55}}
 &&\mathrm{\textbf{(f) }} 35\,\text{\ding{55}}
\end{aligned} $
\\[2.9mm]
\rule[1.45mm]{\linewidth}{0.5pt}
$\boxed{\bm{\theta}} \quad \begin{aligned}
\romannumeral1 : \; &\textbf{R} 
 &&\mathrm{\textbf{(a) }} \text{klesajúca}\,\text{\ding{51}}
 &&\mathrm{\textbf{(b) }} \left(-\infty , 0\right\rangle  ,   \left\langle\nicefrac{2}{5} , \infty\right)\,\text{\ding{51}}
 &&\mathrm{\textbf{(c) }} 1\,\text{\ding{51}}
 &&\mathrm{\textbf{(d) }} \text{áno}\,\text{\ding{55}}
 &&\mathrm{\textbf{(e) }} -28\,\text{\ding{51}}
 &&\mathrm{\textbf{(f) }} \nicefrac{12}{5}\,\text{\ding{55}}
\\[-0.09999999999999998mm]
\romannumeral2 : \; &\textbf{E} 
 &&\mathrm{\textbf{(a) }} \text{rastúca}\,\text{\ding{51}}
 &&\mathrm{\textbf{(b) }} \left\langle-1 , 0\right\rangle\,\text{\ding{55}}
 &&\mathrm{\textbf{(c) }} 1\,\text{\ding{55}}
 &&\mathrm{\textbf{(d) }} \text{áno}\,\text{\ding{55}}
 &&\mathrm{\textbf{(e) }} -4\,\text{\ding{51}}
 &&\mathrm{\textbf{(f) }} \nicefrac{8}{3}\,\text{\ding{55}}
\\[-0.09999999999999998mm]
\romannumeral3 : \; &\textbf{P} 
 &&\mathrm{\textbf{(a) }} \text{klesajúca}\,\text{\ding{51}}
 &&\mathrm{\textbf{(b) }} \left(-\infty , 0\right\rangle  ,   \left\langle\nicefrac{1}{3} , \infty\right)\,\text{\ding{51}}
 &&\mathrm{\textbf{(c) }} 1\,\text{\ding{51}}
 &&\mathrm{\textbf{(d) }} \text{áno}\,\text{\ding{55}}
 &&\mathrm{\textbf{(e) }} -4\,\text{\ding{55}}
 &&\mathrm{\textbf{(f) }} 54\,\text{\ding{51}}
\\[-0.09999999999999998mm]
\romannumeral4 : \; &\textbf{A} 
 &&\mathrm{\textbf{(a) }} \text{rastúca}\,\text{\ding{51}}
 &&\mathrm{\textbf{(b) }} \left(-\infty , 0\right\rangle  ,   \left\langle\nicefrac{1}{3} , \infty\right)\,\text{\ding{55}}
 &&\mathrm{\textbf{(c) }} 1\,\text{\ding{55}}
 &&\mathrm{\textbf{(d) }} \text{áno}\,\text{\ding{55}}
 &&\mathrm{\textbf{(e) }} -4\,\text{\ding{55}}
 &&\mathrm{\textbf{(f) }} 24\,\text{\ding{55}}
\end{aligned} $
\\[2.9mm]
\rule[1.45mm]{\linewidth}{0.5pt}
$\boxed{\bm{\iota}} \quad \begin{aligned}
\romannumeral1 : \; &\textbf{I} 
 &&\mathrm{\textbf{(a) }} \text{klesajúca}\,\text{\ding{55}}
 &&\mathrm{\textbf{(b) }} \left(-\infty , 0\right\rangle  ,   \left\langle\nicefrac{1}{3} , \infty\right)\,\text{\ding{51}}
 &&\mathrm{\textbf{(c) }} 1\,\text{\ding{55}}
 &&\mathrm{\textbf{(d) }} \text{áno}\,\text{\ding{55}}
 &&\mathrm{\textbf{(e) }} -16\,\text{\ding{55}}
 &&\mathrm{\textbf{(f) }} 15\,\text{\ding{51}}
\\[-0.09999999999999998mm]
\romannumeral2 : \; &\textbf{G} 
 &&\mathrm{\textbf{(a) }} \text{rastúca}\,\text{\ding{51}}
 &&\mathrm{\textbf{(b) }} \left(-\infty , 0\right\rangle  ,   \left\langle\nicefrac{2}{7} , \infty\right)\,\text{\ding{51}}
 &&\mathrm{\textbf{(c) }} 1\,\text{\ding{55}}
 &&\mathrm{\textbf{(d) }} \text{áno}\,\text{\ding{55}}
 &&\mathrm{\textbf{(e) }} -12\,\text{\ding{51}}
 &&\mathrm{\textbf{(f) }} \nicefrac{20}{3}\,\text{\ding{51}}
\\[-0.09999999999999998mm]
\romannumeral3 : \; &\textbf{O} 
 &&\mathrm{\textbf{(a) }} \text{rastúca}\,\text{\ding{51}}
 &&\mathrm{\textbf{(b) }} \left(-\infty , 0\right\rangle  ,   \left\langle\nicefrac{1}{2} , \infty\right)\,\text{\ding{55}}
 &&\mathrm{\textbf{(c) }} 1\,\text{\ding{51}}
 &&\mathrm{\textbf{(d) }} \text{áno}\,\text{\ding{55}}
 &&\mathrm{\textbf{(e) }} -4\,\text{\ding{51}}
 &&\mathrm{\textbf{(f) }} 6\,\text{\ding{55}}
\\[-0.09999999999999998mm]
\romannumeral4 : \; &\textbf{R} 
 &&\mathrm{\textbf{(a) }} \text{klesajúca}\,\text{\ding{51}}
 &&\mathrm{\textbf{(b) }} \left\langle-\nicefrac{1}{4} , 0\right\rangle\,\text{\ding{51}}
 &&\mathrm{\textbf{(c) }} 1\,\text{\ding{51}}
 &&\mathrm{\textbf{(d) }} \text{áno}\,\text{\ding{55}}
 &&\mathrm{\textbf{(e) }} -4\,\text{\ding{51}}
 &&\mathrm{\textbf{(f) }} -\nicefrac{7}{6}\,\text{\ding{55}}
\end{aligned} $
\\[2.9mm]
\rule[1.45mm]{\linewidth}{0.5pt}
$\boxed{\bm{\kappa}} \quad \begin{aligned}
\romannumeral1 : \; &\textbf{A} 
 &&\mathrm{\textbf{(a) }} \text{rastúca}\,\text{\ding{51}}
 &&\mathrm{\textbf{(b) }} \left(-\infty , 0\right\rangle  ,   \left\langle\nicefrac{1}{2} , \infty\right)\,\text{\ding{55}}
 &&\mathrm{\textbf{(c) }} 1\,\text{\ding{55}}
 &&\mathrm{\textbf{(d) }} \text{áno}\,\text{\ding{55}}
 &&\mathrm{\textbf{(e) }} -8\,\text{\ding{55}}
 &&\mathrm{\textbf{(f) }} 8\,\text{\ding{55}}
\\[-0.09999999999999998mm]
\romannumeral2 : \; &\textbf{U} 
 &&\mathrm{\textbf{(a) }} \text{rastúca}\,\text{\ding{51}}
 &&\mathrm{\textbf{(b) }} \left(-\infty , 0\right\rangle  ,   \left\langle1 , \infty\right)\,\text{\ding{55}}
 &&\mathrm{\textbf{(c) }} 1\,\text{\ding{51}}
 &&\mathrm{\textbf{(d) }} \text{áno}\,\text{\ding{51}}
 &&\mathrm{\textbf{(e) }} -16\,\text{\ding{55}}
 &&\mathrm{\textbf{(f) }} \nicefrac{9}{4}\,\text{\ding{55}}
\\[-0.09999999999999998mm]
\romannumeral3 : \; &\textbf{T} 
 &&\mathrm{\textbf{(a) }} \text{klesajúca}\,\text{\ding{55}}
 &&\mathrm{\textbf{(b) }} \left(-\infty , 0\right\rangle  ,   \left\langle1 , \infty\right)\,\text{\ding{51}}
 &&\mathrm{\textbf{(c) }} 1\,\text{\ding{51}}
 &&\mathrm{\textbf{(d) }} \text{áno}\,\text{\ding{55}}
 &&\mathrm{\textbf{(e) }} -8\,\text{\ding{51}}
 &&\mathrm{\textbf{(f) }} 15\,\text{\ding{51}}
\\[-0.09999999999999998mm]
\romannumeral4 : \; &\textbf{O} 
 &&\mathrm{\textbf{(a) }} \text{rastúca}\,\text{\ding{51}}
 &&\mathrm{\textbf{(b) }} \left\langle-\nicefrac{2}{5} , 0\right\rangle\,\text{\ding{55}}
 &&\mathrm{\textbf{(c) }} 1\,\text{\ding{51}}
 &&\mathrm{\textbf{(d) }} \text{áno}\,\text{\ding{55}}
 &&\mathrm{\textbf{(e) }} -20\,\text{\ding{51}}
 &&\mathrm{\textbf{(f) }} 3\,\text{\ding{55}}
\end{aligned} $
\\[2.9mm]
\rule[1.45mm]{\linewidth}{0.5pt}
$\boxed{\bm{\lambda}} \quad \begin{aligned}
\romannumeral1 : \; &\textbf{T} 
 &&\mathrm{\textbf{(a) }} \text{rastúca}\,\text{\ding{55}}
 &&\mathrm{\textbf{(b) }} \left\langle-\nicefrac{1}{2} , 0\right\rangle\,\text{\ding{51}}
 &&\mathrm{\textbf{(c) }} 1\,\text{\ding{51}}
 &&\mathrm{\textbf{(d) }} \text{áno}\,\text{\ding{55}}
 &&\mathrm{\textbf{(e) }} -16\,\text{\ding{51}}
 &&\mathrm{\textbf{(f) }} 3\,\text{\ding{51}}
\\[-0.09999999999999998mm]
\romannumeral2 : \; &\textbf{L} 
 &&\mathrm{\textbf{(a) }} \text{rastúca}\,\text{\ding{51}}
 &&\mathrm{\textbf{(b) }} \left(-\infty , 0\right\rangle  ,   \left\langle\nicefrac{2}{7} , \infty\right)\,\text{\ding{51}}
 &&\mathrm{\textbf{(c) }} 1\,\text{\ding{51}}
 &&\mathrm{\textbf{(d) }} \text{áno}\,\text{\ding{55}}
 &&\mathrm{\textbf{(e) }} -20\,\text{\ding{55}}
 &&\mathrm{\textbf{(f) }} \nicefrac{25}{2}\,\text{\ding{55}}
\\[-0.09999999999999998mm]
\romannumeral3 : \; &\textbf{A} 
 &&\mathrm{\textbf{(a) }} \text{rastúca}\,\text{\ding{51}}
 &&\mathrm{\textbf{(b) }} \left(-\infty , 0\right\rangle  ,   \left\langle1 , \infty\right)\,\text{\ding{55}}
 &&\mathrm{\textbf{(c) }} 1\,\text{\ding{55}}
 &&\mathrm{\textbf{(d) }} \text{áno}\,\text{\ding{55}}
 &&\mathrm{\textbf{(e) }} -4\,\text{\ding{55}}
 &&\mathrm{\textbf{(f) }} 24\,\text{\ding{55}}
\\[-0.09999999999999998mm]
\romannumeral4 : \; &\textbf{K} 
 &&\mathrm{\textbf{(a) }} \text{rastúca}\,\text{\ding{51}}
 &&\mathrm{\textbf{(b) }} \left\langle-\nicefrac{2}{5} , 0\right\rangle\,\text{\ding{55}}
 &&\mathrm{\textbf{(c) }} 1\,\text{\ding{51}}
 &&\mathrm{\textbf{(d) }} \text{áno}\,\text{\ding{55}}
 &&\mathrm{\textbf{(e) }} -16\,\text{\ding{55}}
 &&\mathrm{\textbf{(f) }} 24\,\text{\ding{55}}
\end{aligned} $
\\[2.9mm]
\rule[1.45mm]{\linewidth}{0.5pt}
$\boxed{\bm{\mu}} \quad \begin{aligned}
\romannumeral1 : \; &\textbf{P} 
 &&\mathrm{\textbf{(a) }} \text{klesajúca}\,\text{\ding{51}}
 &&\mathrm{\textbf{(b) }} \left(-\infty , 0\right\rangle  ,   \left\langle\nicefrac{2}{3} , \infty\right)\,\text{\ding{51}}
 &&\mathrm{\textbf{(c) }} 1\,\text{\ding{51}}
 &&\mathrm{\textbf{(d) }} \text{áno}\,\text{\ding{55}}
 &&\mathrm{\textbf{(e) }} -28\,\text{\ding{55}}
 &&\mathrm{\textbf{(f) }} 4\,\text{\ding{51}}
\\[-0.09999999999999998mm]
\romannumeral2 : \; &\textbf{E} 
 &&\mathrm{\textbf{(a) }} \text{rastúca}\,\text{\ding{51}}
 &&\mathrm{\textbf{(b) }} \left\langle-\nicefrac{2}{7} , 0\right\rangle\,\text{\ding{55}}
 &&\mathrm{\textbf{(c) }} 1\,\text{\ding{55}}
 &&\mathrm{\textbf{(d) }} \text{áno}\,\text{\ding{55}}
 &&\mathrm{\textbf{(e) }} -12\,\text{\ding{51}}
 &&\mathrm{\textbf{(f) }} -1\,\text{\ding{55}}
\\[-0.09999999999999998mm]
\romannumeral3 : \; &\textbf{N} 
 &&\mathrm{\textbf{(a) }} \text{klesajúca}\,\text{\ding{51}}
 &&\mathrm{\textbf{(b) }} \left(-\infty , 0\right\rangle  ,   \left\langle\nicefrac{1}{3} , \infty\right)\,\text{\ding{55}}
 &&\mathrm{\textbf{(c) }} 1\,\text{\ding{51}}
 &&\mathrm{\textbf{(d) }} \text{áno}\,\text{\ding{55}}
 &&\mathrm{\textbf{(e) }} -8\,\text{\ding{51}}
 &&\mathrm{\textbf{(f) }} 15\,\text{\ding{51}}
\\[-0.09999999999999998mm]
\romannumeral4 : \; &\textbf{A} 
 &&\mathrm{\textbf{(a) }} \text{rastúca}\,\text{\ding{51}}
 &&\mathrm{\textbf{(b) }} \left(-\infty , 0\right\rangle  ,   \left\langle\nicefrac{2}{5} , \infty\right)\,\text{\ding{55}}
 &&\mathrm{\textbf{(c) }} 1\,\text{\ding{55}}
 &&\mathrm{\textbf{(d) }} \text{áno}\,\text{\ding{55}}
 &&\mathrm{\textbf{(e) }} -4\,\text{\ding{55}}
 &&\mathrm{\textbf{(f) }} \nicefrac{40}{3}\,\text{\ding{55}}
\end{aligned} $
\\[2.9mm]
\rule[1.45mm]{\linewidth}{0.5pt}
\end{center}\end{varwidth}\end{center}\clearpage\begin{center}{\huge Monotónnosť a extrémy (riešenia)}\\[4mm]\begin{varwidth}{\linewidth}\begin{center}
\vspace{--2.4000000000000004mm}
\footnotesize
\rule[1mm]{\linewidth}{0.5pt}
$\boxed{\bm{\nu}} \quad \begin{aligned}
\romannumeral1 : \; &\textbf{O} 
 &&\mathrm{\textbf{(a) }} \text{rastúca}\,\text{\ding{51}}
 &&\mathrm{\textbf{(b) }} \left(-\infty , 0\right\rangle  ,   \left\langle\nicefrac{2}{5} , \infty\right)\,\text{\ding{55}}
 &&\mathrm{\textbf{(c) }} 1\,\text{\ding{51}}
 &&\mathrm{\textbf{(d) }} \text{áno}\,\text{\ding{55}}
 &&\mathrm{\textbf{(e) }} -16\,\text{\ding{51}}
 &&\mathrm{\textbf{(f) }} 35\,\text{\ding{55}}
\\[-0.09999999999999998mm]
\romannumeral2 : \; &\textbf{L} 
 &&\mathrm{\textbf{(a) }} \text{klesajúca}\,\text{\ding{51}}
 &&\mathrm{\textbf{(b) }} \left\langle-\nicefrac{1}{4} , 0\right\rangle\,\text{\ding{51}}
 &&\mathrm{\textbf{(c) }} 1\,\text{\ding{51}}
 &&\mathrm{\textbf{(d) }} \text{áno}\,\text{\ding{55}}
 &&\mathrm{\textbf{(e) }} -8\,\text{\ding{55}}
 &&\mathrm{\textbf{(f) }} -\nicefrac{1}{3}\,\text{\ding{55}}
\\[-0.09999999999999998mm]
\romannumeral3 : \; &\textbf{E} 
 &&\mathrm{\textbf{(a) }} \text{rastúca}\,\text{\ding{51}}
 &&\mathrm{\textbf{(b) }} \left\langle-1 , 0\right\rangle\,\text{\ding{55}}
 &&\mathrm{\textbf{(c) }} 1\,\text{\ding{55}}
 &&\mathrm{\textbf{(d) }} \text{áno}\,\text{\ding{55}}
 &&\mathrm{\textbf{(e) }} -8\,\text{\ding{51}}
 &&\mathrm{\textbf{(f) }} 24\,\text{\ding{55}}
\\[-0.09999999999999998mm]
\romannumeral4 : \; &\textbf{J} 
 &&\mathrm{\textbf{(a) }} \text{klesajúca}\,\text{\ding{55}}
 &&\mathrm{\textbf{(b) }} \left\langle-\nicefrac{1}{4} , 0\right\rangle\,\text{\ding{51}}
 &&\mathrm{\textbf{(c) }} 1\,\text{\ding{55}}
 &&\mathrm{\textbf{(d) }} \text{áno}\,\text{\ding{55}}
 &&\mathrm{\textbf{(e) }} -4\,\text{\ding{51}}
 &&\mathrm{\textbf{(f) }} 8\,\text{\ding{51}}
\end{aligned} $
\\[2.9mm]
\rule[1.45mm]{\linewidth}{0.5pt}
$\boxed{\bm{\xi}} \quad \begin{aligned}
\romannumeral1 : \; &\textbf{S} 
 &&\mathrm{\textbf{(a) }} \text{rastúca}\,\text{\ding{55}}
 &&\mathrm{\textbf{(b) }} \left(-\infty , 0\right\rangle  ,   \left\langle2 , \infty\right)\,\text{\ding{51}}
 &&\mathrm{\textbf{(c) }} 1\,\text{\ding{51}}
 &&\mathrm{\textbf{(d) }} \text{áno}\,\text{\ding{55}}
 &&\mathrm{\textbf{(e) }} -4\,\text{\ding{55}}
 &&\mathrm{\textbf{(f) }} \nicefrac{3}{2}\,\text{\ding{51}}
\\[-0.09999999999999998mm]
\romannumeral2 : \; &\textbf{O} 
 &&\mathrm{\textbf{(a) }} \text{rastúca}\,\text{\ding{51}}
 &&\mathrm{\textbf{(b) }} \left\langle-\nicefrac{1}{2} , 0\right\rangle\,\text{\ding{55}}
 &&\mathrm{\textbf{(c) }} 1\,\text{\ding{51}}
 &&\mathrm{\textbf{(d) }} \text{áno}\,\text{\ding{55}}
 &&\mathrm{\textbf{(e) }} -4\,\text{\ding{51}}
 &&\mathrm{\textbf{(f) }} \nicefrac{20}{3}\,\text{\ding{55}}
\\[-0.09999999999999998mm]
\romannumeral3 : \; &\textbf{V} 
 &&\mathrm{\textbf{(a) }} \text{klesajúca}\,\text{\ding{51}}
 &&\mathrm{\textbf{(b) }} \left\langle-2 , 0\right\rangle\,\text{\ding{51}}
 &&\mathrm{\textbf{(c) }} 1\,\text{\ding{51}}
 &&\mathrm{\textbf{(d) }} \text{áno}\,\text{\ding{51}}
 &&\mathrm{\textbf{(e) }} -12\,\text{\ding{55}}
 &&\mathrm{\textbf{(f) }} 28\,\text{\ding{55}}
\\[-0.09999999999999998mm]
\romannumeral4 : \; &\textbf{A} 
 &&\mathrm{\textbf{(a) }} \text{klesajúca}\,\text{\ding{51}}
 &&\mathrm{\textbf{(b) }} \left(-\infty , 0\right\rangle  ,   \left\langle\nicefrac{2}{5} , \infty\right)\,\text{\ding{55}}
 &&\mathrm{\textbf{(c) }} 1\,\text{\ding{55}}
 &&\mathrm{\textbf{(d) }} \text{áno}\,\text{\ding{55}}
 &&\mathrm{\textbf{(e) }} -8\,\text{\ding{55}}
 &&\mathrm{\textbf{(f) }} 15\,\text{\ding{55}}
\end{aligned} $
\\[2.9mm]
\rule[1.45mm]{\linewidth}{0.5pt}
$\boxed{\bm{\omicron}} \quad \begin{aligned}
\romannumeral1 : \; &\textbf{C} 
 &&\mathrm{\textbf{(a) }} \text{klesajúca}\,\text{\ding{51}}
 &&\mathrm{\textbf{(b) }} \left\langle-2 , 0\right\rangle\,\text{\ding{55}}
 &&\mathrm{\textbf{(c) }} 1\,\text{\ding{55}}
 &&\mathrm{\textbf{(d) }} \text{áno}\,\text{\ding{55}}
 &&\mathrm{\textbf{(e) }} -16\,\text{\ding{55}}
 &&\mathrm{\textbf{(f) }} 0\,\text{\ding{51}}
\\[-0.09999999999999998mm]
\romannumeral2 : \; &\textbf{E} 
 &&\mathrm{\textbf{(a) }} \text{klesajúca}\,\text{\ding{51}}
 &&\mathrm{\textbf{(b) }} \left\langle-2 , 0\right\rangle\,\text{\ding{55}}
 &&\mathrm{\textbf{(c) }} 1\,\text{\ding{55}}
 &&\mathrm{\textbf{(d) }} \text{áno}\,\text{\ding{55}}
 &&\mathrm{\textbf{(e) }} -4\,\text{\ding{51}}
 &&\mathrm{\textbf{(f) }} -2\,\text{\ding{55}}
\\[-0.09999999999999998mm]
\romannumeral3 : \; &\textbf{L} 
 &&\mathrm{\textbf{(a) }} \text{rastúca}\,\text{\ding{51}}
 &&\mathrm{\textbf{(b) }} \left\langle-2 , 0\right\rangle\,\text{\ding{51}}
 &&\mathrm{\textbf{(c) }} 1\,\text{\ding{51}}
 &&\mathrm{\textbf{(d) }} \text{áno}\,\text{\ding{55}}
 &&\mathrm{\textbf{(e) }} -8\,\text{\ding{55}}
 &&\mathrm{\textbf{(f) }} -1\,\text{\ding{55}}
\\[-0.09999999999999998mm]
\romannumeral4 : \; &\textbf{A} 
 &&\mathrm{\textbf{(a) }} \text{klesajúca}\,\text{\ding{51}}
 &&\mathrm{\textbf{(b) }} \left(-\infty , 0\right\rangle  ,   \left\langle\nicefrac{2}{3} , \infty\right)\,\text{\ding{55}}
 &&\mathrm{\textbf{(c) }} 1\,\text{\ding{55}}
 &&\mathrm{\textbf{(d) }} \text{áno}\,\text{\ding{55}}
 &&\mathrm{\textbf{(e) }} -8\,\text{\ding{55}}
 &&\mathrm{\textbf{(f) }} 0\,\text{\ding{55}}
\end{aligned} $
\\[2.9mm]
\rule[1.45mm]{\linewidth}{0.5pt}
$\boxed{\bm{\pi}} \quad \begin{aligned}
\romannumeral1 : \; &\textbf{S} 
 &&\mathrm{\textbf{(a) }} \text{klesajúca}\,\text{\ding{55}}
 &&\mathrm{\textbf{(b) }} \left(-\infty , 0\right\rangle  ,   \left\langle2 , \infty\right)\,\text{\ding{51}}
 &&\mathrm{\textbf{(c) }} 1\,\text{\ding{51}}
 &&\mathrm{\textbf{(d) }} \text{áno}\,\text{\ding{55}}
 &&\mathrm{\textbf{(e) }} -16\,\text{\ding{55}}
 &&\mathrm{\textbf{(f) }} -\nicefrac{7}{6}\,\text{\ding{51}}
\\[-0.09999999999999998mm]
\romannumeral2 : \; &\textbf{L} 
 &&\mathrm{\textbf{(a) }} \text{rastúca}\,\text{\ding{51}}
 &&\mathrm{\textbf{(b) }} \left\langle-\nicefrac{2}{7} , 0\right\rangle\,\text{\ding{51}}
 &&\mathrm{\textbf{(c) }} 1\,\text{\ding{51}}
 &&\mathrm{\textbf{(d) }} \text{áno}\,\text{\ding{55}}
 &&\mathrm{\textbf{(e) }} -8\,\text{\ding{55}}
 &&\mathrm{\textbf{(f) }} 54\,\text{\ding{55}}
\\[-0.09999999999999998mm]
\romannumeral3 : \; &\textbf{O} 
 &&\mathrm{\textbf{(a) }} \text{rastúca}\,\text{\ding{51}}
 &&\mathrm{\textbf{(b) }} \left\langle-\nicefrac{2}{5} , 0\right\rangle\,\text{\ding{55}}
 &&\mathrm{\textbf{(c) }} 1\,\text{\ding{51}}
 &&\mathrm{\textbf{(d) }} \text{áno}\,\text{\ding{55}}
 &&\mathrm{\textbf{(e) }} -16\,\text{\ding{51}}
 &&\mathrm{\textbf{(f) }} -2\,\text{\ding{55}}
\\[-0.09999999999999998mm]
\romannumeral4 : \; &\textbf{N} 
 &&\mathrm{\textbf{(a) }} \text{klesajúca}\,\text{\ding{51}}
 &&\mathrm{\textbf{(b) }} \left(-\infty , 0\right\rangle  ,   \left\langle\nicefrac{1}{3} , \infty\right)\,\text{\ding{55}}
 &&\mathrm{\textbf{(c) }} 1\,\text{\ding{51}}
 &&\mathrm{\textbf{(d) }} \text{áno}\,\text{\ding{55}}
 &&\mathrm{\textbf{(e) }} -16\,\text{\ding{51}}
 &&\mathrm{\textbf{(f) }} 24\,\text{\ding{51}}
\end{aligned} $
\\[2.9mm]
\rule[1.45mm]{\linewidth}{0.5pt}
$\boxed{\bm{\rho}} \quad \begin{aligned}
\romannumeral1 : \; &\textbf{P} 
 &&\mathrm{\textbf{(a) }} \text{klesajúca}\,\text{\ding{51}}
 &&\mathrm{\textbf{(b) }} \left(-\infty , 0\right\rangle  ,   \left\langle2 , \infty\right)\,\text{\ding{51}}
 &&\mathrm{\textbf{(c) }} 1\,\text{\ding{51}}
 &&\mathrm{\textbf{(d) }} \text{áno}\,\text{\ding{55}}
 &&\mathrm{\textbf{(e) }} -12\,\text{\ding{55}}
 &&\mathrm{\textbf{(f) }} 80\,\text{\ding{51}}
\\[-0.09999999999999998mm]
\romannumeral2 : \; &\textbf{R} 
 &&\mathrm{\textbf{(a) }} \text{klesajúca}\,\text{\ding{51}}
 &&\mathrm{\textbf{(b) }} \left\langle-\nicefrac{2}{7} , 0\right\rangle\,\text{\ding{51}}
 &&\mathrm{\textbf{(c) }} 1\,\text{\ding{51}}
 &&\mathrm{\textbf{(d) }} \text{áno}\,\text{\ding{55}}
 &&\mathrm{\textbf{(e) }} -4\,\text{\ding{51}}
 &&\mathrm{\textbf{(f) }} 24\,\text{\ding{55}}
\\[-0.09999999999999998mm]
\romannumeral3 : \; &\textbf{S} 
 &&\mathrm{\textbf{(a) }} \text{rastúca}\,\text{\ding{55}}
 &&\mathrm{\textbf{(b) }} \left\langle-2 , 0\right\rangle\,\text{\ding{51}}
 &&\mathrm{\textbf{(c) }} 1\,\text{\ding{51}}
 &&\mathrm{\textbf{(d) }} \text{áno}\,\text{\ding{55}}
 &&\mathrm{\textbf{(e) }} -4\,\text{\ding{55}}
 &&\mathrm{\textbf{(f) }} \nicefrac{25}{2}\,\text{\ding{51}}
\\[-0.09999999999999998mm]
\romannumeral4 : \; &\textbf{T} 
 &&\mathrm{\textbf{(a) }} \text{rastúca}\,\text{\ding{55}}
 &&\mathrm{\textbf{(b) }} \left(-\infty , 0\right\rangle  ,   \left\langle2 , \infty\right)\,\text{\ding{51}}
 &&\mathrm{\textbf{(c) }} 1\,\text{\ding{51}}
 &&\mathrm{\textbf{(d) }} \text{áno}\,\text{\ding{55}}
 &&\mathrm{\textbf{(e) }} -12\,\text{\ding{51}}
 &&\mathrm{\textbf{(f) }} 21\,\text{\ding{51}}
\end{aligned} $
\\[2.9mm]
\rule[1.45mm]{\linewidth}{0.5pt}
$\boxed{\bm{\sigma}} \quad \begin{aligned}
\romannumeral1 : \; &\textbf{F} 
 &&\mathrm{\textbf{(a) }} \text{rastúca}\,\text{\ding{51}}
 &&\mathrm{\textbf{(b) }} \left(-\infty , 0\right\rangle  ,   \left\langle\nicefrac{1}{4} , \infty\right)\,\text{\ding{51}}
 &&\mathrm{\textbf{(c) }} 1\,\text{\ding{55}}
 &&\mathrm{\textbf{(d) }} \text{áno}\,\text{\ding{55}}
 &&\mathrm{\textbf{(e) }} -12\,\text{\ding{55}}
 &&\mathrm{\textbf{(f) }} 2\,\text{\ding{51}}
\\[-0.09999999999999998mm]
\romannumeral2 : \; &\textbf{Ľ} 
 &&\mathrm{\textbf{(a) }} \text{rastúca}\,\text{\ding{55}}
 &&\mathrm{\textbf{(b) }} \left\langle-1 , 0\right\rangle\,\text{\ding{55}}
 &&\mathrm{\textbf{(c) }} 1\,\text{\ding{55}}
 &&\mathrm{\textbf{(d) }} \text{áno}\,\text{\ding{51}}
 &&\mathrm{\textbf{(e) }} -24\,\text{\ding{51}}
 &&\mathrm{\textbf{(f) }} 70\,\text{\ding{51}}
\\[-0.09999999999999998mm]
\romannumeral3 : \; &\textbf{A} 
 &&\mathrm{\textbf{(a) }} \text{klesajúca}\,\text{\ding{51}}
 &&\mathrm{\textbf{(b) }} \left\langle-2 , 0\right\rangle\,\text{\ding{55}}
 &&\mathrm{\textbf{(c) }} 1\,\text{\ding{55}}
 &&\mathrm{\textbf{(d) }} \text{áno}\,\text{\ding{55}}
 &&\mathrm{\textbf{(e) }} -12\,\text{\ding{55}}
 &&\mathrm{\textbf{(f) }} 24\,\text{\ding{55}}
\\[-0.09999999999999998mm]
\romannumeral4 : \; &\textbf{K} 
 &&\mathrm{\textbf{(a) }} \text{klesajúca}\,\text{\ding{51}}
 &&\mathrm{\textbf{(b) }} \left\langle-\nicefrac{2}{7} , 0\right\rangle\,\text{\ding{55}}
 &&\mathrm{\textbf{(c) }} 1\,\text{\ding{51}}
 &&\mathrm{\textbf{(d) }} \text{áno}\,\text{\ding{55}}
 &&\mathrm{\textbf{(e) }} -12\,\text{\ding{55}}
 &&\mathrm{\textbf{(f) }} 0\,\text{\ding{55}}
\end{aligned} $
\\[2.9mm]
\rule[1.45mm]{\linewidth}{0.5pt}
$\boxed{\bm{\tau}} \quad \begin{aligned}
\romannumeral1 : \; &\textbf{O} 
 &&\mathrm{\textbf{(a) }} \text{rastúca}\,\text{\ding{51}}
 &&\mathrm{\textbf{(b) }} \left\langle-\nicefrac{1}{3} , 0\right\rangle\,\text{\ding{55}}
 &&\mathrm{\textbf{(c) }} 1\,\text{\ding{51}}
 &&\mathrm{\textbf{(d) }} \text{áno}\,\text{\ding{55}}
 &&\mathrm{\textbf{(e) }} -12\,\text{\ding{51}}
 &&\mathrm{\textbf{(f) }} 3\,\text{\ding{55}}
\\[-0.09999999999999998mm]
\romannumeral2 : \; &\textbf{S} 
 &&\mathrm{\textbf{(a) }} \text{klesajúca}\,\text{\ding{55}}
 &&\mathrm{\textbf{(b) }} \left\langle-\nicefrac{1}{3} , 0\right\rangle\,\text{\ding{51}}
 &&\mathrm{\textbf{(c) }} 1\,\text{\ding{51}}
 &&\mathrm{\textbf{(d) }} \text{áno}\,\text{\ding{55}}
 &&\mathrm{\textbf{(e) }} -16\,\text{\ding{55}}
 &&\mathrm{\textbf{(f) }} 0\,\text{\ding{51}}
\\[-0.09999999999999998mm]
\romannumeral3 : \; &\textbf{O} 
 &&\mathrm{\textbf{(a) }} \text{klesajúca}\,\text{\ding{51}}
 &&\mathrm{\textbf{(b) }} \left\langle-\nicefrac{1}{3} , 0\right\rangle\,\text{\ding{55}}
 &&\mathrm{\textbf{(c) }} 1\,\text{\ding{51}}
 &&\mathrm{\textbf{(d) }} \text{áno}\,\text{\ding{55}}
 &&\mathrm{\textbf{(e) }} -12\,\text{\ding{51}}
 &&\mathrm{\textbf{(f) }} 8\,\text{\ding{55}}
\\[-0.09999999999999998mm]
\romannumeral4 : \; &\textbf{L} 
 &&\mathrm{\textbf{(a) }} \text{rastúca}\,\text{\ding{51}}
 &&\mathrm{\textbf{(b) }} \left(-\infty , 0\right\rangle  ,   \left\langle2 , \infty\right)\,\text{\ding{51}}
 &&\mathrm{\textbf{(c) }} 1\,\text{\ding{51}}
 &&\mathrm{\textbf{(d) }} \text{áno}\,\text{\ding{55}}
 &&\mathrm{\textbf{(e) }} -4\,\text{\ding{55}}
 &&\mathrm{\textbf{(f) }} 15\,\text{\ding{55}}
\end{aligned} $
\\[2.9mm]
\rule[1.45mm]{\linewidth}{0.5pt}
$\boxed{\bm{\upsilon}} \quad \begin{aligned}
\romannumeral1 : \; &\textbf{Ž} 
 &&\mathrm{\textbf{(a) }} \text{klesajúca}\,\text{\ding{55}}
 &&\mathrm{\textbf{(b) }} \left\langle-\nicefrac{1}{2} , 0\right\rangle\,\text{\ding{51}}
 &&\mathrm{\textbf{(c) }} 1\,\text{\ding{51}}
 &&\mathrm{\textbf{(d) }} \text{áno}\,\text{\ding{51}}
 &&\mathrm{\textbf{(e) }} -8\,\text{\ding{55}}
 &&\mathrm{\textbf{(f) }} 24\,\text{\ding{51}}
\\[-0.09999999999999998mm]
\romannumeral2 : \; &\textbf{A} 
 &&\mathrm{\textbf{(a) }} \text{klesajúca}\,\text{\ding{51}}
 &&\mathrm{\textbf{(b) }} \left(-\infty , 0\right\rangle  ,   \left\langle\nicefrac{1}{3} , \infty\right)\,\text{\ding{55}}
 &&\mathrm{\textbf{(c) }} 1\,\text{\ding{55}}
 &&\mathrm{\textbf{(d) }} \text{áno}\,\text{\ding{55}}
 &&\mathrm{\textbf{(e) }} -8\,\text{\ding{55}}
 &&\mathrm{\textbf{(f) }} 4\,\text{\ding{55}}
\\[-0.09999999999999998mm]
\romannumeral3 : \; &\textbf{B} 
 &&\mathrm{\textbf{(a) }} \text{rastúca}\,\text{\ding{51}}
 &&\mathrm{\textbf{(b) }} \left\langle-\nicefrac{2}{5} , 0\right\rangle\,\text{\ding{51}}
 &&\mathrm{\textbf{(c) }} 1\,\text{\ding{55}}
 &&\mathrm{\textbf{(d) }} \text{áno}\,\text{\ding{55}}
 &&\mathrm{\textbf{(e) }} -28\,\text{\ding{55}}
 &&\mathrm{\textbf{(f) }} 70\,\text{\ding{55}}
\\[-0.09999999999999998mm]
\romannumeral4 : \; &\textbf{A} 
 &&\mathrm{\textbf{(a) }} \text{klesajúca}\,\text{\ding{51}}
 &&\mathrm{\textbf{(b) }} \left(-\infty , 0\right\rangle  ,   \left\langle2 , \infty\right)\,\text{\ding{55}}
 &&\mathrm{\textbf{(c) }} 1\,\text{\ding{55}}
 &&\mathrm{\textbf{(d) }} \text{áno}\,\text{\ding{55}}
 &&\mathrm{\textbf{(e) }} -4\,\text{\ding{55}}
 &&\mathrm{\textbf{(f) }} 3\,\text{\ding{55}}
\end{aligned} $
\\[2.9mm]
\rule[1.45mm]{\linewidth}{0.5pt}
$\boxed{\bm{\phi}} \quad \begin{aligned}
\romannumeral1 : \; &\textbf{Ú} 
 &&\mathrm{\textbf{(a) }} \text{rastúca}\,\text{\ding{55}}
 &&\mathrm{\textbf{(b) }} \left(-\infty , 0\right\rangle  ,   \left\langle\nicefrac{2}{5} , \infty\right)\,\text{\ding{55}}
 &&\mathrm{\textbf{(c) }} 1\,\text{\ding{51}}
 &&\mathrm{\textbf{(d) }} \text{áno}\,\text{\ding{51}}
 &&\mathrm{\textbf{(e) }} -16\,\text{\ding{55}}
 &&\mathrm{\textbf{(f) }} 0\,\text{\ding{51}}
\\[-0.09999999999999998mm]
\romannumeral2 : \; &\textbf{R} 
 &&\mathrm{\textbf{(a) }} \text{klesajúca}\,\text{\ding{51}}
 &&\mathrm{\textbf{(b) }} \left\langle-\nicefrac{1}{2} , 0\right\rangle\,\text{\ding{51}}
 &&\mathrm{\textbf{(c) }} 1\,\text{\ding{51}}
 &&\mathrm{\textbf{(d) }} \text{áno}\,\text{\ding{55}}
 &&\mathrm{\textbf{(e) }} -8\,\text{\ding{51}}
 &&\mathrm{\textbf{(f) }} 8\,\text{\ding{55}}
\\[-0.09999999999999998mm]
\romannumeral3 : \; &\textbf{A} 
 &&\mathrm{\textbf{(a) }} \text{klesajúca}\,\text{\ding{51}}
 &&\mathrm{\textbf{(b) }} \left(-\infty , 0\right\rangle  ,   \left\langle2 , \infty\right)\,\text{\ding{55}}
 &&\mathrm{\textbf{(c) }} 1\,\text{\ding{55}}
 &&\mathrm{\textbf{(d) }} \text{áno}\,\text{\ding{55}}
 &&\mathrm{\textbf{(e) }} -4\,\text{\ding{55}}
 &&\mathrm{\textbf{(f) }} 3\,\text{\ding{55}}
\\[-0.09999999999999998mm]
\romannumeral4 : \; &\textbf{Z} 
 &&\mathrm{\textbf{(a) }} \text{klesajúca}\,\text{\ding{51}}
 &&\mathrm{\textbf{(b) }} \left\langle-1 , 0\right\rangle\,\text{\ding{55}}
 &&\mathrm{\textbf{(c) }} 1\,\text{\ding{51}}
 &&\mathrm{\textbf{(d) }} \text{áno}\,\text{\ding{51}}
 &&\mathrm{\textbf{(e) }} -8\,\text{\ding{51}}
 &&\mathrm{\textbf{(f) }} 24\,\text{\ding{55}}
\end{aligned} $
\\[2.9mm]
\rule[1.45mm]{\linewidth}{0.5pt}
$\boxed{\bm{\chi}} \quad \begin{aligned}
\romannumeral1 : \; &\textbf{M} 
 &&\mathrm{\textbf{(a) }} \text{rastúca}\,\text{\ding{51}}
 &&\mathrm{\textbf{(b) }} \left\langle-\nicefrac{2}{3} , 0\right\rangle\,\text{\ding{55}}
 &&\mathrm{\textbf{(c) }} 1\,\text{\ding{51}}
 &&\mathrm{\textbf{(d) }} \text{áno}\,\text{\ding{55}}
 &&\mathrm{\textbf{(e) }} -4\,\text{\ding{55}}
 &&\mathrm{\textbf{(f) }} 4\,\text{\ding{51}}
\\[-0.09999999999999998mm]
\romannumeral2 : \; &\textbf{A} 
 &&\mathrm{\textbf{(a) }} \text{klesajúca}\,\text{\ding{51}}
 &&\mathrm{\textbf{(b) }} \left(-\infty , 0\right\rangle  ,   \left\langle\nicefrac{1}{2} , \infty\right)\,\text{\ding{55}}
 &&\mathrm{\textbf{(c) }} 1\,\text{\ding{55}}
 &&\mathrm{\textbf{(d) }} \text{áno}\,\text{\ding{55}}
 &&\mathrm{\textbf{(e) }} -16\,\text{\ding{55}}
 &&\mathrm{\textbf{(f) }} 0\,\text{\ding{55}}
\\[-0.09999999999999998mm]
\romannumeral3 : \; &\textbf{P} 
 &&\mathrm{\textbf{(a) }} \text{klesajúca}\,\text{\ding{51}}
 &&\mathrm{\textbf{(b) }} \left(-\infty , 0\right\rangle  ,   \left\langle\nicefrac{2}{5} , \infty\right)\,\text{\ding{51}}
 &&\mathrm{\textbf{(c) }} 1\,\text{\ding{51}}
 &&\mathrm{\textbf{(d) }} \text{áno}\,\text{\ding{55}}
 &&\mathrm{\textbf{(e) }} -20\,\text{\ding{55}}
 &&\mathrm{\textbf{(f) }} \nicefrac{3}{2}\,\text{\ding{51}}
\\[-0.09999999999999998mm]
\romannumeral4 : \; &\textbf{A} 
 &&\mathrm{\textbf{(a) }} \text{klesajúca}\,\text{\ding{51}}
 &&\mathrm{\textbf{(b) }} \left\langle-\nicefrac{2}{3} , 0\right\rangle\,\text{\ding{55}}
 &&\mathrm{\textbf{(c) }} 1\,\text{\ding{55}}
 &&\mathrm{\textbf{(d) }} \text{áno}\,\text{\ding{55}}
 &&\mathrm{\textbf{(e) }} -12\,\text{\ding{55}}
 &&\mathrm{\textbf{(f) }} 15\,\text{\ding{55}}
\end{aligned} $
\\[2.9mm]
\rule[1.45mm]{\linewidth}{0.5pt}
$\boxed{\bm{\psi}} \quad \begin{aligned}
\romannumeral1 : \; &\textbf{U} 
 &&\mathrm{\textbf{(a) }} \text{klesajúca}\,\text{\ding{51}}
 &&\mathrm{\textbf{(b) }} \left\langle-\nicefrac{2}{9} , 0\right\rangle\,\text{\ding{55}}
 &&\mathrm{\textbf{(c) }} 1\,\text{\ding{51}}
 &&\mathrm{\textbf{(d) }} \text{áno}\,\text{\ding{51}}
 &&\mathrm{\textbf{(e) }} -12\,\text{\ding{55}}
 &&\mathrm{\textbf{(f) }} 3\,\text{\ding{55}}
\\[-0.09999999999999998mm]
\romannumeral2 : \; &\textbf{H} 
 &&\mathrm{\textbf{(a) }} \text{rastúca}\,\text{\ding{51}}
 &&\mathrm{\textbf{(b) }} \left\langle-\nicefrac{2}{9} , 0\right\rangle\,\text{\ding{51}}
 &&\mathrm{\textbf{(c) }} 1\,\text{\ding{55}}
 &&\mathrm{\textbf{(d) }} \text{áno}\,\text{\ding{55}}
 &&\mathrm{\textbf{(e) }} -4\,\text{\ding{51}}
 &&\mathrm{\textbf{(f) }} -\nicefrac{2}{5}\,\text{\ding{55}}
\\[-0.09999999999999998mm]
\romannumeral3 : \; &\textbf{O} 
 &&\mathrm{\textbf{(a) }} \text{rastúca}\,\text{\ding{51}}
 &&\mathrm{\textbf{(b) }} \left(-\infty , 0\right\rangle  ,   \left\langle\nicefrac{1}{2} , \infty\right)\,\text{\ding{55}}
 &&\mathrm{\textbf{(c) }} 1\,\text{\ding{51}}
 &&\mathrm{\textbf{(d) }} \text{áno}\,\text{\ding{55}}
 &&\mathrm{\textbf{(e) }} -28\,\text{\ding{51}}
 &&\mathrm{\textbf{(f) }} 21\,\text{\ding{55}}
\\[-0.09999999999999998mm]
\romannumeral4 : \; &\textbf{L} 
 &&\mathrm{\textbf{(a) }} \text{klesajúca}\,\text{\ding{51}}
 &&\mathrm{\textbf{(b) }} \left(-\infty , 0\right\rangle  ,   \left\langle\nicefrac{2}{5} , \infty\right)\,\text{\ding{51}}
 &&\mathrm{\textbf{(c) }} 1\,\text{\ding{51}}
 &&\mathrm{\textbf{(d) }} \text{áno}\,\text{\ding{55}}
 &&\mathrm{\textbf{(e) }} -24\,\text{\ding{55}}
 &&\mathrm{\textbf{(f) }} 0\,\text{\ding{55}}
\end{aligned} $
\\[2.9mm]
\rule[1.45mm]{\linewidth}{0.5pt}
$\boxed{\bm{\omega}} \quad \begin{aligned}
\romannumeral1 : \; &\textbf{V} 
 &&\mathrm{\textbf{(a) }} \text{rastúca}\,\text{\ding{51}}
 &&\mathrm{\textbf{(b) }} \left(-\infty , 0\right\rangle  ,   \left\langle\nicefrac{2}{5} , \infty\right)\,\text{\ding{51}}
 &&\mathrm{\textbf{(c) }} 1\,\text{\ding{51}}
 &&\mathrm{\textbf{(d) }} \text{áno}\,\text{\ding{51}}
 &&\mathrm{\textbf{(e) }} -8\,\text{\ding{55}}
 &&\mathrm{\textbf{(f) }} \nicefrac{18}{7}\,\text{\ding{55}}
\\[-0.09999999999999998mm]
\romannumeral2 : \; &\textbf{L} 
 &&\mathrm{\textbf{(a) }} \text{rastúca}\,\text{\ding{51}}
 &&\mathrm{\textbf{(b) }} \left\langle-2 , 0\right\rangle\,\text{\ding{51}}
 &&\mathrm{\textbf{(c) }} 1\,\text{\ding{51}}
 &&\mathrm{\textbf{(d) }} \text{áno}\,\text{\ding{55}}
 &&\mathrm{\textbf{(e) }} -4\,\text{\ding{55}}
 &&\mathrm{\textbf{(f) }} 54\,\text{\ding{55}}
\\[-0.09999999999999998mm]
\romannumeral3 : \; &\textbf{E} 
 &&\mathrm{\textbf{(a) }} \text{klesajúca}\,\text{\ding{51}}
 &&\mathrm{\textbf{(b) }} \left\langle-\nicefrac{1}{2} , 0\right\rangle\,\text{\ding{55}}
 &&\mathrm{\textbf{(c) }} 1\,\text{\ding{55}}
 &&\mathrm{\textbf{(d) }} \text{áno}\,\text{\ding{55}}
 &&\mathrm{\textbf{(e) }} -12\,\text{\ding{51}}
 &&\mathrm{\textbf{(f) }} -2\,\text{\ding{55}}
\\[-0.09999999999999998mm]
\romannumeral4 : \; &\textbf{K} 
 &&\mathrm{\textbf{(a) }} \text{klesajúca}\,\text{\ding{51}}
 &&\mathrm{\textbf{(b) }} \left\langle-2 , 0\right\rangle\,\text{\ding{55}}
 &&\mathrm{\textbf{(c) }} 1\,\text{\ding{51}}
 &&\mathrm{\textbf{(d) }} \text{áno}\,\text{\ding{55}}
 &&\mathrm{\textbf{(e) }} -28\,\text{\ding{55}}
 &&\mathrm{\textbf{(f) }} 15\,\text{\ding{55}}
\end{aligned} $
\\[1mm]
\rule[-1mm]{\linewidth}{0.5pt}
\end{center}
\end{varwidth}
\end{center}
\end{landscape}
\end{document}
